% !TEX program = xelatex
% =====================================================================
% Διπλωματική Εργασία - Μαρία Ντάφλου
% Πρωτεομική Ταξινόμηση Ογκολογικών Ασθενών με Machine Learning
%
% Μεταγλώττιση: XeLaTeX (όχι pdfLaTeX, όχι LuaLaTeX)
% =====================================================================

\documentclass[12pt, a4paper]{article}

% --- Geometry ---

% --- Math ---
\usepackage{amsmath,amssymb,amsfonts,mathtools}

% --- Text Justification & Underline ---
\usepackage{ragged2e}
\usepackage[normalem]{ulem}

% --- XeLaTeX: fonts & Greek support ---
\usepackage{fontspec}

\IfFontExistsTF{Liberation Serif}{%
    \setmainfont{Liberation Serif}
    \setsansfont{Liberation Sans}
    \setmonofont{Liberation Mono}
    \newcommand{\activeMainFont}{Liberation Serif}
    \newcommand{\activeSansFont}{Liberation Sans}
    \newcommand{\activeMonoFont}{Liberation Mono}
}{%
\IfFontExistsTF{Times New Roman}{%
    \setmainfont{Times New Roman}
    \setsansfont{Arial}
    \setmonofont{Courier New}
    \newcommand{\activeMainFont}{Times New Roman}
    \newcommand{\activeSansFont}{Arial}
    \newcommand{\activeMonoFont}{Courier New}
}{%
\IfFontExistsTF{Latin Modern Roman}{%
    \setmainfont{Latin Modern Roman}
    \setsansfont{Latin Modern Sans}
    \setmonofont{Latin Modern Mono}
    \newcommand{\activeMainFont}{Latin Modern Roman}
    \newcommand{\activeSansFont}{Latin Modern Sans}
    \newcommand{\activeMonoFont}{Latin Modern Mono}
}{%
    \PackageError{fontspec}{No suitable Greek-capable font found}{Install one of: Liberation Serif, Times New Roman, Latin Modern Roman}
}}}

% --- Paragraph Spacing ---
\usepackage{parskip}

% --- Polyglossia για διαχείριση Ελληνικών/Αγγλικών ---
\usepackage{polyglossia}
\setdefaultlanguage{greek}
\setotherlanguage{english}

% Ρητός ορισμός Greek font για polyglossia
\newfontfamily\greekfont{\activeMainFont}
\newfontfamily\greekfontsf{\activeSansFont}
\newfontfamily\greekfonttt{\activeMonoFont}

% --- Code listings ---
\usepackage{listings}
\usepackage{xcolor}

\definecolor{codegreen}{rgb}{0.0, 0.5, 0.0}
\definecolor{codegray}{rgb}{0.5, 0.5, 0.5}
\definecolor{codepurple}{rgb}{0.58, 0, 0.82}
\definecolor{codeblue}{rgb}{0.13, 0.13, 1.0}
\definecolor{backcolour}{rgb}{0.97, 0.97, 0.94}

\lstdefinestyle{pythonstyle}{
  backgroundcolor=\color{backcolour},
  commentstyle=\color{codegreen}\itshape,
  keywordstyle=\color{codeblue}\bfseries,
  numberstyle=\tiny\color{codegray},
  stringstyle=\color{codepurple},
  basicstyle=\ttfamily\scriptsize,
  breakatwhitespace=false,
  breaklines=true,
  captionpos=b,
  keepspaces=true,
  numbers=left,
  numbersep=5pt,
  showspaces=false,
  showstringspaces=false,
  showtabs=false,
  tabsize=2,
  frame=single,
  rulecolor=\color{black!30},
  language=Python,
  inputencoding=utf8,
  extendedchars=true,
}
\lstset{style=pythonstyle}

% --- Tables ---
\usepackage{longtable,booktabs,array,multirow}
\usepackage{makecell}

% --- Graphics ---
\usepackage{graphicx}

% --- Hyperlinks ---
\usepackage{hyperref}
\hypersetup{
  hidelinks,
  unicode=true,
}
\usepackage{url}
\urlstyle{same}

% --- Disable auto numbering ---
\setcounter{secnumdepth}{-\maxdimen}

% --- Metadata ---
\title{Ανάπτυξη μεθοδολογίας τεχνητής νοημοσύνης στην ανάλυση πρωτεομικών προφίλ για την αναγνώριση βιοδεικτών που σχετίζονται με την εξέλιξη του καρκίνου παχέος εντέρου.}
\author{Μαρία Ελπίδα Ντάφλου}
\date{Αθήνα, Ιούλιος 2026}


% --- Αντιμετώπιση Overfull Lines στη Βιβλιογραφία & Links ---
\usepackage{microtype} % Βελτιώνει δραματικά το spacing και το σπάσιμο των λέξεων
\usepackage{xurl}      % Επιτρέπει στα URLs και DOIs να σπάνε σε ΟΠΟΙΟΔΗΠΟΤΕ σημείο της γραμμής

% --- Geometry (Μειώθηκαν τα περιθώρια left/right σε 2.0cm) ---
\usepackage[a4paper, left=2.0cm, right=2.0cm, top=3cm, bottom=3cm, bindingoffset=0.5cm]{geometry}

% --- Πακέτα για Πίνακες (Προσθήκη του tabularx) ---
\usepackage{longtable,booktabs,array,multirow}
\usepackage{tabularx} 
\usepackage{makecell}

% =====================================================================
\begin{document}

% =====================================================================
% ΕΞΩΦΥΛΛΟ
% =====================================================================
\begin{titlepage}
\thispagestyle{empty}
\begin{center}

\begin{minipage}{0.22\textwidth}
    \flushleft
    \includegraphics[width=3.3cm]{NTUA-logo.png}
\end{minipage}%
\begin{minipage}{0.7\textwidth}
    \flushleft
    {\large
    Εθνικό Μετσόβιο Πολυτεχνείο\\
    Σχολή Ηλεκτρολόγων Μηχανικών \&\\
    Μηχανικών Υπολογιστών\\
    \vspace{0.3cm}
    Toμέας Συστημάτων Μετάδοσης Πληροφορίας και Τεχνολογίας Υλικών}
\end{minipage}

\vspace*{3cm}

{\LARGE\bfseries\uline{Ανάπτυξη μεθοδολογίας τεχνητής νοημοσύνης στην ανάλυση πρωτεομικών προφίλ
για την αναγνώριση βιοδεικτών που σχετίζονται
με την εξέλιξη του καρκίνου παχέος εντέρου}\par}

\vspace{2.5cm}
{\LARGE ΔΙΠΛΩΜΑΤΙΚΗ ΕΡΓΑΣΙΑ\par}

\vspace{3cm}
{\Large\underline{Μαρία Ελπίδα Ντάφλου}\par}

\vspace{3cm}

\flushleft{\large
    \begin{tabular}{l l}
    \textbf{Επιβλέπων}: & Γεώργιος Ματσόπουλος \\
                        & Καθηγητής, ΕΜΠ\\
    \end{tabular}
}

\vfill
\begin{center}
{\large Αθήνα, Ιούλιος 2026}
\end{center}

\end{center}
\end{titlepage}

% --- ΛΕΥΚΗ ΣΕΛΙΔΑ ΜΕΤΑ ΤΟ ΕΞΩΦΥΛΛΟ (Σελίδα 2) ---
\clearpage
\thispagestyle{empty}
~
\clearpage

% =====================================================================
% ΔΕΥΤΕΡΟ ΕΞΩΦΥΛΛΟ ΜΕ ΤΡΙΜΕΛΗ ΕΠΙΤΡΟΠΗ (Σελίδα 3)
%// =====================================================================
\thispagestyle{empty}
\begin{minipage}{0.22\textwidth}
    \flushleft
    \includegraphics[width=3.3cm]{NTUA-logo.png}
\end{minipage}%
\begin{minipage}{0.7\textwidth}
    \flushleft
    {\large
    Εθνικό Μετσόβιο Πολυτεχνείο\\
    Σχολή Ηλεκτρολόγων Μηχανικών \&\\
    Μηχανικών Υπολογιστών\\
    \vspace{0.3cm}
    Toμέας Συστημάτων Μετάδοσης Πληροφορίας και Τεχνολογίας Υλικών}
\end{minipage}

\vspace*{2cm}

\begin{center}
{\LARGE\bfseries\uline{Ανάπτυξη μεθοδολογίας τεχνητής νοημοσύνης στην ανάλυση πρωτεομικών προφίλ
για την αναγνώριση βιοδεικτών που σχετίζονται
με την εξέλιξη του καρκίνου παχέος εντέρου}\par}

\vspace{2cm}
{\LARGE ΔΙΠΛΩΜΑΤΙΚΗ ΕΡΓΑΣΙΑ\par}

\vspace{1cm}
{\Large\underline{Μαρία Ελπίδα Ντάφλου}\par}
\end{center}

\vspace{1cm}
\flushleft{\large
    \begin{tabular}{l l}
    \textbf{Επιβλέπων}: & Γεώργιος Ματσόπουλος\\
                        & Καθηγητής, ΕΜΠ\\
    \end{tabular}
}

\vspace{1cm}
{\large Εγκρίθηκε από την τριμελή εξεταστική επιτροπή την Ημέρα, Μήνας, Έτος.}\\
\vspace{1cm}

\begin{center}
\small
\begin{tabular}{p{0.28\textwidth} p{0.28\textwidth} p{0.28\textwidth}}
\centering\dotfill & \centering\dotfill & \centering\dotfill\tabularnewline
\centering Γεώργιος Ματσόπουλος & \centering Αθανάσιος Παναγόπουλος & \centering Ουρανία Πετροπούλου\tabularnewline
\centering Καθηγητής, ΕΜΠ & \centering Καθηγητής, ΕΜΠ & \centering ΕΔΙΠ, ΕΜΠ\tabularnewline
\end{tabular}
\end{center}

\vfill
\begin{center}
{\large Αθήνα, Ιούλιος 2026}
\end{center}

\newpage

% =====================================================================
% COPYRIGHT (Σελίδα 4)
% =====================================================================
\thispagestyle{empty}
\noindent\underline{\hspace{5cm}}\\
Μαρία Ελπίδα Ντάφλου\\
Διπλωματούχος Ηλεκτρολόγος Μηχανικός και Μηχανικός Υπολογιστών Ε.Μ.Π.

\vspace*{\fill}
\noindent Copyright \copyright\ Μαρία Ελπίδα Ντάφλου, 2026.\\
Με επιφύλαξη παντός δικαιώματος. All rights reserved.

\vspace{0.5cm}
Απαγορεύεται η αντιγραφή, αποθήκευση και διανοµή της παρούσας εργασίας, εξ ολοκλήρου ή τµήµατος αυτής, για εµπορικό σκοπό. Επιτρέπεται η ανατύπωση, αποθήκευση και διανοµή για σκοπό µη κερδοσκοπικό, εκπαιδευτικής ή ερευνητικής φύσης, υπό την προϋπόθεση να αναφέρεται η πηγή προέλευσης και να διατηρείται το παρόν µήνυµα. Ερωτήµατα που αφορούν τη χρήση της εργασίας για κερδοσκοπικό σκοπό πρέπει να απευθύνονται προς τον συγγραφέα.

\vspace{0.5cm}
Οι απόψεις και τα συμπεράσματα που περιέχονται σε αυτό το έγγραφο εκφράζουν τον συγγραφέα και δεν πρέπει να ερμηνευθεί ότι αντιπροσωπεύουν τις επίσημες θέσεις του Εθνικού Μετσόβιου Πολυτεχνείου.

\newpage

% =====================================================================
% ΑΦΙΕΡΩΣΕΙΣ (Σελίδα 5)
% =====================================================================
\thispagestyle{empty}
\vspace*{8cm}
\begin{center}
\textit{Στους ανθρώπους\\
που με στήριξαν}
\end{center}

% --- ΛΕΥΚΗ ΣΕΛΙΔΑ ΜΕΤΑ ΤΙΣ ΑΦΙΕΡΩΣΕΙΣ (Σελίδα 6) ---
\clearpage
\thispagestyle{empty}
~
\clearpage

% =====================================================================
% ΠΕΡΙΛΗΨΗ (Σελίδα 7 - ΠΛΗΡΗΣ ΣΤΟΙΧΙΣΗ)
% =====================================================================
\pagenumbering{arabic} 
\setcounter{page}{7}   
\thispagestyle{plain}  

\justifying % <--- Επιβάλλει Justify από εδώ και πέρα στο έγγραφο

\section*{Περίληψη}
\addcontentsline{toc}{section}{Περίληψη}

Ο καρκίνος του παχέος εντέρου (Colorectal Cancer - CRC) αποτελεί την τρίτη συχνότερη μορφή καρκίνου παγκοσμίως, με τη φυσική του εξέλιξη, την Ακολουθία Αδενώματος-Καρκινώματος, να διαρκεί 10 έως 15 έτη και να προσφέρει ένα κρίσιμο παράθυρο για πρώιμη διάγνωση. Η παρούσα διπλωματική εργασία διερευνά τη δυνατότητα διαχρονικής πρωτεομικής ανάλυσης πλάσματος αίματος με μεθόδους Μηχανικής Μάθησης για την τριπλή ταξινόμηση ασθενών σε Healthy, Adenoma και Cancer, αξιοποιώντας μετρήσεις σε δύο χρονικά σημεία ($T_0$, $T_1$). Η μελέτη χρησιμοποιεί πρωτεομικά δεδομένα Φασματομετρίας Μάζας από $N = 62$ ασθενείς, μέσω μη επεμβατικής υγρής βιοψίας. Σχεδιάστηκαν και αξιολογήθηκαν πέντε διαδοχικά πειράματα: Vanilla RNN με μηχανισμό προσοχής, LSTM, σύγκριση μεθόδων συνόλου (Random Forest, Random Subspace, Rotation Forest, Custom Cascade Forest) και Logistic Regression, βελτιστοποίηση με GridSearch, Voting, Stacking και SMOTE, και αυστηρή στατιστική αξιολόγηση μέσω Nested CV, Permutation Tests και Bootstrap Confidence Intervals. Η αρχική υπόθεση ήταν ότι θα κυριαρχούσαν οι αρχιτεκτονικές Βαθιάς Μάθησης, ωστόσο τα πειράματα ανέδειξαν την \textbf{Logistic Regression }ως νικητή, έναντι Vanilla RNN και LSTM. Η χρονική πληροφορία ωφέλησε ιδιαίτερα τη γραμμική μοντελοποίηση ($+8.2\%$ έναντι του στατικού baseline), επιβεβαιώνοντας τη διαγνωστική αξία της διαχρονικής προσέγγισης. Για την κλάση CRC, το μοντέλο πέτυχε Sensitivity $68.8\%$, Specificity $84.8\%$ και NPV $88.6\%$, ενδείξεις δυνητικής χρησιμότητας ως εργαλείο πρώτου σταδίου ελέγχου (triage). Πέρα από την ταξινόμηση, εντοπίστηκαν τρεις πρωτεΐνες με σαφή βιολογική σημασία ως υποψήφιοι βιοδείκτες: η \textbf{MMP12} (Macrophage Metalloelastase), η \textbf{AREG} (Amphiregulin, EGFR ligand) και η \textbf{L1CAM} (L1 Cell Adhesion Molecule). Η L1CAM ιδιαίτερα, αν και κατέλαβε μόλις την 7η θέση στη μονομεταβλητή ANOVA ανάλυση, αναδείχθηκε ως το \#1 χαρακτηριστικό στην πολυμεταβλητή ανάλυση Random Forest, υποδεικνύοντας ότι η σημασία της εκδηλώνεται μέσω συνδυασμού με άλλες πρωτεΐνες.

Η εργασία αναδεικνύει συγκεκριμένους υποψήφιους βιοδείκτες για τον CRC με στατιστικά σημαντικά αποτελέσματα και προτείνει ένα ολοκληρωμένο πλαίσιο σύγκρισης οικογενειών αλγορίθμων και αυστηρής στατιστικής επικύρωσης για βιοϊατρικά datasets μικρού μεγέθους.


\vspace{4cm}
\noindent\textbf{Λέξεις Κλειδιά:} Πρωτεομική, Καρκίνος του Παχέος Εντέρου, Διαχρονική Ανάλυση, Μηχανική Μάθηση, Βιοδείκτες

% --- ΛΕΥΚΗ ΣΕΛΙΔΑ ΜΕΤΑ ΤΗΝ ΠΕΡΙΛΗΨΗ (Σελίδα 8) ---
\clearpage
\thispagestyle{empty}
~
\clearpage

% =====================================================================
% ABSTRACT (ENGLISH) (Σελίδα 9)
% =====================================================================
\thispagestyle{plain} 
\begin{english}
\justifying % Εξασφαλίζει justify και για το αγγλικό περιβάλλον
\section*{Abstract}
\addcontentsline{toc}{section}{Abstract}


 Colorectal Cancer (CRC) is the third most common type of cancer worldwide. Its natural progression, known as the Adenoma-Carcinoma Sequence, typically spans 10 to 15 years, offering a critical window for early diagnosis. This diploma thesis investigates the potential of longitudinal proteomic plasma analysis combined with Machine Learning methods for the three-class classification of patients into Healthy, Adenoma, and Cancer categories, leveraging measurements at two timepoints ($T_0$ and $T_1$). The study uses Mass Spectrometry proteomic data from $N = 62$ patients, collected through non-invasive liquid biopsy. Five successive experiments were designed and evaluated: a Vanilla RNN with attention mechanism, an LSTM, a comparison of ensemble methods (Random Forest, Random Subspace, Rotation Forest, Custom Cascade Forest) and Logistic Regression, hyperparameter optimization with GridSearch, Voting, Stacking, and SMOTE, and rigorous statistical evaluation via Nested Cross-Validation, Permutation Tests, and Bootstrap Confidence Intervals. The initial hypothesis was that Deep Learning architectures would prevail, yet the experiments revealed \textbf{Logistic Regression }as the winning model, outperforming both Vanilla RNN and LSTM. The longitudinal information proved especially beneficial to linear modeling (a gain of $+8.2\%$ over the static baseline), confirming the diagnostic value of the temporal approach. For the minority CRC class, the model achieved Sensitivity $68.8\%$, Specificity $84.8\%$, and NPV $88.6\%$, indicating potential utility as a first-line triage tool. Beyond classification, three proteins emerged with clear biological significance as candidate biomarkers: \textbf{MMP12} (Macrophage Metalloelastase), \textbf{AREG} (Amphiregulin, an EGFR ligand), and \textbf{L1CAM} (L1 Cell Adhesion Molecule). Notably, L1CAM ranked only 7th in univariate ANOVA but emerged as the \#1 feature in multivariate Random Forest analysis, indicating that its importance is expressed through combinations with other proteins.

This work contributes on two levels: empirically, it identifies specific candidate biomarkers for CRC with statistically significant results; methodologically, it proposes an integrated framework for comparing algorithm families and rigorous statistical validation for small-N biomedical datasets.

\vspace{4cm}
\noindent\textbf{Keywords:} Proteomics, Colorectal Cancer, Longitudinal Analysis, Machine Learning, Biomarkers
\end{english}

% --- ΛΕΥΚΗ ΣΕΛΙΔΑ ΜΕΤΑ ΤΟ ABSTRACT (Σελίδα 10) ---
\clearpage
\thispagestyle{empty}
~
\clearpage

% =====================================================================
% ΕΥΧΑΡΙΣΤΙΕΣ (Σελίδα 11)
% =====================================================================
\thispagestyle{plain}
\section*{Ευχαριστίες}
\addcontentsline{toc}{section}{Ευχαριστίες}


Καταρχάς, ευχαριστώ ιδιαίτερα ..............

Ευχαριστώ επίσης ............


\vspace{1cm}
\hfill\textit{Μαρία Ελπίδα Ντάφλου}\\
\hspace*{\fill}\textit{Αθήνα, Ιούλιος 2026}

% --- ΛΕΥΚΗ ΣΕΛΙΔΑ ΜΕΤΑ ΤΙΣ ΕΥΧΑΡΙΣΤΙΕΣ (Σελίδα 12) ---
\clearpage
\thispagestyle{empty}
~
\clearpage

% =====================================================================
% ΠΕΡΙΕΧΟΜΕΝΑ (Σελίδα 13)
% =====================================================================
\thispagestyle{plain}
\renewcommand{\contentsname}{Περιεχόμενα}
\tableofcontents
\newpage

% =====================================================================
% ΣΩΜΑ ΕΡΓΑΣΙΑΣ - ΚΕΦΑΛΑΙΑ
% =====================================================================
\section{Κεφάλαιο 1: Εισαγωγή (Introduction)}

\subsection{1.1 Το Κλινικό Πρόβλημα: Η Πρόκληση της Πρώιμης Διάγνωσης στην Ογκολογία}

Ο καρκίνος αποτελεί μία από τις κύριες αιτίες νοσηρότητας και θνησιμότητας παγκοσμίως, με τον Παγκόσμιο Οργανισμό Υγείας (ΠΟΥ) να καταγράφει σχεδόν 10 εκατομμύρια θανάτους ετησίως. Μεταξύ των διαφόρων τύπων κακοήθειας, ο Καρκίνος του Παχέος Εντέρου (Colorectal Cancer - CRC) κατέχει εξέχουσα θέση, αποτελώντας την τρίτη συχνότερη μορφή καρκίνου και τη δεύτερη αιτία θανάτου από καρκίνο παγκοσμίως [70]. Το παράδοξο του CRC είναι ότι, ενώ είναι εξαιρετικά θανατηφόρος στα προχωρημένα στάδια (μετάσταση), είναι ταυτόχρονα και μία από τις πιο ιάσιμες μορφές αν ανιχνευθεί εγκαίρως. Η φυσική ιστορία της νόσου ακολουθεί μια μακροχρόνια, σταδιακή εξέλιξη γνωστή ως \textbf{"Ακολουθία Αδενώματος-Καρκινώματος" (Adenoma-Carcinoma Sequence)}. Η μετάβαση από το φυσιολογικό επιθήλιο σε καλοήθη πολύποδα (αδένωμα) και τελικά σε διηθητικό καρκίνωμα μπορεί να διαρκέσει 10 έως 15 έτη [71]. Αυτό το χρονικό "παράθυρο ευκαιρίας" καθιστά τον CRC ιδανικό στόχο για προγράμματα προσυμπτωματικού ελέγχου (screening).

\subsubsection{1.1.1 Περιορισμοί των Συμβατικών Διαγνωστικών Μεθόδων}

Σήμερα, η διάγνωση βασίζεται κυρίως σε δύο πυλώνες:

\begin{enumerate}
\item \textbf{Επεμβατικές Μέθοδοι (Κολονοσκόπηση):} Αποτελεί το "Gold Standard" για τη διάγνωση. Ωστόσο, είναι μια επεμβατική διαδικασία, με υψηλό κόστος, δυσφορία για τον ασθενή και μικρό (αλλά υπαρκτό) κίνδυνο επιπλοκών (π.χ. διάτρηση εντέρου). Λόγω αυτών των παραγόντων, η συμμόρφωση του γενικού πληθυσμού στις οδηγίες screening παραμένει χαμηλή.
\item \textbf{Μη Επεμβατικές Μέθοδοι (FIT, CEA):} Τεστ όπως η ανίχνευση αιμοσφαιρίνης στα κοπράνα (FIT) έχουν χαμηλή ευαισθησία (sensitivity) για τα αρχικά στάδια (αδενώματα), ενώ οι κλασικοί καρκινικοί δείκτες στο αίμα (π.χ. CEA, CA19-9) αυξάνονται συνήθως όταν ο όγκος έχει ήδη μεγαλώσει σημαντικά, καθιστώντας τους ακατάλληλους για πρώιμη διάγνωση [74].
\end{enumerate}

Συνεπώς, υπάρχει μια επιτακτική κλινική ανάγκη για την ανάπτυξη νέων, ελάχιστα επεμβατικών και υψηλής ακρίβειας εργαλείων που μπορούν να ανιχνεύσουν τη νόσο στο στάδιο του αδενώματος.

\subsection{1.2 Η Επανάσταση της Υγρή Βιοψίας και της Πρωτεομικής}

Η λύση στο πρόβλημα της χαμηλής συμμόρφωσης έρχεται μέσω της \textbf{Υγρής Βιοψίας (Liquid Biopsy)}. Η έννοια βασίζεται στην παρατήρηση ότι οι όγκοι, ακόμα και σε μικροσκοπικό επίπεδο, αλληλεπιδρούν με το μικροπεριβάλλον τους και απελευθερώνουν βιομόρια (κύτταρα, DNA, εξωσώματα, πρωτεΐνες) στην κυκλοφορία του αίματος.

\subsubsection{1.2.1 Γιατί Πρωτεΐνες;}

Ενώ η γονιδιωματική (Genomics) αποκαλύπτει τι \emph{μπορεί} να συμβεί (βάσει των μεταλλάξεων), η Πρωτεομική (Proteomics) φανερώνει τι \emph{πραγματικά συμβαίνει} σε πραγματικό χρόνο. Οι πρωτεΐνες είναι οι τελικοί εκτελεστές των βιολογικών διεργασιών. Αλλαγές στη συγκέντρωση πρωτεϊνών που σχετίζονται με τη φλεγμονή, την αγγειογένεση (π.χ. VEGF) ή την αναδόμηση του ιστού (π.χ. MMPs) μπορούν να προηγηθούν της κλινικής εμφάνισης του όγκου κατά μήνες ή και χρόνια [75]. Η σύγχρονη τεχνολογία Φασματομετρίας Μάζας (Mass Spectrometry - MS) επιτρέπει τον ταυτόχρονο προσδιορισμό χιλιάδων πρωτεϊνών σε μια σταγόνα πλάσματος, δημιουργώντας ένα πολυδιάστατο "μοριακό αποτύπωμα" του ασθενούς.

\subsection{1.3 Το Πρόβλημα της Στατικής Ανάλυσης και η "Δυναμική" Λύση}

\subsubsection{1.3.1 Η Ανεπάρκεια του "Στιγμιότυπου" (Snapshot)}

Η πλειονότητα των διαγνωστικών μοντέλων Μηχανικής Μάθησης εκπαιδεύεται με στατικά δεδομένα: παίρνουμε μια μέτρηση $T_0$ και προσπαθούμε να ταξινομήσουμε τον ασθενή. Αυτή η προσέγγιση αγνοεί τη θεμελιώδη βιολογική αλήθεια πως η υγεία και η ασθένεια είναι δυναμικές διαδικασίες, όχι στατικές καταστάσεις.

Μια μεμονωμένη τιμή πρωτεΐνης μπορεί να είναι παραπλανητική λόγω της \textbf{Δια-ατομικής Μεταβλητότητας (Inter-patient variability)} αφού το "φυσιολογικό" επίπεδο μιας πρωτεΐνης διαφέρει από άνθρωπο σε άνθρωπο. Μια τιμή $X$ μπορεί να είναι φυσιολογική για τον ασθενή Α, αλλά παθολογική (υψηλή) για τον ασθενή Β, και από \textbf{Θορύβο (Noise)} με παράγοντες όπως ιώσεις, στρες ή διατροφή να επηρεάζουν προσωρινά τα επίπεδα πρωτεϊνών.

\subsubsection{1.3.2 Η Υπόθεση της Χρονικής Εξέλιξης (Research Hypothesis)}

Η παρούσα εργασία ασχολείται με την έννοια της \textbf{Διαχρονικής Πρωτεομικής (Longitudinal Proteomics)}. Η κεντρική υπόθεση εργασίας είναι:

\begin{quote}
\emph{"Η διαγνωστική πληροφορία μεγιστοποιείται όταν παρατηρούμε την τροχιά (trajectory) των βιοδεικτών στον χρόνο, αντί για την απόλυτη τιμή τους σε μια δεδομένη στιγμή."}
\end{quote}

Συγκεκριμένα, εξετάζουμε το σενάριο όπου λαμβάνονται δύο δείγματα ($T_0$ και $T_1$). Η υπόθεση είναι ότι η μετάβαση $T_0 - T_1$ αποκαλύπτει την "κλίση" της νόσου. Ένα σταθερό επίπεδο (ακόμα και ελαφρώς αυξημένο) μπορεί να υποδηλώνει καλοήθεια, ενώ μια αυξητική τάση μπορεί να σημάνει την έναρξη καρκινογένεσης. Για την αποκωδικοποίηση αυτών των χρονικών μοτίβων, οι παραδοσιακές στατιστικές μέθοδοι δεν επαρκούν και απαιτούνται αλγόριθμοι Βαθιάς Μάθησης (Deep Learning) ικανοί να μοντελοποιούν ακολουθίες.

\subsection{1.4 Σκοπός της Εργασίας}

Σκοπός της εργασίας είναι η ανάπτυξη ενός ολοκληρωμένου υπολογιστικού πλαισίου (computational framework) για την πρόβλεψη της κατάστασης του ασθενούς (Healthy / Adenoma / Cancer) βάσει χρονικών πρωτεομικών δεδομένων, καθώς και η ενασχόληση με το ερώτημα ποια οικογένεια αλγορίθμων είναι καταλληλότερη για τέτοιου είδους δεδομένα μικρού μεγέθους ($N < 100$). Η εργασία εκκινεί από την υπόθεση ότι οι αρχιτεκτονικές Βαθιάς Μάθησης (RNN/LSTM) είναι κατάλληλες για χρονοσειρές, και την εξετάζει συστηματικά μέσω συγκριτικών πειραμάτων.

Οι ειδικότεροι στόχοι περιλαμβάνουν:

\begin{enumerate}
\item \textbf{Συστηματική Σύγκριση Οικογενειών Αλγορίθμων:} Εμπειρική αξιολόγηση της απόδοσης Βαθιάς Μάθησης (Vanilla RNN, LSTM με Attention) έναντι Μεθόδων Συνόλου (Random Forest, Random Subspace, Rotation Forest, Custom Cascade Forest) και γραμμικών baseline μοντέλων (Logistic Regression) σε πραγματικό πρωτεομικό dataset.
\item \textbf{Ανάπτυξη Υβριδικών Αρχιτεκτονικών DL:} Σχεδιασμός και υλοποίηση μοντέλων που συνδυάζουν \textbf{1D-CNNs} (για εξαγωγή χαρακτηριστικών από το πρωτεϊνικό φάσμα), \textbf{Recurrent Neural Networks (RNNs/LSTMs)} (για τη μάθηση της χρονικής εξέλιξης), και \textbf{Attention Mechanisms} (για ερμηνευσιμότητα και ανακάλυψη βιοδεικτών).
\item \textbf{Σύγκριση Πειραματικών Παραδειγμάτων (Ablation Study):} Αυστηρή σύγκριση της προτεινόμενης χρονικής μεθόδου (Temporal) έναντι της καθιερωμένης στατικής (Static) και της απλοϊκής διαφορικής μεθόδου (Delta), για να ποσοτικοποιηθεί η αξία της χρονικής πληροφορίας.
\item \textbf{Αντιμετώπιση του "Small Data Problem":} Εφαρμογή εξειδικευμένων τεχνικών (Imputation, Regularization, Stratified Cross-Validation) για την εκπαίδευση μοντέλων σε datasets με περιορισμένο αριθμό ασθενών, ένα συχνό πρόβλημα στην ιατρική έρευνα.
\item \textbf{Αυστηρή Στατιστική Αξιολόγηση:} Πέρα από την απλή cross-validation, εφαρμογή Nested CV (αμερόληπτη εκτίμηση), Repeated CV, Bootstrap Confidence Intervals, και Permutation Tests, ώστε τα συμπεράσματα να διαθέτουν επιστημονική εγκυρότητα ανθεκτική στην κριτική για το μικρό μέγεθος δείγματος.
\item \textbf{Βιολογική Ερμηνευσιμότητα:} Ταυτοποίηση των πρωτεϊνών που οδηγούν τις αποφάσεις του μοντέλου μέσω univariate (ANOVA $F$-score), multivariate (Random Forest MDI) και attention-based αναλύσεων, προσφέροντας στους βιολόγους νέους υποψήφιους βιοδείκτες προς περαιτέρω διερεύνηση.
\end{enumerate}

Η εργασία ξεκίνησε με την προσδοκία ότι το Deep Learning θα κυριαρχούσε στα πειράματα, λόγω της εγγενούς ικανότητάς του να μοντελοποιεί χρονικές εξαρτήσεις. Η πειραματική πραγματικότητα, ωστόσο, οδήγησε στο συμπέρασμα της ότι σε small-N βιοϊατρικά datasets, τα απλούστερα μοντέλα (Logistic Regression) μπορούν να υπερτερούν.

\subsection{1.5 Δομή της Εργασίας}

Η εργασία διαρθρώνεται σε έξι κεφάλαια:

\begin{itemize}
\item \textbf{Κεφάλαιο 2:} Παρουσιάζει το θεωρητικό υπόβαθρο, καλύπτοντας τόσο τα Νευρωνικά Δίκτυα και τα Recurrent αρχιτεκτονικά (RNN, LSTM, Attention) όσο και τις Μεθόδους Συνόλου (Bagging, Random Forest, Rotation Forest, Deep Forest). Ανασκοπεί τη βιβλιογραφία σχετικά με την καταλληλότητα κάθε οικογένειας για datasets μικρού μεγέθους.
\item \textbf{Κεφάλαιο 3:} Περιγράφει αναλυτικά τη μεθοδολογία: προ-επεξεργασία πρωτεομικών δεδομένων, στρατηγικές imputation, feature selection, αρχιτεκτονικές των DL μοντέλων (Vanilla RNN, LSTM), υλοποίηση των ensemble μεθόδων, και τα τρία πειραματικά παραδείγματα (Static, Delta, Temporal). Καλύπτει επίσης τη βελτιστοποίηση υπερπαραμέτρων και την αυστηρή στατιστική αξιολόγηση (Nested CV, Permutation Test, Bootstrap CI).
\item \textbf{Κεφάλαιο 4:} Παραθέτει τα πειραματικά αποτελέσματα, συγκρίνοντας τις επιδόσεις των DL και ensemble μοντέλων, παρουσιάζοντας πίνακες σύγχυσης, στατιστικούς ελέγχους σημαντικότητας, learning curves, calibration analysis, και κλινικά σημαντικές μετρικές ανά κλάση.
\item \textbf{Κεφάλαιο 5:} Συζητά τα ευρήματα υπό το πρίσμα της βιολογίας και της μεθοδολογίας, αναλύοντας τον ρόλο των κορυφαίων πρωτεϊνών (MMP12, AREG, L1CAM) στην παθογένεια του καρκίνου, και ερμηνεύοντας γιατί τα γραμμικά μοντέλα υπερτέρησαν σε αυτή την περίπτωση.
\item \textbf{Κεφάλαιο 6:} Συνοψίζει τα συμπεράσματα, τα μεθοδολογικά διδάγματα, τους περιορισμούς, και προτείνει κατευθύνσεις για μελλοντική έρευνα.
\end{itemize}


\sloppy
\section{Κεφάλαιο 2: Θεωρητικό Υπόβαθρο (Theoretical Background)}

Στο παρόν κεφάλαιο παρουσιάζεται το θεωρητικό πλαίσιο πάνω στο οποίο εδράζονται όλα τα πειράματα της παρούσας διπλωματικής.

\subsection{2.1 Μηχανική Μάθηση σε Βιοϊατρικά Δεδομένα Μικρού Μεγέθους}

Η Μηχανική Μάθηση (Machine Learning - ML), υποπεδίο της Τεχνητής Νοημοσύνης, εστιάζει στην ανάπτυξη αλγορίθμων που μαθαίνουν μοτίβα από δεδομένα και κάνουν προβλέψεις σε άγνωστα δείγματα. Στον χώρο της Υπολογιστικής Βιολογίας και της Ιατρικής Ακριβείας (Precision Medicine), η Επιβλεπόμενη Μάθηση (Supervised Learning) κυριαρχεί. Σε αυτό το παράδειγμα, ο αλγόριθμος εκπαιδεύεται σε ένα σύνολο ζευγών εισόδου-εξόδου $\{(x_1, y_1), \ldots, (x_N, y_N)\}$, με στόχο να προσεγγίσει τη συνάρτηση απεικόνισης $f: X \to Y$ που ελαχιστοποιεί το σφάλμα πρόβλεψης. 

Ωστόσο, η εφαρμογή ML σε βιοϊατρικά δεδομένα, όπως τα πρωτεομικά της παρούσας μελέτης, συνοδεύεται από δομικές προκλήσεις που είναι εγγενώς διαφορετικές από αυτές μιας τυπικής εφαρμογής σε μεγάλα datasets υπολογιστικής όρασης ή επεξεργασίας φυσικής γλώσσας.

Η πρώτη και κρισιμότερη πρόκληση είναι η \textbf{Κατάρα της Διαστασιμότητας} (Curse of Dimensionality), όπως τη διατύπωσε ο Richard Bellman το 1957. Καθώς αυξάνεται ο αριθμός των διαστάσεων (χαρακτηριστικών/πρωτεϊνών), ο όγκος του χώρου αυξάνεται εκθετικά, καθιστώντας τα διαθέσιμα δεδομένα αραιά (sparse). Για να επιτευχθεί στατιστική σημαντικότητα, ο αριθμός των δειγμάτων $N$ θα έπρεπε να αυξάνεται εκθετικά με τις διαστάσεις $P$. Στην περίπτωσή μας, όπου $P \approx 500$ (πρωτεΐνες) και $N < 100$ (ασθενείς), υπάρχει μεγάλος κίνδυνος το μοντέλο να βρει τυχαίες συσχετίσεις (spurious correlations), οδηγώντας σε υπερ-προσαρμογή (overfitting) [25], [30]. 

Η αντιμετώπιση απαιτεί συνδυασμό τεχνικών μείωσης διαστάσεων, regularization και επιλογής χαρακτηριστικών.

Η δεύτερη πρόκληση είναι η \textbf{Ανισορροπία Κλάσεων} (Class Imbalance). Στα ιατρικά σύνολα δεδομένων, οι ασθενείς με τη νόσο (θετικά δείγματα) είναι συχνά μειοψηφία. Ένας ταξινομητής που εκπαιδεύεται ελαχιστοποιώντας απλώς το συνολικό σφάλμα (accuracy), τείνει να μεροληπτεί υπέρ της πλειοψηφούσας κλάσης (Healthy). Αυτό αντιμετωπίζεται μέσω τεχνικών όπως η Σταθμισμένη Συνάρτηση Κόστους (Weighted Loss Function) και η Στρωματοποιημένη Δειγματοληψία (Stratified Sampling) [44], [45].

Σε αυτό το γενικό πλαίσιο, δύο μεγάλες οικογένειες αλγορίθμων κυριαρχούν στη βιβλιογραφία: τα Νευρωνικά Δίκτυα (Neural Networks), τα οποία έχουν αποδειχθεί ιδιαίτερα ισχυρά σε προβλήματα με μεγάλα δεδομένα και πολύπλοκη δομή, και οι Μέθοδοι Συνόλου (Ensemble Methods) βασισμένες σε δέντρα αποφάσεων, οι οποίες είναι θεωρητικά καταλληλότερες για μικρότερα, δομημένα (tabular) datasets όπως αυτό της παρούσας μελέτης. Στις ενότητες που ακολουθούν αναπτύσσονται και οι δύο, στη σειρά που εφαρμόζονται και στα πειράματα.

\subsection{2.2 Νευρωνικά Δίκτυα για Ακολουθιακά Δεδομένα}

Τα Τεχνητά Νευρωνικά Δίκτυα (Artificial Neural Networks - ANNs) είναι μαθηματικά μοντέλα εμπνευσμένα από τη βιολογική αρχιτεκτονική του εγκεφάλου. Η θεμελιώδης μονάδα είναι ο τεχνητός νευρώνας (perceptron). Ένας νευρώνας $k$ δέχεται ένα διάνυσμα εισόδου $\mathbf{x} = [x_1, x_2, \ldots, x_m]^T$, το οποίο πολλαπλασιάζεται με ένα διάνυσμα βαρών $\mathbf{w}_k = [w_{k1}, w_{k2}, \ldots, w_{km}]^T$. Σε αυτό προστίθεται μια σταθερά πόλωσης (bias) $b_k$, και το αποτέλεσμα περνά από μια μη γραμμική συνάρτηση ενεργοποίησης $\phi(\cdot)$:
$$y_k = \phi \left( \sum_{j=1}^{m} w_{kj} x_j + b_k \right) = \phi(\mathbf{w}_k^T \mathbf{x} + b_k)$$
Οι μη γραμμικές συναρτήσεις ενεργοποίησης είναι κρίσιμες, διότι χωρίς αυτές, ένα δίκτυο πολλαπλών επιπέδων θα κατέρρεε σε ένα απλό γραμμικό μοντέλο. Στις πιο διαδεδομένες ανήκουν η \textbf{ReLU} $f(x) = \max(0, x)$ (η πιο δημοφιλής λόγω ταχείας σύγκλισης και αποφυγής vanishing gradient στα θετικά πεδία), η \textbf{Sigmoid} $\sigma(x) = 1 / (1 + e^{-x})$ (χρησιμοποιείται στις πύλες των LSTM επειδή παράγει τιμές στο $(0,1)$, ερμηνεύσιμες ως πιθανότητες), και η \textbf{tanh} $\tanh(x) = (e^x - e^{-x})/(e^x + e^{-x})$ (παράγει τιμές στο $(-1,1)$, κεντράροντας τα δεδομένα στο μηδέν, και χρησιμοποιείται ευρέως στα RNNs).


Σε πολλά ιατρικά προβλήματα, η πληροφορία δεν βρίσκεται μόνο στα μεμονωμένα χαρακτηριστικά αλλά και στη χρονική τους αλληλουχία. Τα \textbf{Αναδρομικά Νευρωνικά Δίκτυα (Recurrent Neural Networks - RNNs)}, που εισήχθησαν στη δεκαετία του 1980 (Rumelhart et al., 1986), σχεδιάστηκαν για να χειρίζονται ακολουθιακά δεδομένα. Η βασική τους καινοτομία είναι η ύπαρξη βρόχων ανάδρασης, που επιτρέπουν στην πληροφορία να παραμένει στο δίκτυο. Σε κάθε χρονική στιγμή $t$, το RNN δέχεται την τρέχουσα είσοδο $x_t$ και την κρυφή κατάσταση (hidden state) $h_{t-1}$ από το προηγούμενο βήμα, η οποία λειτουργεί ως η ``μνήμη'' του δικτύου. Η ενημέρωση της κρυφής κατάστασης περιγράφεται από την εξίσωση:
$$h_t = \tanh(W_{hh} h_{t-1} + W_{xh} x_t + \mathbf{b}_h)$$
όπου $W_{hh} \in \mathbb{R}^{d_h \times d_h}$ είναι ο πίνακας βαρών που ελέγχει πώς η προηγούμενη μνήμη επηρεάζει την τρέχουσα (Recurrent weights), $W_{xh} \in \mathbb{R}^{d_h \times d_x}$ ο πίνακας που ελέγχει την επίδραση της νέας εισόδου (Input weights), $\mathbf{b}_h$ το διάνυσμα πόλωσης, και $\tanh$ η συνάρτηση ενεργοποίησης. Τα RNNs έχουν χρησιμοποιηθεί ευρέως στην Επεξεργασία Φυσικής Γλώσσας (NLP) και τη γονιδιωματική. Πάσχουν όμως από το πρόβλημα της \textbf{Εξαφάνισης της Κλίσης} (Vanishing Gradient): κατά την εκπαίδευση με Backpropagation Through Time (BPTT), οι κλίσεις (gradients) πολλαπλασιάζονται διαδοχικά με τον πίνακα $W_{hh}$. Αν οι ιδιοτιμές του $W_{hh}$ είναι μικρότερες του 1, το gradient τείνει εκθετικά στο μηδέν, καθιστώντας αδύνατη τη μάθηση μακροχρόνιων εξαρτήσεων (long-term dependencies).

Για να αντιμετωπιστεί αυτό το πρόβλημα μνήμης των RNNs, οι Hochreiter και Schmidhuber (1997) παρουσίασαν τα \textbf{Long Short-Term Memory} δίκτυα (LSTMs). Η κεντρική ιδέα είναι η εισαγωγή μιας ``λεωφόρου πληροφορίας'', της Κατάστασης Κελιού (Cell State, $C_t$), η οποία επιτρέπει στην πληροφορία να ρέει σχεδόν ανεμπόδιστα στον χρόνο. Η ροή ελέγχεται από δομές που ονομάζονται \textbf{πύλες} (gates). Συγκεκριμένα, το LSTM χρησιμοποιεί τρεις πύλες σε κάθε χρονικό βήμα. Η \emph{Forget Gate} αποφασίζει ποια πληροφορία από το παρελθόν ($C_{t-1}$) είναι άχρηστη και πρέπει να διαγραφεί, μέσω της $f_t = \sigma(W_f \cdot [h_{t-1}, x_t] + b_f)$. Αν $f_t \approx 0$, η μνήμη διαγράφεται· αν $f_t \approx 1$, διατηρείται. Η \emph{Input Gate} αποφασίζει ποια νέα πληροφορία είναι σημαντική για να αποθηκευτεί: $i_t = \sigma(W_i \cdot [h_{t-1}, x_t] + b_i)$, ενώ μια υποψήφια μνήμη δημιουργείται ως $\tilde{C}_t = \tanh(W_C \cdot [h_{t-1}, x_t] + b_C)$. Η ενημέρωση της κατάστασης κελιού συνδυάζει την παλιά μνήμη (αφού φιλτραριστεί από το forget gate) με τη νέα (αφού φιλτραριστεί από το input gate):
$$C_t = f_t \odot C_{t-1} + i_t \odot \tilde{C}_t$$
όπου $\odot$ είναι το γινόμενο κατά στοιχείο (Hadamard product). Η πρόσθεση είναι κρίσιμη: επιτρέπει στο gradient να ρέει χωρίς να πολλαπλασιάζεται συνεχώς, λύνοντας το πρόβλημα του Vanishing Gradient. Τέλος, η \emph{Output Gate} καθορίζει τι θα δει το επόμενο επίπεδο: $o_t = \sigma(W_o \cdot [h_{t-1}, x_t] + b_o)$, και η κρυφή κατάσταση εξάγεται ως $h_t = o_t \odot \tanh(C_t)$. Στην παρούσα μελέτη, η χρήση LSTM συγκρίνεται με το Vanilla RNN ακριβώς για να διαπιστωθεί αν η αυξημένη πολυπλοκότητα (περίπου τετραπλάσιες παράμετροι ανά κελί) δικαιολογείται σε μικρές χρονοσειρές ($T=2$).

Πέρα από την επεξεργασία της χρονικής διαδοχής, ένα νευρωνικό μοντέλο μπορεί να ενισχυθεί με \textbf{Μηχανισμούς Προσοχής} (Attention Mechanisms), εμπνευσμένους από την ανθρώπινη οπτική αντίληψη. Ο μηχανισμός αυτός επιτρέπει στο δίκτυο να εστιάζει επιλεκτικά σε συγκεκριμένα τμήματα της εισόδου. Αρχικά αναπτύχθηκε για τη Μηχανική Μετάφραση (Bahdanau et al., 2014) για να λύσει το πρόβλημα της συμπίεσης μεγάλων προτάσεων σε ένα σταθερό διάνυσμα. Σήμερα, με την έλευση των Transformers (Vaswani et al., 2017), αποτελεί τη βάση των σύγχρονων μοντέλων AI (π.χ. GPT, BERT). Στην ιατρική, ένα μοντέλο ``Black Box'' που δίνει απλώς μια πρόβλεψη (``Καρκίνος: 90\%'') δεν είναι αρκετό, οι κλινικοί απαιτούν να γνωρίζουν το ``Γιατί;''. Ο Attention γεφυρώνει αυτό το χάσμα, αποδίδοντας ένα ``βάρος'' σε κάθε χρονικό βήμα ή χαρακτηριστικό και υποδεικνύοντας πού ``κοίταξε'' το μοντέλο [17], [19]. Αν το βάρος στο $T_1$ είναι υψηλό, η εξέλιξη είναι κρίσιμη· αν το βάρος σε μια πρωτεΐνη είναι υψηλό, αυτή αποτελεί πιθανό βιοδείκτη. Στην εργασία αυτή χρησιμοποιείται ο \emph{προσθετικός μηχανισμός προσοχής} (Additive Attention) του Bahdanau. Έστω $h_1, \ldots, h_T$ οι κρυφές καταστάσεις του RNN. Υπολογίζεται πρώτα ένα alignment score μέσω ενός μικρού νευρωνικού δικτύου, $e_t = \mathbf{v}_a^T \tanh(W_a h_t + b_a)$, όπου $W_a, b_a, \mathbf{v}_a$ είναι εκπαιδεύσιμες παράμετροι. Στη συνέχεια εφαρμόζεται Softmax για να μετατραπούν τα scores σε κατανομή πιθανοτήτων, $\alpha_t = \exp(e_t) / \sum_k \exp(e_k)$, και τέλος υπολογίζεται το context vector ως σταθμισμένο άθροισμα, $c = \sum_t \alpha_t h_t$, το οποίο τροφοδοτείται στον τελικό ταξινομητή.

Τέλος, στο ίδιο νευρωνικό πλαίσιο εντάσσονται και τα \textbf{Συνελικτικά Δίκτυα μιας Διάστασης} (1D-CNNs). Ενώ τα κλασικά Convolutional Neural Networks είναι διάσημα για την επεξεργασία εικόνας (2D), η μονοδιάστατη εκδοχή τους είναι εξαιρετικά ισχυρή στην ανάλυση ακολουθιακών δεδομένων όπως σήματα ΗΚΓ, γονιδιωματικές αλληλουχίες και φασματοσκοπία [5]. Η λειτουργία της συνέλιξης περιλαμβάνει την ολίσθηση ενός φίλτρου (kernel) μήκους $K$ πάνω στο διάνυσμα εισόδου $x$. 

Για κάθε θέση $i$, υπολογίζεται το εσωτερικό γινόμενο
$$y[i] = (x * w)[i] = \sum_{k=0}^{K-1} x[i+k] \cdot w[k] + b$$
Στις προτεινόμενες αρχιτεκτονικές της παρούσας μελέτης, ένα 1D-CNN τοποθετείται \emph{πριν} το RNN, ως επίπεδο εξαγωγής χαρακτηριστικών (feature extraction): μαθαίνει να αναγνωρίζει τοπικά μοτίβα (συσχετίσεις μεταξύ γειτονικών τιμών στο διάνυσμα πρωτεϊνών, π.χ. ομάδες πρωτεϊνών που εκφράζονται μαζί), λειτουργεί ως φίλτρο εξομάλυνσης μειώνοντας τον υψίσυχνο θόρυβο της φασματομετρίας μάζας, και προσφέρει translation invariance (αναγνωρίζει ένα μοτίβο ανεξάρτητα από την ακριβή θέση του στην είσοδο).

\subsection{2.3 Μέθοδοι Συνόλου και Δέντρα Αποφάσεων}

Όπως αναφέρθηκε στην αρχή του κεφαλαίου, σε datasets με μικρό αριθμό δειγμάτων ($N < 100$) τα βαθιά νευρωνικά δίκτυα αντιμετωπίζουν σοβαρούς περιορισμούς. Η βιβλιογραφία υποδεικνύει ότι σε τέτοιες περιπτώσεις οι \textbf{Μέθοδοι Συνόλου} (Ensemble Methods) που βασίζονται σε δέντρα αποφάσεων αποτελούν μια αξιόπιστη εναλλακτική, καθώς παρέχουν ανθεκτικότητα στην υπερ-προσαρμογή, ερμηνευσιμότητα και υψηλή απόδοση χωρίς να απαιτούν μεγάλο πλήθος δειγμάτων [79], [80]. 

Τα \textbf{δέντρα αποφάσεων} (Decision Trees) είναι μη-παραμετρικοί αλγόριθμοι επιβλεπόμενης μάθησης που διαμερίζουν αναδρομικά τον χώρο των χαρακτηριστικών σε υποπεριοχές ομοιογενών κλάσεων. Σε κάθε εσωτερικό κόμβο, ο αλγόριθμος επιλέγει το χαρακτηριστικό $X_j$ και το κατώφλι $\theta$ που μεγιστοποιούν τον διαχωρισμό των κλάσεων. 

Η διαδικασία επαναλαμβάνεται αναδρομικά μέχρι να ικανοποιηθεί κάποιο κριτήριο τερματισμού (π.χ. ελάχιστος αριθμός δειγμάτων ανά φύλλο ή μέγιστο βάθος). Για την επιλογή του βέλτιστου διαχωρισμού χρησιμοποιούνται δύο κύρια κριτήρια: η \emph{Gini Impurity} $\text{Gini}(t) = 1 - \sum_{k=1}^{C} p_k^2$ (όπου $p_k$ είναι η αναλογία δειγμάτων της κλάσης $k$ στον κόμβο $t$), η οποία μετράει την πιθανότητα ένα τυχαία επιλεγμένο δείγμα να ταξινομηθεί λανθασμένα αν ταξινομηθεί σύμφωνα με την κατανομή των κλάσεων στον κόμβο, και η \emph{Information Gain (Entropy)} που βασίζεται στη θεωρία πληροφοριών του Shannon, με την εντροπία να ορίζεται ως $H(t) = -\sum_{k=1}^{C} p_k \log_2(p_k)$ και το κέρδος πληροφορίας από τη διαίρεση βάσει του χαρακτηριστικού $A$ ως $IG(t, A) = H(t) - \sum_{v} (|t_v|/|t|) H(t_v)$. 

Παρά την ερμηνευσιμότητά τους, τα μεμονωμένα δέντρα πάσχουν από \emph{υψηλή διακύμανση} (high variance): μικρές αλλαγές στα δεδομένα εκπαίδευσης μπορούν να οδηγήσουν σε εντελώς διαφορετικά δέντρα. Αυτός ο περιορισμός είναι ο λόγος ύπαρξης των μεθόδων συνόλου.

Ο Breiman (1996) εισήγαγε την τεχνική \textbf{Bagging} (Bootstrap Aggregating) ως μέθοδο μείωσης της διακύμανσης ασταθών εκτιμητών [79]. Η μέθοδος βασίζεται σε δύο θεμελιώδεις αρχές. Πρώτον, στο \emph{Bootstrap Sampling}: δημιουργία $B$ συνόλων εκπαίδευσης μέσω δειγματοληψίας με επανάθεση (sampling with replacement) από το αρχικό σύνολο μεγέθους $N$. Κάθε bootstrap sample περιέχει κατά μέσο όρο το $63.2\%$ των μοναδικών παρατηρήσεων του αρχικού συνόλου, όπως προκύπτει από τη σχέση $\lim_{N \to \infty} 1 - (1 - 1/N)^N = 1 - 1/e \approx 0.632$. Τα εναπομείναντα δείγματα, γνωστά ως \emph{Out-of-Bag (OOB) samples}, χρησιμοποιούνται ως εσωτερικό σύνολο επικύρωσης. Δεύτερον, στο \emph{Aggregation}: εκπαίδευση ενός μοντέλου βάσης (base learner) σε κάθε bootstrap sample και συνδυασμός των προβλέψεων μέσω πλειοψηφικής ψήφου για ταξινόμηση ή μέσου όρου για παλινδρόμηση,
$$\hat{y}_{\text{bag}}(x) = \text{mode}\{f_1(x), f_2(x), \ldots, f_B(x)\}$$
Για $B$ ανεξάρτητους εκτιμητές με διακύμανση $\sigma^2$, η διακύμανση του μέσου τους είναι $\sigma^2 / B$. Στην πράξη, οι εκτιμητές εμφανίζουν συσχέτιση $\rho$, οπότε η διακύμανση του ensemble δίνεται από
$$\text{Var}(\hat{y}_{\text{bag}}) = \rho \sigma^2 + \frac{1 - \rho}{B} \sigma^2$$
Ο δεύτερος όρος μειώνεται καθώς αυξάνεται το $B$, αλλά ο πρώτος παραμένει σταθερός. Η μείωση της συσχέτισης $\rho$ αποτελεί ακριβώς το κίνητρο για την επόμενη μέθοδο.

Το \textbf{Random Forest} (Breiman, 2001) επεκτείνει το Bagging με την εισαγωγή \emph{τυχαιότητας στην επιλογή χαρακτηριστικών} σε κάθε κόμβο διαχωρισμού [80]. Αντί να εξετάζονται όλα τα $P$ χαρακτηριστικά, σε κάθε split επιλέγεται ένα τυχαίο υποσύνολο μεγέθους $m$, όπου τυπικά $m = \lfloor \sqrt{P} \rfloor$ για classification και $m = \lfloor P/3 \rfloor$ για regression. Αυτός ο μηχανισμός μειώνει τη συσχέτιση μεταξύ των δέντρων: ακόμα και αν υπάρχει ένα πολύ ισχυρό χαρακτηριστικό (π.χ.~η MMP12 στο dataset μας), δεν θα χρησιμοποιείται σε κάθε δέντρο. Αυτό οδηγεί σε χαμηλότερη συνολική διακύμανση και καλύτερη γενίκευση. Ένα σημαντικό πλεονέκτημα του Random Forest, ιδιαίτερα στο πλαίσιο ανακάλυψης βιοδεικτών, είναι η δυνατότητα εκτίμησης της σημαντικότητας κάθε χαρακτηριστικού μέσω του δείκτη Mean Decrease in Impurity (MDI):
$$\text{Imp}(X_j) = \frac{1}{B} \sum_{b=1}^{B} \sum_{t \in T_b} \mathbb{I}(v(t) = j) \cdot p(t) \cdot \Delta\text{Gini}(t)$$
όπου $v(t)$ είναι το χαρακτηριστικό που χρησιμοποιείται στον κόμβο $t$, $p(t)$ η αναλογία δειγμάτων που φτάνουν στον κόμβο $t$, και $\Delta\text{Gini}(t)$ η μείωση ακαθαρσίας Gini λόγω του διαχωρισμού. Η μετρική αυτή παρέχει μια εναλλακτική, ανεξάρτητη μέθοδο ταυτοποίησης βιοδεικτών η οποία μπορεί να συγκριθεί με την ANOVA $F$-scores ανάλυση. Παράλληλα, το \emph{OOB error} προσφέρει εκτίμηση του σφάλματος γενίκευσης χωρίς ξεχωριστό validation set, ιδιαίτερα χρήσιμο σε datasets μικρού μεγέθους [83].

Πέρα από το κλασικό Random Forest, αναπτύχθηκαν και πιο σύνθετες μέθοδοι συνόλου που αυξάνουν περαιτέρω την ποικιλομορφία μεταξύ των δέντρων. Το \textbf{Rotation Forest} (Rodriguez et al., 2006) διαμερίζει τυχαία τα χαρακτηριστικά σε ομάδες, εφαρμόζει PCA σε κάθε ομάδα, και συνενώνει τα principal components ως νέα χαρακτηριστικά. Κάθε δέντρο, έτσι, ``βλέπει'' τα δεδομένα από μια εντελώς διαφορετική οπτική γωνία. Το \textbf{Random Subspace} εκπαιδεύει κάθε δέντρο σε ένα τυχαίο υποσύνολο όλων των χαρακτηριστικών (όχι μόνο σε κάθε κόμβο όπως το RF), αυξάνοντας τη συνολική ποικιλομορφία.

Μια ιδιαίτερα ενδιαφέρουσα πρόσφατη εξέλιξη είναι το \textbf{Deep Forest} (Zhou \& Feng, 2017), επίσης γνωστό ως \emph{gcForest}. Οι συγγραφείς παρατήρησαν ότι η επιτυχία της Βαθιάς Μάθησης βασίζεται σε τρεις βασικές αρχές: (α) επεξεργασία κατά στρώματα (layer-by-layer processing), (β) μετασχηματισμός χαρακτηριστικών εντός κάθε στρώματος (in-model feature transformation), και (γ) επαρκής πολυπλοκότητα μοντέλου. Αυτές οι αρχές δεν είναι, όμως, αποκλειστικά συνδεδεμένες με τα νευρωνικά δίκτυα. Το Deep Forest ενσωματώνει τις ίδιες αρχές χρησιμοποιώντας αποκλειστικά δέντρα αποφάσεων και μεθόδους συνόλου, χωρίς backpropagation [81], [82]. Τα κύρια πλεονεκτήματα είναι σημαντικά λιγότερες υπερπαράμετροι (δεν απαιτείται ρύθμιση learning rate, batch size ή αρχιτεκτονικής), αυτόματη ρύθμιση βάθους (μέσω early stopping βάσει cross-validation), και εξαιρετική απόδοση σε μικρά datasets (σχεδιασμένο για $N < 1000$). Η αρχιτεκτονική αποτελείται από μια σειρά επιπέδων (cascade) όπου κάθε επίπεδο περιέχει πολλαπλούς και διαφορετικούς τύπους forests, τυπικά Random Forests, Completely Random Forests (στα οποία τόσο τα χαρακτηριστικά όσο και τα κατώφλια επιλέγονται τυχαία) και Extra Trees [84]. Έστω $x \in \mathbb{R}^d$ το αρχικό διάνυσμα εισόδου και $C$ ο αριθμός κλάσεων. Σε κάθε επίπεδο $l$ του cascade, κάθε forest $f$ ($f = 1, \ldots, F$) παράγει ένα class probability vector $p_f^{(l)} \in \mathbb{R}^C$ μέσω εσωτερικού $k$-fold cross-validation στο σύνολο εκπαίδευσης (αυτό αποτρέπει information leakage μεταξύ επιπέδων). Τα probability vectors συνενώνονται με το αρχικό διάνυσμα εισόδου, σχηματίζοντας ένα augmented feature vector
$$x^{(l+1)} = \left[ x \; ; \; \underset{f=1}{\overset{F}{\text{concat}}} \; p_f^{(l)} \right] \in \mathbb{R}^{d + F \cdot C}$$
και η διαδικασία επαναλαμβάνεται. Η προσθήκη νέων επιπέδων σταματά αυτόματα όταν η απόδοση στο validation set δεν βελτιώνεται, αποφεύγοντας την υπερ-προσαρμογή. Στο τελικό επίπεδο $L$, η πρόβλεψη υπολογίζεται ως ο μέσος όρος των class probability vectors:
$$\hat{y} = \arg\max_c \frac{1}{F} \sum_{f=1}^{F} p_{f,c}^{(L)}$$

Η επιλογή μεθοδολογίας, νευρωνικά δίκτυα έναντι μεθόδων συνόλου, εξαρτάται κρίσιμα από τα χαρακτηριστικά του dataset. Συνοπτικά, τα tree-based μοντέλα απαιτούν τυπικά 50--100 δείγματα, παρουσιάζουν χαμηλό κίνδυνο overfitting, δεν χρειάζονται κανονικοποίηση δεδομένων, προσφέρουν υψηλή ερμηνευσιμότητα μέσω Feature Importance, αλλά μοντελοποιούν χρονικές δομές μόνο χειροκίνητα (μέσω feature engineering). Τα νευρωνικά δίκτυα (RNN/LSTM), αντιθέτως, τυπικά απαιτούν $\sim 1000+$ δείγματα, έχουν υψηλό κίνδυνο overfitting, απαιτούν κανονικοποίηση, λειτουργούν ως black box (χαμηλή ερμηνευσιμότητα χωρίς attention), αλλά μοντελοποιούν χρονικές δομές εγγενώς. Πρόσφατες μελέτες [85] έχουν αναδείξει ότι σε ένα ευρύ φάσμα δομημένων (tabular) datasets με $N < 10{,}000$, τα tree-based models (ιδίως XGBoost και Random Forest) υπερτερούν σταθερά των νευρωνικών δικτύων, εύρημα ιδιαίτερα σχετικό με τη φύση του πρωτεομικού dataset της παρούσας μελέτης.

\input{chapters/chapter3_methodology}
\input{chapters/chapter4_experiments}
\input{chapters/chapter5_discussion}
\input{chapters/chapter6_conclusions}

% =====================================================================
% ΒΙΒΛΙΟΓΡΑΦΙΑ
% =====================================================================
\input{chapters/references}

% =====================================================================
% ΠΑΡΑΡΤΗΜΑ
% =====================================================================
\section{Παράρτημα Α: Πλήρης Κώδικας Υλοποίησης}
 
Το παρόν παράρτημα περιέχει τον πλήρη κώδικα Python που αναπτύχθηκε στο πλαίσιο της διατριβής. Ο κώδικας οργανώνεται σε πέντε Jupyter Notebooks, τα οποία παρουσιάζονται στις ακόλουθες πέντε υποενότητες.
 
\textbf{Περιβάλλον εκτέλεσης:}
\begin{itemize}
\item Python 3.12
\item \texttt{numpy}, \texttt{pandas}, \texttt{scipy}: Βασικές βιβλιοθήκες
\item \texttt{scikit-learn}: Tree-based models, evaluation
\item \texttt{pytorch}: Deep learning (RNN, LSTM)
\item \texttt{xgboost}, \texttt{lightgbm}: Gradient boosting
\item \texttt{imbalanced-learn}: SMOTE
\item \texttt{matplotlib}, \texttt{seaborn}: Visualization
\end{itemize}
 
Για πλήρη αναπαραγωγιμότητα, όλα τα πειράματα χρησιμοποιούν σταθερό \texttt{random\_state=42}.
 
 
 

\subsection{A.1 Vanilla RNN με Attention}
 
\noindent \textbf{Πηγή:} \texttt{Vanilla\_RNN\_Temporal\_Protein.ipynb}
 
 
\subsubsection*{Temporal Protein Classification with Vanilla RNN}
 
\subsubsection*{ΜΕΡΟΣ 1: SETUP \& IMPORTS}
 
\begin{lstlisting}
# Imports
import numpy as np
import pandas as pd
import torch
import torch.nn as nn
import torch.nn.functional as F
from torch.utils.data import Dataset, DataLoader
from sklearn.model_selection import StratifiedKFold
from sklearn.preprocessing import StandardScaler, LabelEncoder
from sklearn.feature_selection import SelectKBest, f_classif
from sklearn.impute import SimpleImputer
from sklearn.metrics import classification_report, confusion_matrix, accuracy_score
import matplotlib.pyplot as plt
import seaborn as sns
import warnings
warnings.filterwarnings('ignore')
 
# Device selection
device = torch.device('cuda' if torch.cuda.is_available() else 'cpu')
print(f"="*50)
print(f"Using device: {device}")
if torch.cuda.is_available():
    print(f"GPU: {torch.cuda.get_device_name(0)}")
print(f"="*50)
 
# Reproducibility
def set_seed(seed=42):
    """Set seed for reproducibility"""
    np.random.seed(seed)
    torch.manual_seed(seed)
    if torch.cuda.is_available():
        torch.cuda.manual_seed_all(seed)
    torch.backends.cudnn.deterministic = True
    torch.backends.cudnn.benchmark = False
 
set_seed(42)
print("\n✓ Setup complete!")
\end{lstlisting}
 
\subsubsection*{ΜΕΡΟΣ 2: CONFIGURATION}
 
\begin{lstlisting}
class Config:
    """
    Configuration class - Όλες οι hyperparameters εδώ
    """
    # === Random Seed ===
    RANDOM_SEED = 42
 
    # === Data ===
    N_FOLDS = 5                    # Cross-validation folds
    N_TOP_FEATURES = 50            # Πόσες πρωτεΐνες να κρατήσουμε
 
    # === Model Architecture ===
    # CNN
    CNN_FILTERS = 32               # Αριθμός filters στο CNN
    CNN_KERNEL_SIZE = 5            # Kernel size
 
    # RNN
    RNN_HIDDEN_SIZE = 32           # Hidden size του RNN
    RNN_NUM_LAYERS = 1             # Αριθμός RNN layers
    RNN_BIDIRECTIONAL = True       # Bidirectional RNN
 
    # Attention
    ATTENTION_DIM = 32             # Attention dimension
 
    # Regularization
    DROPOUT = 0.5                  # Dropout rate
 
    # === Training ===
    BATCH_SIZE = 8                 # Batch size
    LEARNING_RATE = 1e-4           # Learning rate (μικρό για stability)
    WEIGHT_DECAY = 1e-3            # L2 regularization
    EPOCHS = 100                   # Max epochs
    PATIENCE = 15                  # Early stopping patience
 
    # === Device ===
    DEVICE = device
 
# Print configuration
print("Configuration:")
print(f"  Features: {Config.N_TOP_FEATURES}")
print(f"  RNN Hidden: {Config.RNN_HIDDEN_SIZE}")
print(f"  Bidirectional: {Config.RNN_BIDIRECTIONAL}")
print(f"  Dropout: {Config.DROPOUT}")
print(f"  Learning Rate: {Config.LEARNING_RATE}")
print(f"  Cross-validation Folds: {Config.N_FOLDS}")
print("\n✓ Configuration loaded!")
\end{lstlisting}
 
\subsubsection*{ΜΕΡΟΣ 3: DATA LOADING}
 
\begin{lstlisting}
# ╔══════════════════════════════════════════════════════════════╗
# ║  ΑΛΛΑΞΕ ΑΥΤΑ ΤΑ PATHS ΣΤΑ ΔΙΚΑ ΣΟΥ ΑΡΧΕΙΑ                   ║
# ╚══════════════════════════════════════════════════════════════╝
 
SAMPLES_PATH = '/content/DIOPTRA_Samples_list_Oncology.xlsx'
PROTEINS_PATH = '/content/DIOPTRA_Protein_results_Oncology.xlsx'
\end{lstlisting}
 
\begin{lstlisting}
def load_data(samples_path, proteins_path, include_t00=True):
    """
    Φορτώνει και προετοιμάζει τα δεδομένα.
 
    Parameters:
    -----------
    samples_path : str
        Path στο αρχείο με τα sample metadata
    proteins_path : str
        Path στο αρχείο με τις τιμές πρωτεϊνών
    include_t00 : bool
        Αν θα συμπεριλάβουμε τα single-timepoint samples (T00)
 
    Returns:
    --------
    X : np.ndarray
        Shape (n_samples, 2, n_proteins)
    y : np.ndarray
        Labels
    protein_names : list
        Ονόματα πρωτεϊνών
    label_encoder : LabelEncoder
        Fitted encoder
    """
    print("Loading data...")
    print("="*50)
 
    # Load Excel files
    samples = pd.read_excel(samples_path)
    proteins = pd.read_excel(proteins_path)
 
    print(f"Samples file: {len(samples)} rows")
    print(f"Proteins file: {len(proteins)} rows, {len(proteins.columns)-1} proteins")
 
    # Merge data
    data = samples.merge(proteins, on='SampleID')
 
    # Get protein column names
    protein_cols = [c for c in proteins.columns if c != 'SampleID']
 
    # === Create paired samples (T0 → T1) ===
    print("\nCreating paired samples...")
 
    t0_data = data[data['Timepoint'] == 'T0'].copy()
    t1_data = data[data['Timepoint'] == 'T1'].copy()
 
    # Patient ID extraction
    t0_data['PatientID'] = t0_data['SampleID']
    t1_data['PatientID'] = t1_data['SampleID'].str.replace('-F', '', regex=False)
 
    # Merge T0 and T1
    paired = t0_data.merge(
        t1_data[['PatientID'] + protein_cols],
        on='PatientID',
        suffixes=('_T0', '_T1')
    )
    print(f"  Paired samples (T0→T1): {len(paired)}")
 
    # === Add T00 samples (single timepoint) ===
    if include_t00:
        t00_data = data[data['Timepoint'] == 'T00'].copy()
 
        # Duplicate as T0 and T1 (no temporal change)
        for col in protein_cols:
            t00_data[f'{col}_T0'] = t00_data[col]
            t00_data[f'{col}_T1'] = t00_data[col]
        t00_data['PatientID'] = t00_data['SampleID']
 
        # Combine
        cols_needed = ['PatientID', 'Group'] + [f'{c}_T0' for c in protein_cols] + [f'{c}_T1' for c in protein_cols]
        combined = pd.concat([paired[cols_needed], t00_data[cols_needed]], ignore_index=True)
        print(f"  Single timepoint (T00): {len(t00_data)}")
        print(f"  Total samples: {len(combined)}")
    else:
        combined = paired
 
    # === Remove NAA class (too few samples) ===
    combined = combined[combined['Group'] != 'NAA']
    print(f"  After removing NAA: {len(combined)}")
 
    # === Create numpy arrays ===
    X_T0 = combined[[f'{c}_T0' for c in protein_cols]].values
    X_T1 = combined[[f'{c}_T1' for c in protein_cols]].values
    X = np.stack([X_T0, X_T1], axis=1)  # Shape: (n_samples, 2, n_proteins)
 
    # === Encode labels ===
    label_encoder = LabelEncoder()
    y = label_encoder.fit_transform(combined['Group'].values)
 
    # === Print class distribution ===
    print(f"\nClass distribution:")
    for i, cls in enumerate(label_encoder.classes_):
        count = (y == i).sum()
        pct = count / len(y) * 100
        print(f"  {cls}: {count} samples ({pct:.1f}%)")
 
    print(f"\nFinal shapes:")
    print(f"  X: {X.shape} (samples, timepoints, proteins)")
    print(f"  y: {y.shape}")
 
    return X, y, protein_cols, label_encoder
\end{lstlisting}
 
\begin{lstlisting}
# Load the data
X_raw, y, protein_names, label_encoder = load_data(SAMPLES_PATH, PROTEINS_PATH)
class_names = list(label_encoder.classes_)
 
print("\n" + "="*50)
print("✓ Data loaded successfully!")
print(f"  Classes: {class_names}")
\end{lstlisting}
 
\subsubsection*{ΜΕΡΟΣ 4: DATA PREPROCESSING}
 
\subsubsection*{4.1 Handle NaN Values}
 
\begin{lstlisting}
def handle_nan_values(X, protein_names):
    """
    Χειρίζεται τα NaN values:
    1. Αφαιρεί πρωτεΐνες που έχουν μόνο NaN
    2. Κάνει imputation με median στα υπόλοιπα NaN
 
    Parameters:
    -----------
    X : np.ndarray
        Shape (n_samples, n_timepoints, n_proteins)
    protein_names : list
        Ονόματα πρωτεϊνών
 
    Returns:
    --------
    X_clean : np.ndarray
        Cleaned data
    clean_protein_names : list
        Updated protein names
    """
    print("Handling NaN values...")
    print("="*50)
 
    n_samples, n_timepoints, n_proteins = X.shape
    total_nan = np.isnan(X).sum()
    print(f"  Total NaN values: {total_nan}")
    print(f"  Original shape: {X.shape}")
 
    if total_nan == 0:
        print("  No NaN values found!")
        return X, protein_names
 
    # Step 1: Find proteins with ALL NaN values
    valid_protein_mask = ~np.all(np.isnan(X), axis=(0, 1))
    n_all_nan = (~valid_protein_mask).sum()
    print(f"  Proteins with ALL NaN (removed): {n_all_nan}")
 
    # Remove all-NaN proteins
    X = X[:, :, valid_protein_mask]
    protein_names = [p for p, valid in zip(protein_names, valid_protein_mask) if valid]
 
    print(f"  Shape after removing all-NaN proteins: {X.shape}")
 
    # Step 2: Impute remaining NaN with median
    n_samples, n_timepoints, n_proteins = X.shape
    X_imputed = np.zeros_like(X)
 
    imputer = SimpleImputer(strategy='median')
 
    for t in range(n_timepoints):
        X_imputed[:, t, :] = imputer.fit_transform(X[:, t, :])
 
    remaining_nan = np.isnan(X_imputed).sum()
    print(f"  NaN after imputation: {remaining_nan}")
    print(f"  Final shape: {X_imputed.shape}")
 
    return X_imputed, protein_names
\end{lstlisting}
 
\begin{lstlisting}
# Apply NaN handling
X_clean, protein_names = handle_nan_values(X_raw, protein_names)
print("\n✓ NaN handling complete!")
\end{lstlisting}
 
\subsubsection*{4.2 Feature Selection}
 
\begin{lstlisting}
def select_top_features(X, y, protein_names, n_features=50):
    """
    Επιλέγει τις top-K πρωτεΐνες με βάση ANOVA F-score.
 
    Γιατί feature selection:
    - 368 features με 62 samples = overfitting
    - Κρατάμε μόνο τις πιο discriminative
 
    Parameters:
    -----------
    X : np.ndarray
        Shape (n_samples, n_timepoints, n_proteins)
    y : np.ndarray
        Labels
    protein_names : list
        Ονόματα πρωτεϊνών
    n_features : int
        Πόσα features να κρατήσουμε
 
    Returns:
    --------
    X_selected : np.ndarray
        Selected features
    selected_proteins : list
        Names of selected proteins
    selector : SelectKBest
        Fitted selector
    """
    print("Feature Selection...")
    print("="*50)
 
    # Use T0 for feature selection
    X_t0 = X[:, 0, :]
 
    # Adjust n_features if needed
    n_available = X_t0.shape[1]
    n_features = min(n_features, n_available)
    print(f"  Available proteins: {n_available}")
    print(f"  Selecting top: {n_features}")
 
    # ANOVA F-test
    selector = SelectKBest(f_classif, k=n_features)
    selector.fit(X_t0, y)
 
    # Get selected indices
    selected_mask = selector.get_support()
    selected_indices = np.where(selected_mask)[0]
 
    # Apply to both timepoints
    X_selected = X[:, :, selected_indices]
 
    # Get selected protein names
    selected_proteins = [protein_names[i] for i in selected_indices]
 
    # Get scores for ranking
    scores = selector.scores_[selected_indices]
    protein_scores = sorted(zip(selected_proteins, scores), key=lambda x: x[1], reverse=True)
 
    print(f"\n  Top 10 discriminative proteins:")
    for i, (prot, score) in enumerate(protein_scores[:10]):
        print(f"    {i+1}. {prot:<15} (F-score: {score:.2f})")
 
    print(f"\n  Final shape: {X_selected.shape}")
 
    return X_selected, selected_proteins, selector, selected_indices
\end{lstlisting}
 
\begin{lstlisting}
# Apply feature selection
X_selected, selected_proteins, feature_selector, selected_indices = select_top_features(
    X_clean, y, protein_names, n_features=Config.N_TOP_FEATURES
)
print("\n✓ Feature selection complete!")
\end{lstlisting}
 
\subsubsection*{4.3 Normalization}
 
\begin{lstlisting}
def normalize_data(X):
    """
    Κανονικοποιεί τα δεδομένα (mean=0, std=1).
 
    Parameters:
    -----------
    X : np.ndarray
        Shape (n_samples, n_timepoints, n_features)
 
    Returns:
    --------
    X_normalized : np.ndarray
        Normalized data
    scaler : StandardScaler
        Fitted scaler
    """
    print("Normalizing data...")
    print("="*50)
 
    n_samples, n_timepoints, n_features = X.shape
 
    # Flatten to 2D for normalization
    X_flat = X.reshape(-1, n_features)
 
    # Normalize
    scaler = StandardScaler()
    X_flat_scaled = scaler.fit_transform(X_flat)
 
    # Clip extreme values (για numerical stability)
    X_flat_scaled = np.clip(X_flat_scaled, -5, 5)
 
    # Reshape back
    X_normalized = X_flat_scaled.reshape(n_samples, n_timepoints, n_features)
 
    print(f"  Original range: [{X.min():.2f}, {X.max():.2f}]")
    print(f"  Normalized range: [{X_normalized.min():.2f}, {X_normalized.max():.2f}]")
    print(f"  Mean: {X_normalized.mean():.4f}")
    print(f"  Std: {X_normalized.std():.4f}")
 
    return X_normalized, scaler
\end{lstlisting}
 
\begin{lstlisting}
# Apply normalization
X, scaler = normalize_data(X_selected)
print("\n✓ Normalization complete!")
print(f"\nFinal data shape: {X.shape}")
print(f"Any NaN: {np.isnan(X).any()}")
\end{lstlisting}
 
\subsubsection*{ΜΕΡΟΣ 5: DATASET CLASS}
 
\begin{lstlisting}
class ProteinDataset(Dataset):
    """
    PyTorch Dataset για temporal protein data.
 
    Modes:
    - 'temporal': [T0, T1] - πλήρης χρονική πληροφορία
    - 'static': [T0] μόνο - baseline
    - 'delta': [T1-T0] - μόνο η αλλαγή
    """
 
    def __init__(self, X, y, mode='temporal'):
        """
        Parameters:
        -----------
        X : np.ndarray
            Shape (n_samples, 2, n_features)
        y : np.ndarray
            Labels
        mode : str
            'temporal', 'static', or 'delta'
        """
        self.mode = mode
        self.y = torch.LongTensor(y)
 
        if mode == 'temporal':
            # Both timepoints: (n_samples, 2, n_features)
            self.X = torch.FloatTensor(X)
        elif mode == 'static':
            # Only T0: (n_samples, 1, n_features)
            self.X = torch.FloatTensor(X[:, 0:1, :])
        elif mode == 'delta':
            # T1 - T0: (n_samples, 1, n_features)
            delta = X[:, 1, :] - X[:, 0, :]
            self.X = torch.FloatTensor(delta[:, np.newaxis, :])
        else:
            raise ValueError(f"Unknown mode: {mode}. Use 'temporal', 'static', or 'delta'")
 
    def __len__(self):
        return len(self.y)
 
    def __getitem__(self, idx):
        return self.X[idx], self.y[idx]
 
 
# Test dataset
print("Testing Dataset class...")
for mode in ['temporal', 'static', 'delta']:
    ds = ProteinDataset(X, y, mode=mode)
    x_sample, y_sample = ds[0]
    print(f"  {mode:10s}: X shape = {x_sample.shape}")
 
print("\n✓ Dataset class ready!")
\end{lstlisting}
 
\subsubsection*{ΜΕΡΟΣ 6: VANILLA RNN MODEL}
 
\subsubsection*{6.1 Vanilla RNN Cell (Educational Implementation)}
 
\begin{lstlisting}
class VanillaRNNCell(nn.Module):
    """
    Vanilla RNN Cell - Educational Implementation
 
    Εξίσωση:
    h_t = tanh(W_ih · x_t + W_hh · h_{t-1} + b)
 
    Όπου:
    - x_t: input στο χρόνο t
    - h_{t-1}: hidden state από το προηγούμενο βήμα
    - W_ih: weights για το input
    - W_hh: weights για το hidden state
    - b: bias
    - tanh: activation function
    """
 
    def __init__(self, input_size, hidden_size):
        super().__init__()
 
        self.input_size = input_size
        self.hidden_size = hidden_size
 
        # Weights
        self.W_ih = nn.Linear(input_size, hidden_size)   # Input to hidden
        self.W_hh = nn.Linear(hidden_size, hidden_size)  # Hidden to hidden
 
    def forward(self, x, h_prev):
        """
        Forward pass για ένα timestep.
 
        Parameters:
        -----------
        x : torch.Tensor
            Input at time t, shape (batch, input_size)
        h_prev : torch.Tensor
            Previous hidden state, shape (batch, hidden_size)
 
        Returns:
        --------
        h_new : torch.Tensor
            New hidden state, shape (batch, hidden_size)
        """
        # h_t = tanh(W_ih · x_t + W_hh · h_{t-1})
        h_new = torch.tanh(self.W_ih(x) + self.W_hh(h_prev))
 
        return h_new
 
 
# Test
print("Testing VanillaRNNCell...")
cell = VanillaRNNCell(input_size=32, hidden_size=16)
x = torch.randn(4, 32)  # batch=4, input=32
h = torch.zeros(4, 16)  # initial hidden state
h_new = cell(x, h)
print(f"  Input: {x.shape}")
print(f"  Hidden: {h.shape} → {h_new.shape}")
print("✓ VanillaRNNCell works!")
\end{lstlisting}
 
\subsubsection*{6.2 Full Vanilla RNN Layer}
 
\begin{lstlisting}
class VanillaRNN(nn.Module):
    """
    Complete Vanilla RNN Layer.
 
    Επεξεργάζεται ολόκληρη την ακολουθία και επιστρέφει
    τα hidden states για κάθε timestep.
    """
 
    def __init__(self, input_size, hidden_size, num_layers=1,
                 bidirectional=False, dropout=0.0):
        super().__init__()
 
        self.input_size = input_size
        self.hidden_size = hidden_size
        self.num_layers = num_layers
        self.bidirectional = bidirectional
        self.num_directions = 2 if bidirectional else 1
 
        # RNN Cells for each layer and direction
        self.cells_forward = nn.ModuleList()
        self.cells_backward = nn.ModuleList()
 
        for layer in range(num_layers):
            layer_input_size = input_size if layer == 0 else hidden_size * self.num_directions
 
            self.cells_forward.append(VanillaRNNCell(layer_input_size, hidden_size))
 
            if bidirectional:
                self.cells_backward.append(VanillaRNNCell(layer_input_size, hidden_size))
 
        self.dropout = nn.Dropout(dropout) if dropout > 0 else None
 
    def forward(self, x, h_0=None):
        """
        Forward pass.
 
        Parameters:
        -----------
        x : torch.Tensor
            Input sequence, shape (batch, seq_len, input_size)
        h_0 : torch.Tensor, optional
            Initial hidden state
 
        Returns:
        --------
        output : torch.Tensor
            Hidden states for all timesteps, shape (batch, seq_len, hidden_size * num_directions)
        h_n : torch.Tensor
            Final hidden state
        """
        batch_size, seq_len, _ = x.shape
 
        # Initialize hidden state
        if h_0 is None:
            h_0 = torch.zeros(self.num_layers * self.num_directions,
                             batch_size, self.hidden_size, device=x.device)
 
        # Process each layer
        current_input = x
        final_hidden = []
 
        for layer in range(self.num_layers):
            # Forward direction
            h_forward = h_0[layer * self.num_directions]
            outputs_forward = []
 
            for t in range(seq_len):
                h_forward = self.cells_forward[layer](current_input[:, t, :], h_forward)
                outputs_forward.append(h_forward)
 
            outputs_forward = torch.stack(outputs_forward, dim=1)  # (batch, seq_len, hidden)
            final_hidden.append(h_forward)
 
            # Backward direction
            if self.bidirectional:
                h_backward = h_0[layer * self.num_directions + 1]
                outputs_backward = []
 
                for t in range(seq_len - 1, -1, -1):
                    h_backward = self.cells_backward[layer](current_input[:, t, :], h_backward)
                    outputs_backward.append(h_backward)
 
                outputs_backward = torch.stack(outputs_backward[::-1], dim=1)
                final_hidden.append(h_backward)
 
                # Concatenate forward and backward
                layer_output = torch.cat([outputs_forward, outputs_backward], dim=2)
            else:
                layer_output = outputs_forward
 
            # Apply dropout between layers
            if self.dropout is not None and layer < self.num_layers - 1:
                layer_output = self.dropout(layer_output)
 
            current_input = layer_output
 
        output = current_input
        h_n = torch.stack(final_hidden, dim=0)
 
        return output, h_n
 
 
# Test
print("\nTesting VanillaRNN...")
rnn = VanillaRNN(input_size=32, hidden_size=16, bidirectional=True)
x = torch.randn(4, 10, 32)  # batch=4, seq_len=10, input=32
output, h_n = rnn(x)
print(f"  Input: {x.shape}")
print(f"  Output: {output.shape}")
print(f"  Final hidden: {h_n.shape}")
print("✓ VanillaRNN works!")
\end{lstlisting}
 
\subsubsection*{6.3 Attention Mechanism}
 
\begin{lstlisting}
class Attention(nn.Module):
    """
    Attention Mechanism.
 
    Υπολογίζει ποια positions είναι σημαντικά και
    δημιουργεί weighted sum.
 
    Πλεονέκτημα: Interpretability!
    Μπορούμε να δούμε ποιες πρωτεΐνες "κοιτάει" το model.
    """
 
    def __init__(self, input_dim, attention_dim):
        super().__init__()
 
        self.attention = nn.Sequential(
            nn.Linear(input_dim, attention_dim),
            nn.Tanh(),
            nn.Linear(attention_dim, 1)
        )
 
    def forward(self, x):
        """
        Parameters:
        -----------
        x : torch.Tensor
            Shape (batch, seq_len, input_dim)
 
        Returns:
        --------
        context : torch.Tensor
            Weighted sum, shape (batch, input_dim)
        weights : torch.Tensor
            Attention weights, shape (batch, seq_len)
        """
        # Calculate attention scores
        scores = self.attention(x)  # (batch, seq_len, 1)
 
        # Softmax to get weights
        weights = F.softmax(scores, dim=1)  # (batch, seq_len, 1)
 
        # Weighted sum
        context = (x * weights).sum(dim=1)  # (batch, input_dim)
 
        return context, weights.squeeze(-1)
 
 
# Test
print("\nTesting Attention...")
attn = Attention(input_dim=32, attention_dim=16)
x = torch.randn(4, 10, 32)
context, weights = attn(x)
print(f"  Input: {x.shape}")
print(f"  Context: {context.shape}")
print(f"  Weights: {weights.shape}")
print(f"  Weights sum: {weights[0].sum().item():.4f} (should be 1.0)")
print("✓ Attention works!")
\end{lstlisting}
 
\subsubsection*{6.4 Complete Model: CNN + Vanilla RNN + Attention}
 
\begin{lstlisting}
class TemporalProteinRNN(nn.Module):
    """
    Complete Model: 1D-CNN → Vanilla RNN → Attention → Classifier
 
    Architecture:
    1. CNN: Εξάγει local patterns από τις πρωτεΐνες
    2. RNN: Μοντελοποιεί την χρονική εξέλιξη
    3. Attention: Εστιάζει στα σημαντικά features
    4. Classifier: Τελική πρόβλεψη
    """
 
    def __init__(
        self,
        n_features,
        n_timepoints=2,
        cnn_filters=32,
        cnn_kernel_size=5,
        rnn_hidden_size=32,
        rnn_num_layers=1,
        rnn_bidirectional=True,
        attention_dim=32,
        n_classes=3,
        dropout=0.5,
        use_pytorch_rnn=False  # True για PyTorch built-in, False για custom
    ):
        super().__init__()
 
        self.n_features = n_features
        self.n_timepoints = n_timepoints
        self.use_pytorch_rnn = use_pytorch_rnn
 
        # ============ 1. CNN Feature Extractor ============
        self.cnn = nn.Sequential(
            nn.Conv1d(
                in_channels=n_timepoints,
                out_channels=cnn_filters,
                kernel_size=cnn_kernel_size,
                padding=cnn_kernel_size // 2
            ),
            nn.BatchNorm1d(cnn_filters),
            nn.ReLU(),
            nn.Dropout(dropout)
        )
 
        # ============ 2. Vanilla RNN ============
        if use_pytorch_rnn:
            # PyTorch built-in RNN
            self.rnn = nn.RNN(
                input_size=cnn_filters,
                hidden_size=rnn_hidden_size,
                num_layers=rnn_num_layers,
                batch_first=True,
                bidirectional=rnn_bidirectional,
                dropout=dropout if rnn_num_layers > 1 else 0
            )
        else:
            # Custom Vanilla RNN
            self.rnn = VanillaRNN(
                input_size=cnn_filters,
                hidden_size=rnn_hidden_size,
                num_layers=rnn_num_layers,
                bidirectional=rnn_bidirectional,
                dropout=dropout if rnn_num_layers > 1 else 0
            )
 
        # RNN output size
        rnn_output_size = rnn_hidden_size * (2 if rnn_bidirectional else 1)
 
        # ============ 3. Attention ============
        self.attention = Attention(rnn_output_size, attention_dim)
 
        # ============ 4. Classifier ============
        self.classifier = nn.Sequential(
            nn.Dropout(dropout),
            nn.Linear(rnn_output_size, n_classes)
        )
 
        # Store attention weights for interpretability
        self.attention_weights = None
 
    def forward(self, x):
        """
        Forward pass.
 
        Parameters:
        -----------
        x : torch.Tensor
            Shape (batch, n_timepoints, n_features)
 
        Returns:
        --------
        logits : torch.Tensor
            Shape (batch, n_classes)
        """
        # x: (batch, timepoints, features)
 
        # 1. CNN: Extract local patterns
        # CNN expects (batch, channels, length)
        cnn_out = self.cnn(x)  # (batch, cnn_filters, seq_len)
 
        # 2. RNN: Process sequence
        # RNN expects (batch, seq_len, features)
        rnn_in = cnn_out.transpose(1, 2)  # (batch, seq_len, cnn_filters)
        rnn_out, _ = self.rnn(rnn_in)  # (batch, seq_len, rnn_hidden * directions)
 
        # 3. Attention: Focus on important parts
        context, attn_weights = self.attention(rnn_out)  # (batch, rnn_hidden * directions)
        self.attention_weights = attn_weights  # Store for interpretability
 
        # 4. Classifier: Final prediction
        logits = self.classifier(context)  # (batch, n_classes)
 
        return logits
 
    def get_attention_weights(self):
        """Returns the last computed attention weights"""
        return self.attention_weights
\end{lstlisting}
 
\begin{lstlisting}
# Test the complete model
print("Testing Complete Model...")
print("="*50)
 
n_features = X.shape[2]
n_classes = len(class_names)
 
# Create model
model = TemporalProteinRNN(
    n_features=n_features,
    n_timepoints=2,
    cnn_filters=Config.CNN_FILTERS,
    cnn_kernel_size=Config.CNN_KERNEL_SIZE,
    rnn_hidden_size=Config.RNN_HIDDEN_SIZE,
    rnn_num_layers=Config.RNN_NUM_LAYERS,
    rnn_bidirectional=Config.RNN_BIDIRECTIONAL,
    attention_dim=Config.ATTENTION_DIM,
    n_classes=n_classes,
    dropout=Config.DROPOUT,
    use_pytorch_rnn=False  # Use our custom Vanilla RNN
)
 
# Count parameters
total_params = sum(p.numel() for p in model.parameters())
trainable_params = sum(p.numel() for p in model.parameters() if p.requires_grad)
 
print(f"\nModel Architecture:")
print(f"  Input: ({Config.BATCH_SIZE}, 2, {n_features})")
print(f"  CNN filters: {Config.CNN_FILTERS}")
print(f"  RNN hidden: {Config.RNN_HIDDEN_SIZE}")
print(f"  Bidirectional: {Config.RNN_BIDIRECTIONAL}")
print(f"  Output: ({Config.BATCH_SIZE}, {n_classes})")
print(f"\n  Total parameters: {total_params:,}")
print(f"  Trainable parameters: {trainable_params:,}")
print(f"  Samples / Parameters ratio: {len(y) / total_params:.3f}")
 
# Test forward pass
dummy_input = torch.randn(4, 2, n_features)
dummy_output = model(dummy_input)
print(f"\nForward pass test:")
print(f"  Input shape: {dummy_input.shape}")
print(f"  Output shape: {dummy_output.shape}")
print(f"  Attention weights shape: {model.get_attention_weights().shape}")
 
print("\n✓ Model ready!")
\end{lstlisting}
 
\subsubsection*{ΜΕΡΟΣ 7: TRAINING FUNCTIONS}
 
\begin{lstlisting}
def train_epoch(model, loader, optimizer, criterion, device):
    """
    Training για ένα epoch.
 
    Returns:
    --------
    avg_loss : float
    accuracy : float
    """
    model.train()
    total_loss = 0
    correct = 0
    total = 0
 
    for X_batch, y_batch in loader:
        X_batch = X_batch.to(device)
        y_batch = y_batch.to(device)
 
        # Forward pass
        optimizer.zero_grad()
        outputs = model(X_batch)
        loss = criterion(outputs, y_batch)
 
        # Check for NaN
        if torch.isnan(loss):
            print("WARNING: NaN loss!")
            continue
 
        # Backward pass
        loss.backward()
 
        # Gradient clipping (prevents exploding gradients)
        torch.nn.utils.clip_grad_norm_(model.parameters(), max_norm=1.0)
 
        # Update weights
        optimizer.step()
 
        # Statistics
        total_loss += loss.item()
        _, predicted = outputs.max(1)
        correct += (predicted == y_batch).sum().item()
        total += y_batch.size(0)
 
    avg_loss = total_loss / len(loader) if len(loader) > 0 else 0
    accuracy = correct / total if total > 0 else 0
 
    return avg_loss, accuracy
\end{lstlisting}
 
\begin{lstlisting}
def evaluate(model, loader, criterion, device):
    """
    Evaluation.
 
    Returns:
    --------
    avg_loss : float
    accuracy : float
    predictions : np.ndarray
    targets : np.ndarray
    """
    model.eval()
    total_loss = 0
    all_preds = []
    all_targets = []
 
    with torch.no_grad():
        for X_batch, y_batch in loader:
            X_batch = X_batch.to(device)
            y_batch = y_batch.to(device)
 
            # Forward pass
            outputs = model(X_batch)
            loss = criterion(outputs, y_batch)
 
            if not torch.isnan(loss):
                total_loss += loss.item()
 
            # Predictions
            _, predicted = outputs.max(1)
            all_preds.extend(predicted.cpu().numpy())
            all_targets.extend(y_batch.cpu().numpy())
 
    avg_loss = total_loss / len(loader) if len(loader) > 0 else 0
    accuracy = accuracy_score(all_targets, all_preds)
 
    return avg_loss, accuracy, np.array(all_preds), np.array(all_targets)
 
print("✓ Training functions defined!")
\end{lstlisting}
 
\subsubsection*{ΜΕΡΟΣ 8: CROSS-VALIDATION}
 
\begin{lstlisting}
def run_cross_validation(X, y, mode, config, n_folds=5, use_pytorch_rnn=False):
    """
    Stratified K-Fold Cross-Validation.
 
    Parameters:
    -----------
    X : np.ndarray
        Data
    y : np.ndarray
        Labels
    mode : str
        'temporal', 'static', or 'delta'
    config : Config
        Configuration object
    n_folds : int
        Number of CV folds
    use_pytorch_rnn : bool
        Use PyTorch RNN or custom implementation
 
    Returns:
    --------
    results : dict
        Results including accuracy, predictions, etc.
    """
    print(f"\n{'='*60}")
    print(f"EXPERIMENT: {mode.upper()} MODEL")
    print(f"{'='*60}")
    print(f"Mode: {mode}")
    print(f"Folds: {n_folds}")
    print(f"RNN: {'PyTorch' if use_pytorch_rnn else 'Custom Vanilla'}")
 
    # Determine n_timepoints based on mode
    n_timepoints = 2 if mode == 'temporal' else 1
    n_features = X.shape[2]
    n_classes = len(np.unique(y))
 
    # Class weights for imbalanced data
    class_counts = np.bincount(y)
    class_weights = torch.FloatTensor(
        len(y) / (len(class_counts) * class_counts)
    ).to(config.DEVICE)
 
    # Cross-validation
    kfold = StratifiedKFold(n_splits=n_folds, shuffle=True, random_state=config.RANDOM_SEED)
 
    fold_results = []
    all_preds = []
    all_targets = []
    histories = []
 
    for fold, (train_idx, val_idx) in enumerate(kfold.split(X, y)):
        print(f"\n--- Fold {fold+1}/{n_folds} ---")
 
        # Split data
        X_train, X_val = X[train_idx], X[val_idx]
        y_train, y_val = y[train_idx], y[val_idx]
 
        # Create datasets
        train_dataset = ProteinDataset(X_train, y_train, mode=mode)
        val_dataset = ProteinDataset(X_val, y_val, mode=mode)
 
        train_loader = DataLoader(
            train_dataset,
            batch_size=config.BATCH_SIZE,
            shuffle=True
        )
        val_loader = DataLoader(
            val_dataset,
            batch_size=config.BATCH_SIZE
        )
 
        # Create model
        model = TemporalProteinRNN(
            n_features=n_features,
            n_timepoints=n_timepoints,
            cnn_filters=config.CNN_FILTERS,
            cnn_kernel_size=config.CNN_KERNEL_SIZE,
            rnn_hidden_size=config.RNN_HIDDEN_SIZE,
            rnn_num_layers=config.RNN_NUM_LAYERS,
            rnn_bidirectional=config.RNN_BIDIRECTIONAL,
            attention_dim=config.ATTENTION_DIM,
            n_classes=n_classes,
            dropout=config.DROPOUT,
            use_pytorch_rnn=use_pytorch_rnn
        ).to(config.DEVICE)
 
        # Loss and optimizer
        criterion = nn.CrossEntropyLoss(weight=class_weights)
        optimizer = torch.optim.AdamW(
            model.parameters(),
            lr=config.LEARNING_RATE,
            weight_decay=config.WEIGHT_DECAY
        )
        scheduler = torch.optim.lr_scheduler.ReduceLROnPlateau(
            optimizer, mode='min', factor=0.5, patience=5
        )
 
        # Training loop
        best_val_acc = 0
        best_model_state = None
        patience_counter = 0
        history = {'train_loss': [], 'train_acc': [], 'val_loss': [], 'val_acc': []}
 
        for epoch in range(config.EPOCHS):
            # Train
            train_loss, train_acc = train_epoch(
                model, train_loader, optimizer, criterion, config.DEVICE
            )
 
            # Validate
            val_loss, val_acc, _, _ = evaluate(
                model, val_loader, criterion, config.DEVICE
            )
 
            # Update scheduler
            scheduler.step(val_loss)
 
            # Save history
            history['train_loss'].append(train_loss)
            history['train_acc'].append(train_acc)
            history['val_loss'].append(val_loss)
            history['val_acc'].append(val_acc)
 
            # Check for improvement
            if val_acc > best_val_acc:
                best_val_acc = val_acc
                best_model_state = {k: v.cpu().clone() for k, v in model.state_dict().items()}
                patience_counter = 0
            else:
                patience_counter += 1
 
            # Early stopping
            if patience_counter >= config.PATIENCE:
                print(f"  Early stopping at epoch {epoch+1}")
                break
 
            # Print progress
            if (epoch + 1) % 20 == 0:
                print(f"  Epoch {epoch+1}: Train Acc={train_acc:.4f}, Val Acc={val_acc:.4f}")
 
        # Load best model and evaluate
        if best_model_state is not None:
            model.load_state_dict(best_model_state)
 
        _, val_acc, preds, targets = evaluate(
            model, val_loader, criterion, config.DEVICE
        )
 
        fold_results.append(val_acc)
        all_preds.extend(preds)
        all_targets.extend(targets)
        histories.append(history)
 
        print(f"  Fold {fold+1} Best Accuracy: {val_acc:.4f}")
 
    # Summary
    mean_acc = np.mean(fold_results)
    std_acc = np.std(fold_results)
 
    print(f"\n{'='*40}")
    print(f"RESULTS: {mode.upper()}")
    print(f"{'='*40}")
    print(f"Mean Accuracy: {mean_acc:.4f} ± {std_acc:.4f}")
    print(f"Per-fold: {[f'{acc:.3f}' for acc in fold_results]}")
 
    return {
        'mean_acc': mean_acc,
        'std_acc': std_acc,
        'fold_results': fold_results,
        'all_preds': np.array(all_preds),
        'all_targets': np.array(all_targets),
        'histories': histories
    }
\end{lstlisting}
 
\subsubsection*{ΜΕΡΟΣ 9: RUN EXPERIMENTS ��}
 
\begin{lstlisting}
# Run all experiments
print("\n" + "#"*60)
print("#" + " "*20 + "RUNNING EXPERIMENTS" + " "*19 + "#")
print("#"*60)
 
results = {}
 
for mode in ['static', 'delta', 'temporal']:
    results[mode] = run_cross_validation(
        X, y,
        mode=mode,
        config=Config,
        n_folds=Config.N_FOLDS,
        use_pytorch_rnn=False  # Use custom Vanilla RNN
    )
\end{lstlisting}
 
\subsubsection*{ΜΕΡΟΣ 10: RESULTS \& VISUALIZATION}
 
\begin{lstlisting}
# Summary Table
print("\n" + "="*60)
print("FINAL RESULTS COMPARISON")
print("="*60)
print(f"\n{'Model':<15} {'Mean Accuracy':<18} {'Std':<10} {'Per Fold'}")
print("-" * 70)
 
for mode in ['static', 'delta', 'temporal']:
    mean = results[mode]['mean_acc']
    std = results[mode]['std_acc']
    folds = [f"{x:.3f}" for x in results[mode]['fold_results']]
    print(f"{mode:<15} {mean:<18.4f} ±{std:<9.4f} {folds}")
 
# Random baseline
random_baseline = 1.0 / len(class_names)
print(f"\n{'random guess':<15} {random_baseline:<18.4f}")
 
# Conclusion
print("\n" + "="*60)
print("CONCLUSION")
print("="*60)
 
temporal_acc = results['temporal']['mean_acc']
static_acc = results['static']['mean_acc']
delta_acc = results['delta']['mean_acc']
 
best_mode = max(['static', 'delta', 'temporal'], key=lambda m: results[m]['mean_acc'])
print(f"\nBest Model: {best_mode.upper()} ({results[best_mode]['mean_acc']:.4f})")
 
if temporal_acc > static_acc + 0.02:
    diff = (temporal_acc - static_acc) * 100
    print(f"\n✓ HYPOTHESIS SUPPORTED!")
    print(f"  Temporal model outperforms Static by {diff:.1f}%")
    print(f"  → Temporal information adds diagnostic value!")
elif static_acc > temporal_acc + 0.02:
    print(f"\n⚠ Static model performs better")
    print(f"  → May need more data for temporal patterns")
else:
    print(f"\n~ Results are similar (within noise)")
    print(f"  → Need more data for definitive conclusions")
\end{lstlisting}
 
\begin{lstlisting}
# Visualization
fig, axes = plt.subplots(1, 2, figsize=(14, 5))
 
# 1. Bar plot comparison
modes = ['static', 'delta', 'temporal']
means = [results[m]['mean_acc'] for m in modes]
stds = [results[m]['std_acc'] for m in modes]
 
colors = ['#3498db', '#e74c3c', '#2ecc71']
x_pos = np.arange(len(modes))
 
bars = axes[0].bar(x_pos, means, yerr=stds, capsize=5, color=colors, alpha=0.8, edgecolor='black')
axes[0].set_xticks(x_pos)
axes[0].set_xticklabels(['Static\n(T0 only)', 'Delta\n(T1-T0)', 'Temporal\n([T0,T1])'])
axes[0].set_ylabel('Accuracy', fontsize=12)
axes[0].set_title('Vanilla RNN: Model Comparison\n(5-Fold Cross-Validation)', fontsize=14)
axes[0].set_ylim(0, 1)
axes[0].axhline(y=1/len(class_names), color='gray', linestyle='--', label='Random guess', alpha=0.7)
axes[0].legend(loc='upper right')
axes[0].grid(axis='y', alpha=0.3)
 
# Add values on bars
for bar, mean, std in zip(bars, means, stds):
    axes[0].text(bar.get_x() + bar.get_width()/2, mean + std + 0.03,
                 f'{mean:.3f}', ha='center', fontsize=11, fontweight='bold')
 
# 2. Confusion matrix for best model
best_mode = max(modes, key=lambda m: results[m]['mean_acc'])
cm = confusion_matrix(results[best_mode]['all_targets'], results[best_mode]['all_preds'])
 
sns.heatmap(cm, annot=True, fmt='d', cmap='Blues',
            xticklabels=class_names, yticklabels=class_names, ax=axes[1],
            annot_kws={"size": 14})
axes[1].set_xlabel('Predicted', fontsize=12)
axes[1].set_ylabel('True', fontsize=12)
axes[1].set_title(f'Confusion Matrix\n({best_mode.capitalize()} Model)', fontsize=14)
 
plt.tight_layout()
plt.show()
\end{lstlisting}
 
\begin{lstlisting}
# Classification Report
best_mode = max(modes, key=lambda m: results[m]['mean_acc'])
 
print(f"\nClassification Report ({best_mode.capitalize()} Model):")
print("="*55)
print(classification_report(
    results[best_mode]['all_targets'],
    results[best_mode]['all_preds'],
    target_names=class_names
))
\end{lstlisting}
 
\begin{lstlisting}
# Training History (from last fold)
fig, axes = plt.subplots(1, 3, figsize=(15, 4))
 
for idx, mode in enumerate(['static', 'delta', 'temporal']):
    history = results[mode]['histories'][-1]  # Last fold
 
    axes[idx].plot(history['train_acc'], label='Train', linewidth=2)
    axes[idx].plot(history['val_acc'], label='Validation', linewidth=2)
    axes[idx].set_xlabel('Epoch')
    axes[idx].set_ylabel('Accuracy')
    axes[idx].set_title(f'{mode.capitalize()} Model')
    axes[idx].legend()
    axes[idx].grid(True, alpha=0.3)
    axes[idx].set_ylim(0, 1)
 
plt.suptitle('Training History (Last Fold)', fontsize=14, y=1.02)
plt.tight_layout()
plt.show()
\end{lstlisting}
 
\subsubsection*{ΜΕΡΟΣ 11: TOP DISCRIMINATIVE PROTEINS}
 
\begin{lstlisting}
# Get protein importance from feature selection
scores = feature_selector.scores_[selected_indices]
protein_importance = sorted(
    zip(selected_proteins, scores),
    key=lambda x: x[1],
    reverse=True
)
 
print("\nTop 20 Most Discriminative Proteins:")
print("="*50)
for i, (protein, score) in enumerate(protein_importance[:20]):
    print(f"{i+1:2d}. {protein:<20} (F-score: {score:.2f})")
\end{lstlisting}
 
\begin{lstlisting}
# Visualize
top_n = 15
top_proteins = [p for p, _ in protein_importance[:top_n]]
top_scores = [s for _, s in protein_importance[:top_n]]
 
fig, ax = plt.subplots(figsize=(10, 6))
bars = ax.barh(range(top_n), top_scores[::-1], color='steelblue', alpha=0.8, edgecolor='black')
ax.set_yticks(range(top_n))
ax.set_yticklabels(top_proteins[::-1])
ax.set_xlabel('F-score (ANOVA)', fontsize=12)
ax.set_title('Top 15 Discriminative Proteins', fontsize=14)
ax.grid(axis='x', alpha=0.3)
 
plt.tight_layout()
plt.show()
\end{lstlisting}
 
\subsubsection*{ΜΕΡΟΣ 12: SAVE RESULTS}
 
\begin{lstlisting}
# Save results to CSV
results_df = pd.DataFrame({
    'Model': ['Static', 'Delta', 'Temporal'],
    'Mean_Accuracy': [results['static']['mean_acc'], results['delta']['mean_acc'], results['temporal']['mean_acc']],
    'Std_Accuracy': [results['static']['std_acc'], results['delta']['std_acc'], results['temporal']['std_acc']],
})
 
# Add fold results
for i in range(Config.N_FOLDS):
    results_df[f'Fold_{i+1}'] = [
        results['static']['fold_results'][i],
        results['delta']['fold_results'][i],
        results['temporal']['fold_results'][i]
    ]
 
results_df.to_csv('vanilla_rnn_results.csv', index=False)
print("Results saved to 'vanilla_rnn_results.csv'")
display(results_df)
\end{lstlisting}
 
\begin{lstlisting}
# Save top proteins
proteins_df = pd.DataFrame({
    'Rank': range(1, len(protein_importance) + 1),
    'Protein': [p for p, _ in protein_importance],
    'F_score': [s for _, s in protein_importance]
})
 
proteins_df.to_csv('top_proteins_vanilla_rnn.csv', index=False)
print("\nTop proteins saved to 'top_proteins_vanilla_rnn.csv'")
display(proteins_df.head(15))
\end{lstlisting}
 
\subsubsection*{SUMMARY}
 

\subsection{A.2 LSTM Implementation}
 
\noindent \textbf{Πηγή:} \texttt{Protein\_Network\_FIXED.ipynb}
 
 
\subsubsection*{Temporal Protein Network - FIXED VERSION}
 
\subsubsection*{1. Setup}
 
\begin{lstlisting}
import numpy as np
import pandas as pd
import torch
import torch.nn as nn
import torch.nn.functional as F
from torch.utils.data import Dataset, DataLoader
from sklearn.model_selection import StratifiedKFold
from sklearn.preprocessing import StandardScaler, LabelEncoder
from sklearn.feature_selection import SelectKBest, f_classif
from sklearn.impute import SimpleImputer
from sklearn.metrics import classification_report, confusion_matrix, accuracy_score
import matplotlib.pyplot as plt
import seaborn as sns
import warnings
warnings.filterwarnings('ignore')
 
# Device
device = torch.device('cuda' if torch.cuda.is_available() else 'cpu')
print(f"Using device: {device}")
 
# Reproducibility
def set_seed(seed=42):
    np.random.seed(seed)
    torch.manual_seed(seed)
    if torch.cuda.is_available():
        torch.cuda.manual_seed_all(seed)
 
set_seed(42)
print("Setup complete ✓")
\end{lstlisting}
 
\subsubsection*{2. Configuration}
 
\begin{lstlisting}
class Config:
    RANDOM_SEED = 42
    N_FOLDS = 5
    N_TOP_FEATURES = 50
 
    # Model - Small
    HIDDEN_DIM = 32
    DROPOUT = 0.5
 
    # Training - Conservative
    BATCH_SIZE = 8
    LEARNING_RATE = 1e-4
    WEIGHT_DECAY = 1e-3
    EPOCHS = 50
    PATIENCE = 10
 
    DEVICE = device
 
print("Config loaded ✓")
\end{lstlisting}
 
\subsubsection*{3. Data Loading}
 
\begin{lstlisting}
# ============================================================
# ΑΛΛΑΞΕ ΑΥΤΑ ΤΑ PATHS ΣΤΑ ΔΙΚΑ ΣΟΥ ΑΡΧΕΙΑ
# ============================================================
SAMPLES_PATH = '/content/DIOPTRA_Samples_list_Oncology.xlsx'
PROTEINS_PATH = '/content/DIOPTRA_Protein_results_Oncology.xlsx'
# ============================================================
\end{lstlisting}
 
\begin{lstlisting}
def load_data(samples_path, proteins_path, include_t00=True):
    """
    Φορτώνει τα δεδομένα.
    """
    samples = pd.read_excel(samples_path)
    proteins = pd.read_excel(proteins_path)
    data = samples.merge(proteins, on='SampleID')
 
    protein_cols = [c for c in proteins.columns if c != 'SampleID']
 
    # Paired samples (T0 → T1)
    t0_data = data[data['Timepoint'] == 'T0'].copy()
    t1_data = data[data['Timepoint'] == 'T1'].copy()
 
    t0_data['PatientID'] = t0_data['SampleID']
    t1_data['PatientID'] = t1_data['SampleID'].str.replace('-F', '', regex=False)
 
    paired = t0_data.merge(
        t1_data[['PatientID'] + protein_cols],
        on='PatientID',
        suffixes=('_T0', '_T1')
    )
    print(f"Paired samples: {len(paired)}")
 
    # Add T00 samples
    if include_t00:
        t00_data = data[data['Timepoint'] == 'T00'].copy()
        for col in protein_cols:
            t00_data[f'{col}_T0'] = t00_data[col]
            t00_data[f'{col}_T1'] = t00_data[col]
        t00_data['PatientID'] = t00_data['SampleID']
 
        cols_needed = ['PatientID', 'Group'] + [f'{c}_T0' for c in protein_cols] + [f'{c}_T1' for c in protein_cols]
        combined = pd.concat([paired[cols_needed], t00_data[cols_needed]], ignore_index=True)
        print(f"After adding T00: {len(combined)}")
    else:
        combined = paired
 
    # Remove NAA
    combined = combined[combined['Group'] != 'NAA']
    print(f"After removing NAA: {len(combined)}")
 
    # Create arrays
    X_T0 = combined[[f'{c}_T0' for c in protein_cols]].values
    X_T1 = combined[[f'{c}_T1' for c in protein_cols]].values
    X = np.stack([X_T0, X_T1], axis=1)  # (n_samples, 2, n_proteins)
 
    # Labels
    label_encoder = LabelEncoder()
    y = label_encoder.fit_transform(combined['Group'].values)
 
    print(f"\nClass distribution:")
    for i, cls in enumerate(label_encoder.classes_):
        print(f"  {cls}: {(y == i).sum()}")
 
    return X, y, protein_cols, label_encoder
 
 
# Load data
X_raw, y, protein_names, label_encoder = load_data(SAMPLES_PATH, PROTEINS_PATH)
class_names = list(label_encoder.classes_)
 
print(f"\nX shape: {X_raw.shape}")
print(f"Classes: {class_names}")
\end{lstlisting}
 
\subsubsection*{4. Handle NaN Values}
 
\begin{lstlisting}
# Check for NaN
print(f"NaN values in X_raw: {np.isnan(X_raw).sum()}")
print(f"Shape before: {X_raw.shape}")
 
n_samples, n_timepoints, n_proteins = X_raw.shape
 
# Step 1: Find proteins that have at least some valid values
valid_protein_mask = ~np.all(np.isnan(X_raw), axis=(0, 1))  # True if protein has any valid value
print(f"\nProteins with ALL NaN (will be removed): {(~valid_protein_mask).sum()}")
print(f"Valid proteins: {valid_protein_mask.sum()}")
 
# Step 2: Keep only valid proteins
X_raw = X_raw[:, :, valid_protein_mask]
protein_names = [p for p, valid in zip(protein_names, valid_protein_mask) if valid]
 
print(f"Shape after removing all-NaN proteins: {X_raw.shape}")
 
# Step 3: Impute remaining NaN with median
n_samples, n_timepoints, n_proteins = X_raw.shape
X_imputed = np.zeros_like(X_raw)
 
imputer = SimpleImputer(strategy='median')
 
for t in range(n_timepoints):
    X_imputed[:, t, :] = imputer.fit_transform(X_raw[:, t, :])
 
X_raw = X_imputed
 
print(f"\nNaN values after imputation: {np.isnan(X_raw).sum()}")
print(f"Final shape: {X_raw.shape}")
\end{lstlisting}
 
\subsubsection*{5. Feature Selection}
 
\begin{lstlisting}
def select_top_features(X, y, protein_names, n_features=50):
    """
    Επιλέγει τις top-K πρωτεΐνες με βάση ANOVA F-score.
    """
    # Use T0 for feature selection
    X_t0 = X[:, 0, :]
 
    # Adjust n_features if we have fewer proteins
    n_features = min(n_features, X_t0.shape[1])
 
    # Select top features
    selector = SelectKBest(f_classif, k=n_features)
    selector.fit(X_t0, y)
 
    # Get selected indices
    selected_mask = selector.get_support()
    selected_indices = np.where(selected_mask)[0]
 
    # Apply to both timepoints
    X_selected = X[:, :, selected_indices]
 
    # Get selected protein names
    selected_proteins = [protein_names[i] for i in selected_indices]
 
    return X_selected, selected_indices, selected_proteins, selector
 
 
# Apply feature selection
X, selected_indices, selected_proteins, feature_selector = select_top_features(
    X_raw, y, protein_names, n_features=Config.N_TOP_FEATURES
)
 
print(f"Selected {len(selected_proteins)} proteins")
print(f"New X shape: {X.shape}")
print(f"\nTop 10 selected proteins:")
for i, p in enumerate(selected_proteins[:10]):
    print(f"  {i+1}. {p}")
\end{lstlisting}
 
\subsubsection*{6. Normalize Data}
 
\begin{lstlisting}
def normalize_data(X):
    """Normalize per feature across all samples and timepoints"""
    n_samples, n_timepoints, n_features = X.shape
    X_flat = X.reshape(-1, n_features)
 
    scaler = StandardScaler()
    X_flat_scaled = scaler.fit_transform(X_flat)
 
    # Clip extreme values
    X_flat_scaled = np.clip(X_flat_scaled, -5, 5)
 
    return X_flat_scaled.reshape(n_samples, n_timepoints, n_features), scaler
 
 
X, scaler = normalize_data(X)
print(f"Data normalized.")
print(f"Range: [{X.min():.2f}, {X.max():.2f}]")
print(f"Any NaN: {np.isnan(X).any()}")
\end{lstlisting}
 
\subsubsection*{7. Dataset Class}
 
\begin{lstlisting}
class ProteinDataset(Dataset):
    def __init__(self, X, y, mode='temporal'):
        self.y = torch.LongTensor(y)
 
        if mode == 'temporal':
            self.X = torch.FloatTensor(X)
        elif mode == 'static':
            self.X = torch.FloatTensor(X[:, 0:1, :])
        elif mode == 'delta':
            delta = X[:, 1, :] - X[:, 0, :]
            self.X = torch.FloatTensor(delta[:, np.newaxis, :])
        else:
            raise ValueError(f"Unknown mode: {mode}")
 
    def __len__(self):
        return len(self.y)
 
    def __getitem__(self, idx):
        return self.X[idx], self.y[idx]
 
print("Dataset class defined ✓")
\end{lstlisting}
 
\subsubsection*{8. Model Architecture}
 
\begin{lstlisting}
class SimpleTemporalNetwork(nn.Module):
    """
    Simplified model for small datasets.
    """
 
    def __init__(self, n_features, n_timepoints=2, hidden_dim=32, n_classes=3, dropout=0.5):
        super().__init__()
 
        # Simple 1D Conv
        self.conv = nn.Sequential(
            nn.Conv1d(n_timepoints, hidden_dim, kernel_size=5, padding=2),
            nn.BatchNorm1d(hidden_dim),
            nn.ReLU(),
            nn.Dropout(dropout)
        )
 
        # Simple LSTM
        self.lstm = nn.LSTM(
            input_size=hidden_dim,
            hidden_size=hidden_dim // 2,
            num_layers=1,
            batch_first=True,
            bidirectional=True
        )
 
        # Attention
        self.attention = nn.Sequential(
            nn.Linear(hidden_dim, hidden_dim // 2),
            nn.Tanh(),
            nn.Linear(hidden_dim // 2, 1)
        )
 
        # Classifier
        self.classifier = nn.Sequential(
            nn.Dropout(dropout),
            nn.Linear(hidden_dim, n_classes)
        )
 
    def forward(self, x):
        # CNN
        conv_out = self.conv(x)
 
        # LSTM
        lstm_in = conv_out.transpose(1, 2)
        lstm_out, _ = self.lstm(lstm_in)
 
        # Attention
        attn_scores = self.attention(lstm_out)
        attn_weights = F.softmax(attn_scores, dim=1)
        context = (lstm_out * attn_weights).sum(dim=1)
 
        # Classify
        logits = self.classifier(context)
 
        return logits
 
 
# Test model
n_features = X.shape[2]
test_model = SimpleTemporalNetwork(n_features=n_features, n_classes=len(class_names))
 
n_params = sum(p.numel() for p in test_model.parameters())
print(f"Model parameters: {n_params:,}")
print(f"Samples: {len(y)}")
print(f"Ratio (samples/params): {len(y)/n_params:.2f}")
 
# Test forward
dummy = torch.randn(4, 2, n_features)
out = test_model(dummy)
print(f"Forward pass OK: {dummy.shape} → {out.shape}")
\end{lstlisting}
 
\subsubsection*{9. Training Functions}
 
\begin{lstlisting}
def train_epoch(model, loader, optimizer, criterion, device):
    model.train()
    total_loss = 0
    correct = 0
    total = 0
 
    for X_batch, y_batch in loader:
        X_batch, y_batch = X_batch.to(device), y_batch.to(device)
 
        optimizer.zero_grad()
        outputs = model(X_batch)
        loss = criterion(outputs, y_batch)
 
        if torch.isnan(loss):
            continue
 
        loss.backward()
        torch.nn.utils.clip_grad_norm_(model.parameters(), max_norm=1.0)
        optimizer.step()
 
        total_loss += loss.item()
        _, predicted = outputs.max(1)
        correct += (predicted == y_batch).sum().item()
        total += y_batch.size(0)
 
    avg_loss = total_loss / len(loader) if len(loader) > 0 else 0
    accuracy = correct / total if total > 0 else 0
    return avg_loss, accuracy
 
 
def evaluate(model, loader, criterion, device):
    model.eval()
    total_loss = 0
    all_preds = []
    all_targets = []
 
    with torch.no_grad():
        for X_batch, y_batch in loader:
            X_batch, y_batch = X_batch.to(device), y_batch.to(device)
            outputs = model(X_batch)
            loss = criterion(outputs, y_batch)
 
            if not torch.isnan(loss):
                total_loss += loss.item()
 
            _, predicted = outputs.max(1)
            all_preds.extend(predicted.cpu().numpy())
            all_targets.extend(y_batch.cpu().numpy())
 
    avg_loss = total_loss / len(loader) if len(loader) > 0 else 0
    accuracy = accuracy_score(all_targets, all_preds)
    return avg_loss, accuracy, np.array(all_preds), np.array(all_targets)
 
print("Training functions defined ✓")
\end{lstlisting}
 
\subsubsection*{10. Cross-Validation Function}
 
\begin{lstlisting}
def run_cross_validation(X, y, mode, config, n_folds=5):
    """
    Cross-validation για robust evaluation.
    """
    print(f"\n{'='*60}")
    print(f"EXPERIMENT: {mode.upper()} MODEL ({n_folds}-fold CV)")
    print(f"{'='*60}")
 
    n_timepoints = 2 if mode == 'temporal' else 1
    n_features = X.shape[2]
    n_classes = len(np.unique(y))
 
    # Class weights
    class_counts = np.bincount(y)
    class_weights = torch.FloatTensor(len(y) / (len(class_counts) * class_counts)).to(config.DEVICE)
 
    # Cross-validation
    kfold = StratifiedKFold(n_splits=n_folds, shuffle=True, random_state=config.RANDOM_SEED)
 
    fold_results = []
    all_preds = []
    all_targets = []
 
    for fold, (train_idx, val_idx) in enumerate(kfold.split(X, y)):
        print(f"\nFold {fold+1}/{n_folds}", end=" ")
 
        # Split
        X_train, X_val = X[train_idx], X[val_idx]
        y_train, y_val = y[train_idx], y[val_idx]
 
        # Datasets
        train_dataset = ProteinDataset(X_train, y_train, mode=mode)
        val_dataset = ProteinDataset(X_val, y_val, mode=mode)
 
        train_loader = DataLoader(train_dataset, batch_size=config.BATCH_SIZE, shuffle=True)
        val_loader = DataLoader(val_dataset, batch_size=config.BATCH_SIZE)
 
        # Model
        model = SimpleTemporalNetwork(
            n_features=n_features,
            n_timepoints=n_timepoints,
            hidden_dim=config.HIDDEN_DIM,
            n_classes=n_classes,
            dropout=config.DROPOUT
        ).to(config.DEVICE)
 
        # Training setup
        criterion = nn.CrossEntropyLoss(weight=class_weights)
        optimizer = torch.optim.AdamW(
            model.parameters(),
            lr=config.LEARNING_RATE,
            weight_decay=config.WEIGHT_DECAY
        )
        scheduler = torch.optim.lr_scheduler.ReduceLROnPlateau(
            optimizer, mode='min', factor=0.5, patience=5
        )
 
        # Training loop
        best_val_acc = 0
        best_model_state = None
        patience_counter = 0
 
        for epoch in range(config.EPOCHS):
            train_loss, train_acc = train_epoch(model, train_loader, optimizer, criterion, config.DEVICE)
            val_loss, val_acc, _, _ = evaluate(model, val_loader, criterion, config.DEVICE)
 
            scheduler.step(val_loss)
 
            if val_acc > best_val_acc:
                best_val_acc = val_acc
                best_model_state = {k: v.cpu().clone() for k, v in model.state_dict().items()}
                patience_counter = 0
            else:
                patience_counter += 1
 
            if patience_counter >= config.PATIENCE:
                break
 
        # Evaluate best model
        if best_model_state is not None:
            model.load_state_dict(best_model_state)
        _, val_acc, preds, targets = evaluate(model, val_loader, criterion, config.DEVICE)
 
        fold_results.append(val_acc)
        all_preds.extend(preds)
        all_targets.extend(targets)
 
        print(f"→ Accuracy: {val_acc:.4f}")
 
    # Summary
    mean_acc = np.mean(fold_results)
    std_acc = np.std(fold_results)
 
    print(f"\n" + "-"*40)
    print(f"Mean Accuracy: {mean_acc:.4f} ± {std_acc:.4f}")
 
    return {
        'mean_acc': mean_acc,
        'std_acc': std_acc,
        'fold_results': fold_results,
        'all_preds': np.array(all_preds),
        'all_targets': np.array(all_targets)
    }
\end{lstlisting}
 
\subsubsection*{11. Run All Experiments ��}
 
\begin{lstlisting}
# Run all three experiments
results = {}
 
for mode in ['static', 'delta', 'temporal']:
    results[mode] = run_cross_validation(X, y, mode, Config, n_folds=Config.N_FOLDS)
\end{lstlisting}
 
\subsubsection*{12. Results Comparison ��}
 
\begin{lstlisting}
# Summary table
print("\n" + "="*60)
print("FINAL COMPARISON (5-Fold Cross-Validation)")
print("="*60)
print(f"\n{'Model':<15} {'Mean Accuracy':<18} {'Std':<10}")
print("-" * 45)
 
for mode in ['static', 'delta', 'temporal']:
    mean = results[mode]['mean_acc']
    std = results[mode]['std_acc']
    print(f"{mode:<15} {mean:<18.4f} ±{std:<10.4f}")
 
# Baseline comparison
random_baseline = 1.0 / len(class_names)
print(f"\n{'random guess':<15} {random_baseline:<18.4f}")
 
# Conclusion
print("\n" + "-"*45)
temporal_acc = results['temporal']['mean_acc']
static_acc = results['static']['mean_acc']
delta_acc = results['delta']['mean_acc']
 
best_mode = max(['static', 'delta', 'temporal'], key=lambda m: results[m]['mean_acc'])
print(f"Best model: {best_mode.upper()} ({results[best_mode]['mean_acc']:.4f})")
 
if temporal_acc > static_acc + 0.02:
    improvement = (temporal_acc - static_acc) * 100
    print(f"\n✓ Temporal > Static by {improvement:.1f}%")
    print("  → Evidence that temporal info helps!")
elif static_acc > temporal_acc + 0.02:
    print("\n⚠ Static performs better than Temporal")
    print("  → May need more data for temporal patterns")
else:
    print("\n~ Results are similar (within noise)")
    print("  → Need more data for definitive conclusions")
\end{lstlisting}
 
\begin{lstlisting}
# Visualization
fig, axes = plt.subplots(1, 2, figsize=(14, 5))
 
# Bar plot comparison
modes = ['static', 'delta', 'temporal']
means = [results[m]['mean_acc'] for m in modes]
stds = [results[m]['std_acc'] for m in modes]
 
colors = ['#3498db', '#e74c3c', '#2ecc71']
bars = axes[0].bar(modes, means, yerr=stds, capsize=5, color=colors, alpha=0.8)
axes[0].set_ylabel('Accuracy', fontsize=12)
axes[0].set_title('Model Comparison (5-Fold CV)', fontsize=14)
axes[0].set_ylim(0, 1)
axes[0].axhline(y=1/len(class_names), color='gray', linestyle='--', label='Random guess')
axes[0].legend()
axes[0].grid(axis='y', alpha=0.3)
 
# Add values on bars
for bar, mean, std in zip(bars, means, stds):
    axes[0].text(bar.get_x() + bar.get_width()/2, mean + std + 0.03,
                 f'{mean:.3f}', ha='center', fontsize=11, fontweight='bold')
 
# Confusion matrix for best model
best_mode = max(modes, key=lambda m: results[m]['mean_acc'])
cm = confusion_matrix(results[best_mode]['all_targets'], results[best_mode]['all_preds'])
 
sns.heatmap(cm, annot=True, fmt='d', cmap='Blues',
            xticklabels=class_names, yticklabels=class_names, ax=axes[1],
            annot_kws={"size": 14})
axes[1].set_xlabel('Predicted', fontsize=12)
axes[1].set_ylabel('True', fontsize=12)
axes[1].set_title(f'Confusion Matrix ({best_mode.capitalize()} Model)', fontsize=14)
 
plt.tight_layout()
plt.show()
\end{lstlisting}
 
\begin{lstlisting}
# Classification report for best model
best_mode = max(modes, key=lambda m: results[m]['mean_acc'])
 
print(f"\nClassification Report ({best_mode.capitalize()} Model):")
print("="*50)
print(classification_report(
    results[best_mode]['all_targets'],
    results[best_mode]['all_preds'],
    target_names=class_names
))
\end{lstlisting}
 
\subsubsection*{13. Top Discriminative Proteins}
 
\begin{lstlisting}
# Get protein importance scores
scores = feature_selector.scores_[selected_indices]
protein_importance = sorted(zip(selected_proteins, scores), key=lambda x: x[1], reverse=True)
 
print("\nTop 20 Most Discriminative Proteins:")
print("="*45)
for i, (protein, score) in enumerate(protein_importance[:20]):
    print(f"{i+1:2d}. {protein:<20} (F-score: {score:.2f})")
\end{lstlisting}
 
\begin{lstlisting}
# Visualize top proteins
top_n = 15
top_proteins = [p for p, _ in protein_importance[:top_n]]
top_scores = [s for _, s in protein_importance[:top_n]]
 
fig, ax = plt.subplots(figsize=(10, 6))
bars = ax.barh(range(top_n), top_scores[::-1], color='steelblue', alpha=0.8)
ax.set_yticks(range(top_n))
ax.set_yticklabels(top_proteins[::-1])
ax.set_xlabel('F-score (ANOVA)', fontsize=12)
ax.set_title('Top 15 Discriminative Proteins', fontsize=14)
ax.grid(axis='x', alpha=0.3)
 
plt.tight_layout()
plt.show()
\end{lstlisting}
 
\subsubsection*{14. Save Results}
 
\begin{lstlisting}
# Save results to CSV
results_df = pd.DataFrame({
    'Model': ['Static', 'Delta', 'Temporal'],
    'Mean_Accuracy': [results['static']['mean_acc'], results['delta']['mean_acc'], results['temporal']['mean_acc']],
    'Std_Accuracy': [results['static']['std_acc'], results['delta']['std_acc'], results['temporal']['std_acc']],
    'Fold_1': [results['static']['fold_results'][0], results['delta']['fold_results'][0], results['temporal']['fold_results'][0]],
    'Fold_2': [results['static']['fold_results'][1], results['delta']['fold_results'][1], results['temporal']['fold_results'][1]],
    'Fold_3': [results['static']['fold_results'][2], results['delta']['fold_results'][2], results['temporal']['fold_results'][2]],
    'Fold_4': [results['static']['fold_results'][3], results['delta']['fold_results'][3], results['temporal']['fold_results'][3]],
    'Fold_5': [results['static']['fold_results'][4], results['delta']['fold_results'][4], results['temporal']['fold_results'][4]],
})
 
results_df.to_csv('experiment_results.csv', index=False)
print("Results saved to 'experiment_results.csv'")
display(results_df)
\end{lstlisting}
 
\begin{lstlisting}
# Save top proteins
proteins_df = pd.DataFrame({
    'Rank': range(1, len(protein_importance) + 1),
    'Protein': [p for p, _ in protein_importance],
    'F_score': [s for _, s in protein_importance]
})
 
proteins_df.to_csv('top_proteins.csv', index=False)
print("Top proteins saved to 'top_proteins.csv'")
display(proteins_df.head(20))
\end{lstlisting}
 
\subsubsection*{Summary}
 

\subsection{A.3 Ensemble Methods \& Custom Cascade Forest}
 
\noindent \textbf{Πηγή:} \texttt{Ensemble\_DeepForest\_Experiments.ipynb}
 
 
\subsubsection*{Ensemble \& Deep Forest Experiments}
 
\begin{lstlisting}
# === Imports ===
import numpy as np
import pandas as pd
import warnings
warnings.filterwarnings('ignore')
 
from sklearn.model_selection import StratifiedKFold
from sklearn.preprocessing import StandardScaler, LabelEncoder
from sklearn.feature_selection import SelectKBest, f_classif
from sklearn.impute import SimpleImputer
from sklearn.metrics import (
    classification_report, confusion_matrix, accuracy_score,
    f1_score, precision_score, recall_score
)
 
# Models
from sklearn.linear_model import LogisticRegression
from sklearn.ensemble import (
    RandomForestClassifier, ExtraTreesClassifier, BaggingClassifier
)
from sklearn.tree import DecisionTreeClassifier
from sklearn.decomposition import PCA
from scipy.stats import mode
 
import matplotlib.pyplot as plt
import seaborn as sns
 
# Reproducibility
SEED = 42
np.random.seed(SEED)
print('Setup complete ✓')
\end{lstlisting}
 
\begin{lstlisting}
# === Configuration ===
class Config:
    RANDOM_SEED = 42
    N_FOLDS = 5
    N_TOP_FEATURES = 15          # 15 πρωτεΐνες
 
    # Random Forest
    RF_N_ESTIMATORS = 300
    RF_MAX_DEPTH = None
 
    # Random Subspace
    RS_N_ESTIMATORS = 200
    RS_MAX_FEATURES = 0.5        # 50% features per tree
 
    # Rotation Forest
    ROT_N_ESTIMATORS = 100
    ROT_N_GROUPS = 3             # 15 features / 3 = 5 per group
 
    # Custom Cascade Forest
    DF_N_TREES = 100
    DF_MAX_LAYERS = 10
 
print('Config loaded ✓')
\end{lstlisting}
 
\begin{lstlisting}
# ============================================================
# ΑΛΛΑΞΕ ΑΥΤΑ ΤΑ PATHS ΣΤΑ ΔΙΚΑ ΣΟΥ ΑΡΧΕΙΑ
# ============================================================
SAMPLES_PATH = '/content/DIOPTRA_Samples_list_Oncology.xlsx'
PROTEINS_PATH = '/content/DIOPTRA_Protein_results_Oncology.xlsx'
# ============================================================
\end{lstlisting}
 
\begin{lstlisting}
# === Φόρτωση Δεδομένων (ίδιο pipeline με τα RNN notebooks) ===
 
def load_data(samples_path, proteins_path, include_t00=True):
    samples = pd.read_excel(samples_path)
    proteins = pd.read_excel(proteins_path)
    data = samples.merge(proteins, on='SampleID')
    protein_cols = [c for c in proteins.columns if c != 'SampleID']
 
    # Paired samples (T0 → T1)
    t0_data = data[data['Timepoint'] == 'T0'].copy()
    t1_data = data[data['Timepoint'] == 'T1'].copy()
    t0_data['PatientID'] = t0_data['SampleID']
    t1_data['PatientID'] = t1_data['SampleID'].str.replace('-F', '', regex=False)
 
    paired = t0_data.merge(
        t1_data[['PatientID'] + protein_cols],
        on='PatientID', suffixes=('_T0', '_T1')
    )
    print(f'Paired samples: {len(paired)}')
 
    if include_t00:
        t00_data = data[data['Timepoint'] == 'T00'].copy()
        for col in protein_cols:
            t00_data[f'{col}_T0'] = t00_data[col]
            t00_data[f'{col}_T1'] = t00_data[col]
        t00_data['PatientID'] = t00_data['SampleID']
        cols_needed = ['PatientID', 'Group'] + [f'{c}_T0' for c in protein_cols] + [f'{c}_T1' for c in protein_cols]
        combined = pd.concat([paired[cols_needed], t00_data[cols_needed]], ignore_index=True)
        print(f'After adding T00: {len(combined)}')
    else:
        combined = paired
 
    combined = combined[combined['Group'] != 'NAA']
    print(f'After removing NAA: {len(combined)}')
 
    X_T0 = combined[[f'{c}_T0' for c in protein_cols]].values
    X_T1 = combined[[f'{c}_T1' for c in protein_cols]].values
    X = np.stack([X_T0, X_T1], axis=1)
 
    label_encoder = LabelEncoder()
    y = label_encoder.fit_transform(combined['Group'].values)
 
    print(f'\nClass distribution:')
    for i, cls in enumerate(label_encoder.classes_):
        print(f'  {cls}: {(y == i).sum()}')
 
    return X, y, protein_cols, label_encoder
 
X_raw, y, protein_names, label_encoder = load_data(SAMPLES_PATH, PROTEINS_PATH)
class_names = list(label_encoder.classes_)
print(f'\nX shape: {X_raw.shape}  |  Classes: {class_names}')
\end{lstlisting}
 
\begin{lstlisting}
# === Προεπεξεργασία: NaN handling + Feature Selection (15 proteins) ===
 
# Βήμα 1: Αφαίρεση πρωτεϊνών με μόνο NaN
n_samples, n_timepoints, n_proteins = X_raw.shape
valid_mask = ~np.all(np.isnan(X_raw), axis=(0, 1))
X_raw = X_raw[:, :, valid_mask]
protein_names = [p for p, v in zip(protein_names, valid_mask) if v]
print(f'Valid proteins: {X_raw.shape[2]}')
 
# Βήμα 2: Median Imputation
imputer = SimpleImputer(strategy='median')
X_imp = np.zeros_like(X_raw)
for t in range(n_timepoints):
    X_imp[:, t, :] = imputer.fit_transform(X_raw[:, t, :])
X_raw = X_imp
print(f'NaN after imputation: {np.isnan(X_raw).sum()}')
 
# Βήμα 3: Feature Selection - Top 15 πρωτεΐνες (ANOVA F-score)
selector = SelectKBest(f_classif, k=Config.N_TOP_FEATURES)
selector.fit(X_raw[:, 0, :], y)  # Χρήση T0 για selection
selected_mask = selector.get_support()
selected_idx = np.where(selected_mask)[0]
 
X_selected = X_raw[:, :, selected_idx]
selected_proteins = [protein_names[i] for i in selected_idx]
 
print(f'\nSelected {len(selected_proteins)} proteins:')
f_scores = selector.scores_[selected_idx]
for rank, (p, f) in enumerate(sorted(zip(selected_proteins, f_scores), key=lambda x: -x[1]), 1):
    print(f'  {rank:2d}. {p:<12s} F={f:.2f}')
 
print(f'\nX_selected shape: {X_selected.shape}')
\end{lstlisting}
 
\begin{lstlisting}
# === Κατασκευή Feature Matrices για κάθε Mode ===
 
X_T0 = X_selected[:, 0, :]                           # (N, 15) - Static
X_T1 = X_selected[:, 1, :]                           # (N, 15)
X_delta = X_T1 - X_T0                                # (N, 15) - Delta
X_temporal = np.hstack([X_T0, X_T1])                 # (N, 30) - Temporal
X_full = np.hstack([X_T0, X_T1, X_delta])            # (N, 45) - Full
 
print('Feature matrices:')
print(f'  X_T0 (Static):      {X_T0.shape}')
print(f'  X_delta (Delta):    {X_delta.shape}')
print(f'  X_temporal (T0+T1): {X_temporal.shape}')
print(f'  X_full (T0+T1+Δ):  {X_full.shape}')
\end{lstlisting}
 
\subsubsection*{Βοηθητικές Συναρτήσεις \& Custom Models}
 
\begin{lstlisting}
# === Evaluation Framework ===
 
def evaluate_model(model_fn, X, y, n_folds=5, model_name='Model', normalize=True):
    """
    Αξιολόγηση μοντέλου με Stratified K-Fold CV.
    """
    skf = StratifiedKFold(n_splits=n_folds, shuffle=True, random_state=Config.RANDOM_SEED)
 
    fold_acc, fold_f1, fold_precision, fold_recall = [], [], [], []
    all_y_true, all_y_pred = [], []
    feature_importances = []
 
    for fold, (train_idx, val_idx) in enumerate(skf.split(X, y), 1):
        X_train, X_val = X[train_idx], X[val_idx]
        y_train, y_val = y[train_idx], y[val_idx]
 
        # Normalization ανά fold (αποφυγή data leakage)
        if normalize:
            scaler = StandardScaler()
            X_train = scaler.fit_transform(X_train)
            X_val = scaler.transform(X_val)
 
        # Train & Predict
        model = model_fn()
        model.fit(X_train, y_train)
        y_pred = model.predict(X_val)
 
        # Metrics
        acc = accuracy_score(y_val, y_pred)
        f1 = f1_score(y_val, y_pred, average='weighted', zero_division=0)
        prec = precision_score(y_val, y_pred, average='weighted', zero_division=0)
        rec = recall_score(y_val, y_pred, average='weighted', zero_division=0)
 
        fold_acc.append(acc)
        fold_f1.append(f1)
        fold_precision.append(prec)
        fold_recall.append(rec)
        all_y_true.extend(y_val)
        all_y_pred.extend(y_pred)
 
        if hasattr(model, 'feature_importances_'):
            feature_importances.append(model.feature_importances_)
 
        print(f'  Fold {fold}: Acc={acc:.3f}  F1={f1:.3f}')
 
    results = {
        'model': model_name,
        'mean_acc': np.mean(fold_acc),
        'std_acc': np.std(fold_acc),
        'mean_f1': np.mean(fold_f1),
        'std_f1': np.std(fold_f1),
        'mean_precision': np.mean(fold_precision),
        'mean_recall': np.mean(fold_recall),
        'fold_accs': fold_acc,
        'y_true': all_y_true,
        'y_pred': all_y_pred,
        'feature_importances': np.mean(feature_importances, axis=0) if feature_importances else None
    }
 
    print(f'  >>> {model_name}: Acc={results["mean_acc"]:.3f}±{results["std_acc"]:.3f}  '
          f'F1={results["mean_f1"]:.3f}±{np.std(fold_f1):.3f}')
    return results
 
print('Evaluation framework ready ✓')
\end{lstlisting}
 
\begin{lstlisting}
# === Rotation Forest Implementation ===
 
class RotationForestClassifier:
    """
    Rotation Forest (Rodriguez et al., 2006)
    Χωρίζει features σε ομάδες → PCA ανά ομάδα → Decision Tree
    """
    def __init__(self, n_estimators=100, n_groups=3, random_state=42):
        self.n_estimators = n_estimators
        self.n_groups = n_groups
        self.random_state = random_state
        self.trees = []
        self.rotations = []
        self.feature_groups = []
 
    def fit(self, X, y):
        rng = np.random.RandomState(self.random_state)
        n_features = X.shape[1]
        self.trees = []
        self.rotations = []
        self.feature_groups = []
 
        for i in range(self.n_estimators):
            # Τυχαίος χωρισμός features σε ομάδες
            perm = rng.permutation(n_features)
            groups = np.array_split(perm, self.n_groups)
 
            # PCA σε κάθε ομάδα (σε bootstrap subset)
            pca_list = []
            for group in groups:
                boot_idx = rng.choice(len(X), size=int(0.75 * len(X)), replace=True)
                X_subset = X[boot_idx][:, group]
                pca = PCA(n_components=len(group))
                pca.fit(X_subset)
                pca_list.append(pca)
 
            # Μετασχηματισμός & εκπαίδευση
            X_rotated = self._transform(X, groups, pca_list)
            tree = DecisionTreeClassifier(random_state=self.random_state + i)
            tree.fit(X_rotated, y)
 
            self.trees.append(tree)
            self.rotations.append(pca_list)
            self.feature_groups.append(groups)
        return self
 
    def _transform(self, X, groups, pca_list):
        parts = [pca.transform(X[:, group]) for group, pca in zip(groups, pca_list)]
        return np.hstack(parts)
 
    def predict(self, X):
        predictions = np.array([
            tree.predict(self._transform(X, groups, pcas))
            for tree, groups, pcas in zip(self.trees, self.feature_groups, self.rotations)
        ])
        return mode(predictions, axis=0, keepdims=False).mode
 
print('Rotation Forest ready ✓')
\end{lstlisting}
 
\begin{lstlisting}
# === Custom Cascade Forest Implementation ===
 
class CustomCascadeForest:
    """
    Custom Cascade Forest (Zhou & Feng, 2017) - sklearn only.
    Κάθε layer: [RF, RF, ExtraTrees, ExtraTrees]
    Probabilities (μέσω CV) concatenate με αρχικά features → επόμενο layer.
    Auto-stop όταν δεν βελτιώνεται.
    """
    def __init__(self, n_trees=100, max_layers=10, random_state=42):
        self.n_trees = n_trees
        self.max_layers = max_layers
        self.random_state = random_state
        self.layers = []
        self.n_classes_ = None
 
    def _create_forests(self, seed):
        return [
            RandomForestClassifier(n_estimators=self.n_trees, random_state=seed),
            RandomForestClassifier(n_estimators=self.n_trees, random_state=seed+1),
            ExtraTreesClassifier(n_estimators=self.n_trees, random_state=seed+2),
            ExtraTreesClassifier(n_estimators=self.n_trees, random_state=seed+3),
        ]
 
    def _get_cv_probas(self, X, y, seed):
        """Class probabilities μέσω 3-fold CV (αποφυγή leakage)"""
        n_forests = 4
        probas = np.zeros((len(X), self.n_classes_ * n_forests))
        kf = StratifiedKFold(n_splits=3, shuffle=True, random_state=self.random_state)
 
        for f_idx in range(n_forests):
            proba_col = np.zeros((len(X), self.n_classes_))
            for tr_idx, va_idx in kf.split(X, y):
                if f_idx < 2:
                    clf = RandomForestClassifier(n_estimators=self.n_trees, random_state=seed+f_idx)
                else:
                    clf = ExtraTreesClassifier(n_estimators=self.n_trees, random_state=seed+f_idx)
                clf.fit(X[tr_idx], y[tr_idx])
                proba_col[va_idx] = clf.predict_proba(X[va_idx])
            probas[:, f_idx*self.n_classes_:(f_idx+1)*self.n_classes_] = proba_col
 
        return probas
 
    def fit(self, X, y):
        self.layers = []
        self.n_classes_ = len(np.unique(y))
        self.X_original_dim = X.shape[1]
 
        X_current = X.copy()
        best_acc = 0
        patience = 0
 
        for layer_idx in range(self.max_layers):
            seed = self.random_state + layer_idx * 10
 
            # CV probabilities για augmentation
            probas = self._get_cv_probas(X_current, y, seed)
 
            # Layer accuracy (average of 4 forests)
            avg_proba = np.zeros((len(X), self.n_classes_))
            for f_idx in range(4):
                avg_proba += probas[:, f_idx*self.n_classes_:(f_idx+1)*self.n_classes_]
            avg_proba /= 4
            layer_acc = accuracy_score(y, np.argmax(avg_proba, axis=1))
            print(f'    Layer {layer_idx+1}: Acc={layer_acc:.3f}')
 
            # Train forests on full data (for prediction)
            forests = self._create_forests(seed)
            for forest in forests:
                forest.fit(X_current, y)
            self.layers.append(forests)
 
            # Early stopping
            if layer_acc > best_acc + 0.001:
                best_acc = layer_acc
                patience = 0
            else:
                patience += 1
                if patience >= 2:
                    print(f'    Early stop at layer {layer_idx+1}')
                    break
 
            # Augment: original features + probabilities
            X_current = np.hstack([X, probas])
 
        return self
 
    def predict(self, X):
        X_current = X.copy()
 
        for forests in self.layers:
            probas_list = []
            for forest in forests:
                # Fit ήταν σε X_current, predict πρέπει στο ίδιο dim
                n_expected = forest.n_features_in_
                X_input = X_current[:, :n_expected] if X_current.shape[1] > n_expected else X_current
                probas_list.append(forest.predict_proba(X_input))
 
            all_probas = np.hstack(probas_list)
            X_current = np.hstack([X, all_probas])
 
        # Final: average probabilities from last layer
        avg_proba = np.mean(probas_list, axis=0)
        return np.argmax(avg_proba, axis=1)
 
print('Custom Cascade Forest ready ✓')
\end{lstlisting}
 
\subsubsection*{Εκτέλεση Πειραμάτων}
 
\begin{lstlisting}
# Αποθήκευση αποτελεσμάτων
all_results = []
 
# Κύριο input: T0 + T1 concatenated
X_main = X_temporal
input_label = 'Temporal (T0+T1)'
print(f'Main input: {input_label}, shape={X_main.shape}')
\end{lstlisting}
 
\begin{lstlisting}
# === 1. Logistic Regression ===
print('='*60)
print('1. LOGISTIC REGRESSION')
print('='*60)
 
res = evaluate_model(
    model_fn=lambda: LogisticRegression(
        max_iter=1000,
        class_weight='balanced',
        random_state=Config.RANDOM_SEED,
        C=1.0,
        solver='lbfgs',
        multi_class='multinomial'
    ),
    X=X_main, y=y,
    model_name='Logistic Regression',
    normalize=True
)
all_results.append(res)
\end{lstlisting}
 
\begin{lstlisting}
# === 2. Random Forest ===
print('='*60)
print('2. RANDOM FOREST')
print('='*60)
 
res = evaluate_model(
    model_fn=lambda: RandomForestClassifier(
        n_estimators=Config.RF_N_ESTIMATORS,
        max_depth=Config.RF_MAX_DEPTH,
        class_weight='balanced',
        random_state=Config.RANDOM_SEED,
        n_jobs=-1
    ),
    X=X_main, y=y,
    model_name='Random Forest',
    normalize=False
)
all_results.append(res)
\end{lstlisting}
 
\begin{lstlisting}
# === 3. Random Subspace (Feature Bagging) ===
print('='*60)
print('3. RANDOM SUBSPACE (Feature Bagging)')
print('='*60)
 
res = evaluate_model(
    model_fn=lambda: BaggingClassifier(
        estimator=DecisionTreeClassifier(),
        n_estimators=Config.RS_N_ESTIMATORS,
        max_features=Config.RS_MAX_FEATURES,
        bootstrap=True,
        bootstrap_features=True,
        random_state=Config.RANDOM_SEED,
        n_jobs=-1
    ),
    X=X_main, y=y,
    model_name='Random Subspace',
    normalize=False
)
all_results.append(res)
\end{lstlisting}
 
\begin{lstlisting}
# === 4. Rotation Forest ===
print('='*60)
print('4. ROTATION FOREST')
print('='*60)
 
res = evaluate_model(
    model_fn=lambda: RotationForestClassifier(
        n_estimators=Config.ROT_N_ESTIMATORS,
        n_groups=Config.ROT_N_GROUPS,
        random_state=Config.RANDOM_SEED
    ),
    X=X_main, y=y,
    model_name='Rotation Forest',
    normalize=True
)
all_results.append(res)
\end{lstlisting}
 
\begin{lstlisting}
# === 5. Custom Cascade Forest (Deep Forest) ===
print('='*60)
print('5. CUSTOM CASCADE FOREST (Deep Forest)')
print('='*60)
 
res = evaluate_model(
    model_fn=lambda: CustomCascadeForest(
        n_trees=Config.DF_N_TREES,
        max_layers=Config.DF_MAX_LAYERS,
        random_state=Config.RANDOM_SEED
    ),
    X=X_main, y=y,
    model_name='Cascade Forest (Custom)',
    normalize=False
)
all_results.append(res)
\end{lstlisting}
 
\subsubsection*{Συγκριτικά Αποτελέσματα}
 
\begin{lstlisting}
# === Συγκριτικός Πίνακας ===
 
# Προσθήκη RNN/LSTM αποτελεσμάτων (από προηγούμενα πειράματα)
rnn_results = [
    {'model': 'Vanilla RNN (50 prot)', 'mean_acc': 0.513, 'std_acc': 0.116},
    {'model': 'LSTM (50 prot)', 'mean_acc': 0.336, 'std_acc': 0.081},
]
 
# Δημιουργία DataFrame
rows = []
for r in all_results:
    rows.append({
        'Model': r['model'],
        'Mean Acc': f"{r['mean_acc']:.3f}",
        'Std Acc': f"{r['std_acc']:.3f}",
        'Mean F1': f"{r['mean_f1']:.3f}",
        'Mean Prec': f"{r['mean_precision']:.3f}",
        'Mean Rec': f"{r['mean_recall']:.3f}",
    })
for r in rnn_results:
    rows.append({
        'Model': r['model'],
        'Mean Acc': f"{r['mean_acc']:.3f}",
        'Std Acc': f"{r['std_acc']:.3f}",
        'Mean F1': '-', 'Mean Prec': '-', 'Mean Rec': '-',
    })
 
df_results = pd.DataFrame(rows)
print('\n' + '='*70)
print('ΣΥΓΚΡΙΤΙΚΟΣ ΠΙΝΑΚΑΣ ΑΠΟΤΕΛΕΣΜΑΤΩΝ')
print(f'Input: {input_label} | Features: {X_main.shape[1]} | N={len(y)} | 5-Fold CV')
print('='*70)
print(df_results.to_string(index=False))
print('='*70)
\end{lstlisting}
 
\begin{lstlisting}
# === Bar Chart Σύγκρισης ===
 
fig, ax = plt.subplots(figsize=(12, 6))
 
models = [r['model'] for r in all_results]
accs = [r['mean_acc'] for r in all_results]
stds = [r['std_acc'] for r in all_results]
 
# Add RNN/LSTM
models += ['Vanilla RNN', 'LSTM']
accs += [0.513, 0.336]
stds += [0.116, 0.081]
 
# Colors: ensemble = teal, DL = coral
colors = ['#1D9E75'] * len(all_results) + ['#D85A30'] * 2
 
bars = ax.bar(range(len(models)), accs, yerr=stds, capsize=4,
              color=colors, edgecolor='white', linewidth=0.5, alpha=0.85)
 
ax.axhline(y=1/3, color='gray', linestyle='--', alpha=0.5, label='Random (33%)')
 
ax.set_xticks(range(len(models)))
ax.set_xticklabels(models, rotation=35, ha='right', fontsize=10)
ax.set_ylabel('Mean Accuracy (5-Fold CV)', fontsize=12)
ax.set_title('Σύγκριση Μεθόδων: Ensemble vs Deep Learning\n'
             f'(15 πρωτεΐνες, N={len(y)}, 3 κλάσεις)', fontsize=13)
ax.set_ylim(0, 1.0)
ax.legend()
 
for bar, acc in zip(bars, accs):
    ax.text(bar.get_x() + bar.get_width()/2, bar.get_height() + 0.03,
            f'{acc:.1%}', ha='center', va='bottom', fontsize=9, fontweight='bold')
 
plt.tight_layout()
plt.savefig('model_comparison.png', dpi=150, bbox_inches='tight')
plt.show()
print('Saved: model_comparison.png')
\end{lstlisting}
 
\begin{lstlisting}
# === Confusion Matrices ===
 
n_models = len(all_results)
n_cols = 3
n_rows = (n_models + n_cols - 1) // n_cols
fig, axes = plt.subplots(n_rows, n_cols, figsize=(16, 5 * n_rows))
axes = axes.flatten()
 
for idx, r in enumerate(all_results):
    cm = confusion_matrix(r['y_true'], r['y_pred'])
    sns.heatmap(cm, annot=True, fmt='d', cmap='Blues',
                xticklabels=class_names, yticklabels=class_names,
                ax=axes[idx])
    axes[idx].set_title(f"{r['model']}\nAcc={r['mean_acc']:.3f}", fontsize=11)
    axes[idx].set_xlabel('Predicted')
    axes[idx].set_ylabel('True')
 
for idx in range(n_models, len(axes)):
    axes[idx].set_visible(False)
 
plt.suptitle('Confusion Matrices - Ensemble Methods (5-Fold CV, Aggregated)', fontsize=14, y=1.02)
plt.tight_layout()
plt.savefig('confusion_matrices.png', dpi=150, bbox_inches='tight')
plt.show()
\end{lstlisting}
 
\begin{lstlisting}
# === Feature Importance (Random Forest) ===
 
rf_result = [r for r in all_results if r['model'] == 'Random Forest'][0]
 
if rf_result['feature_importances'] is not None:
    feat_names = ([f'{p}_T0' for p in selected_proteins] +
                  [f'{p}_T1' for p in selected_proteins])
 
    importances = rf_result['feature_importances']
    sorted_idx = np.argsort(importances)[::-1][:20]
 
    fig, ax = plt.subplots(figsize=(10, 8))
    ax.barh(range(len(sorted_idx)), importances[sorted_idx][::-1],
            color='#1D9E75', alpha=0.8)
    ax.set_yticks(range(len(sorted_idx)))
    ax.set_yticklabels([feat_names[i] for i in sorted_idx][::-1], fontsize=10)
    ax.set_xlabel('Feature Importance (MDI)', fontsize=12)
    ax.set_title('Top 20 Feature Importances (Random Forest)', fontsize=13)
    plt.tight_layout()
    plt.savefig('feature_importance_rf.png', dpi=150, bbox_inches='tight')
    plt.show()
    print('Saved: feature_importance_rf.png')
\end{lstlisting}
 
\begin{lstlisting}
# === Αποθήκευση σε CSV ===
 
detailed_rows = []
for r in all_results:
    row = {
        'Model': r['model'],
        'Mean_Accuracy': r['mean_acc'],
        'Std_Accuracy': r['std_acc'],
        'Mean_F1': r['mean_f1'],
        'Mean_Precision': r['mean_precision'],
        'Mean_Recall': r['mean_recall'],
    }
    for i, acc in enumerate(r['fold_accs'], 1):
        row[f'Fold_{i}'] = acc
    detailed_rows.append(row)
 
df_detailed = pd.DataFrame(detailed_rows)
df_detailed.to_csv('ensemble_results.csv', index=False)
print('Saved: ensemble_results.csv')
print(df_detailed.to_string(index=False))
\end{lstlisting}
 
\begin{lstlisting}
# === BONUS: Static vs Temporal σύγκριση ===
print('\n' + '='*70)
print('BONUS: Static (T0) vs Temporal (T0+T1) σύγκριση')
print('='*70)
 
comparison_rows = []
 
models_to_compare = [
    ('Logistic Regression', lambda: LogisticRegression(
        max_iter=1000, class_weight='balanced', random_state=42,
        solver='lbfgs', multi_class='multinomial'), True),
    ('Random Forest', lambda: RandomForestClassifier(
        n_estimators=300, class_weight='balanced', random_state=42, n_jobs=-1), False),
    ('Cascade Forest', lambda: CustomCascadeForest(
        n_trees=100, max_layers=10, random_state=42), False),
]
 
for name, model_fn, needs_norm in models_to_compare:
    print(f'\n--- {name} ---')
 
    print('  [Static T0]')
    r_static = evaluate_model(model_fn, X_T0, y,
        model_name=f'{name} (Static)', normalize=needs_norm)
 
    print('  [Temporal T0+T1]')
    r_temporal = evaluate_model(model_fn, X_temporal, y,
        model_name=f'{name} (Temporal)', normalize=needs_norm)
 
    comparison_rows.append({
        'Model': name,
        'Static Acc': f"{r_static['mean_acc']:.3f}±{r_static['std_acc']:.3f}",
        'Temporal Acc': f"{r_temporal['mean_acc']:.3f}±{r_temporal['std_acc']:.3f}",
        'Δ Acc': f"{r_temporal['mean_acc'] - r_static['mean_acc']:+.3f}",
    })
 
print('\n' + pd.DataFrame(comparison_rows).to_string(index=False))
\end{lstlisting}
 
\subsubsection*{Συμπεράσματα}
 

\subsection{A.4 Βελτιστοποίηση \& Hyperparameter Tuning}
 
\noindent \textbf{Πηγή:} \texttt{Optimization\_Experiments.ipynb}
 
 
\subsubsection*{Optimization Experiments v2}
 
\begin{lstlisting}
!pip install xgboost lightgbm imbalanced-learn -q
print('Installations complete ✓')
\end{lstlisting}
 
\begin{lstlisting}
import numpy as np
import pandas as pd
import warnings
warnings.filterwarnings('ignore')
 
from sklearn.model_selection import StratifiedKFold, GridSearchCV
from sklearn.preprocessing import StandardScaler, LabelEncoder
from sklearn.feature_selection import SelectKBest, f_classif
from sklearn.impute import SimpleImputer
from sklearn.metrics import accuracy_score, f1_score, precision_score, recall_score, confusion_matrix
from sklearn.linear_model import LogisticRegression
from sklearn.ensemble import RandomForestClassifier, ExtraTreesClassifier, BaggingClassifier, StackingClassifier, VotingClassifier
from sklearn.tree import DecisionTreeClassifier
from sklearn.svm import SVC
from sklearn.decomposition import PCA
from xgboost import XGBClassifier
from lightgbm import LGBMClassifier
from imblearn.over_sampling import SMOTE
from scipy.stats import mode
import matplotlib.pyplot as plt
import seaborn as sns
 
SEED = 42
np.random.seed(SEED)
print('Setup complete ✓')
\end{lstlisting}
 
\begin{lstlisting}
SAMPLES_PATH = '/content/DIOPTRA_Samples_list_Oncology.xlsx'
PROTEINS_PATH = '/content/DIOPTRA_Protein_results_Oncology.xlsx'
\end{lstlisting}
 
\begin{lstlisting}
def load_data(samples_path, proteins_path, include_t00=True):
    samples = pd.read_excel(samples_path)
    proteins = pd.read_excel(proteins_path)
    data = samples.merge(proteins, on='SampleID')
    protein_cols = [c for c in proteins.columns if c != 'SampleID']
    t0_data = data[data['Timepoint'] == 'T0'].copy()
    t1_data = data[data['Timepoint'] == 'T1'].copy()
    t0_data['PatientID'] = t0_data['SampleID']
    t1_data['PatientID'] = t1_data['SampleID'].str.replace('-F', '', regex=False)
    paired = t0_data.merge(t1_data[['PatientID'] + protein_cols], on='PatientID', suffixes=('_T0', '_T1'))
    print(f'Paired samples: {len(paired)}')
    if include_t00:
        t00_data = data[data['Timepoint'] == 'T00'].copy()
        for col in protein_cols:
            t00_data[f'{col}_T0'] = t00_data[col]
            t00_data[f'{col}_T1'] = t00_data[col]
        t00_data['PatientID'] = t00_data['SampleID']
        cols_needed = ['PatientID', 'Group'] + [f'{c}_T0' for c in protein_cols] + [f'{c}_T1' for c in protein_cols]
        combined = pd.concat([paired[cols_needed], t00_data[cols_needed]], ignore_index=True)
        print(f'After adding T00: {len(combined)}')
    else:
        combined = paired
    combined = combined[combined['Group'] != 'NAA']
    print(f'After removing NAA: {len(combined)}')
    X_T0 = combined[[f'{c}_T0' for c in protein_cols]].values
    X_T1 = combined[[f'{c}_T1' for c in protein_cols]].values
    X = np.stack([X_T0, X_T1], axis=1)
    label_encoder = LabelEncoder()
    y = label_encoder.fit_transform(combined['Group'].values)
    print(f'\nClass distribution:')
    for i, cls in enumerate(label_encoder.classes_):
        print(f'  {cls}: {(y == i).sum()}')
    return X, y, protein_cols, label_encoder
 
X_raw, y, protein_names, label_encoder = load_data(SAMPLES_PATH, PROTEINS_PATH)
class_names = list(label_encoder.classes_)
print(f'\nX shape: {X_raw.shape}  |  Classes: {class_names}')
\end{lstlisting}
 
\begin{lstlisting}
# Preprocessing
n_samples, n_timepoints, n_proteins = X_raw.shape
valid_mask = ~np.all(np.isnan(X_raw), axis=(0, 1))
X_raw = X_raw[:, :, valid_mask]
protein_names = [p for p, v in zip(protein_names, valid_mask) if v]
imputer = SimpleImputer(strategy='median')
X_imp = np.zeros_like(X_raw)
for t in range(n_timepoints):
    X_imp[:, t, :] = imputer.fit_transform(X_raw[:, t, :])
X_raw = X_imp
print(f'Preprocessed: {X_raw.shape}')
\end{lstlisting}
 
\begin{lstlisting}
# Feature Selection: Top 15
K = 15
selector = SelectKBest(f_classif, k=K)
selector.fit(X_raw[:, 0, :], y)
selected_idx = np.where(selector.get_support())[0]
X_selected = X_raw[:, :, selected_idx]
selected_proteins = [protein_names[i] for i in selected_idx]
X_T0 = X_selected[:, 0, :]
X_T1 = X_selected[:, 1, :]
print(f'Selected {K} proteins')
\end{lstlisting}
 
\begin{lstlisting}
# Feature Engineering
X_delta = X_T1 - X_T0
X_temporal = np.hstack([X_T0, X_T1])
X_ratio = np.divide(X_T1, X_T0, out=np.ones_like(X_T0), where=X_T0 != 0)
X_mean = (X_T0 + X_T1) / 2
X_abs_delta = np.abs(X_delta)
 
feature_sets = {
    'Temporal (T0+T1)':      X_temporal,
    'Full (T0+T1+D)':       np.hstack([X_T0, X_T1, X_delta]),
    'Enriched (T0+T1+D+R)': np.hstack([X_T0, X_T1, X_delta, X_ratio]),
    'Kitchen Sink':          np.hstack([X_T0, X_T1, X_delta, X_ratio, X_mean, X_abs_delta]),
    'Compact (T0+T1+R)':    np.hstack([X_T0, X_T1, X_ratio]),
    'Ratios Only':           X_ratio,
    'Delta + Ratio':         np.hstack([X_delta, X_ratio]),
}
 
print('Feature sets:')
for name, X in feature_sets.items():
    print(f'  {name:25s} -> {X.shape[1]} features')
\end{lstlisting}
 
\begin{lstlisting}
# Evaluation function - FIXED: all params explicit
def evaluate_model(model_fn, X, y, model_name='Model', n_folds=5,
                   normalize=True, use_smote=False):
    skf = StratifiedKFold(n_splits=n_folds, shuffle=True, random_state=SEED)
    fold_acc, fold_f1 = [], []
    all_y_true, all_y_pred = [], []
    feature_importances = []
 
    for fold, (train_idx, val_idx) in enumerate(skf.split(X, y), 1):
        X_train, X_val = X[train_idx], X[val_idx]
        y_train, y_val = y[train_idx], y[val_idx]
        if normalize:
            scaler = StandardScaler()
            X_train = scaler.fit_transform(X_train)
            X_val = scaler.transform(X_val)
        if use_smote:
            try:
                min_class = min(np.bincount(y_train))
                k = min(5, min_class - 1)
                if k >= 1:
                    sm = SMOTE(random_state=SEED, k_neighbors=k)
                    X_train, y_train = sm.fit_resample(X_train, y_train)
            except Exception:
                pass
        model = model_fn()
        model.fit(X_train, y_train)
        y_pred = model.predict(X_val)
        acc = accuracy_score(y_val, y_pred)
        f1 = f1_score(y_val, y_pred, average='weighted', zero_division=0)
        fold_acc.append(acc)
        fold_f1.append(f1)
        all_y_true.extend(y_val)
        all_y_pred.extend(y_pred)
        if hasattr(model, 'feature_importances_'):
            feature_importances.append(model.feature_importances_)
 
    results = {
        'model': model_name,
        'mean_acc': np.mean(fold_acc),
        'std_acc': np.std(fold_acc),
        'mean_f1': np.mean(fold_f1),
        'std_f1': np.std(fold_f1),
        'fold_accs': fold_acc,
        'y_true': all_y_true,
        'y_pred': all_y_pred,
    }
    print(f'  {model_name:45s} Acc={results["mean_acc"]:.3f}+/-{results["std_acc"]:.3f}  '
          f'F1={results["mean_f1"]:.3f}')
    return results
 
print('Evaluation ready ✓')
\end{lstlisting}
 
\begin{lstlisting}
# === ΠΕΙΡΑΜΑ 1: Feature Set Comparison ===
print('='*70)
print('ΠΕΙΡΑΜΑ 1: Feature Set Comparison (Logistic Regression)')
print('='*70)
 
feat_results = []
for name, X_feat in feature_sets.items():
    r = evaluate_model(
        lambda: LogisticRegression(max_iter=2000, class_weight='balanced',
                                   random_state=SEED, solver='lbfgs',
                                   multi_class='multinomial', C=1.0),
        X_feat, y, model_name=name, normalize=True
    )
    feat_results.append(r)
 
best_feat = max(feat_results, key=lambda x: x['mean_acc'])
print(f'\n>>> Best feature set: {best_feat["model"]} ({best_feat["mean_acc"]:.3f})')
\end{lstlisting}
 
\begin{lstlisting}
best_feat_name = best_feat['model']
X_best = feature_sets[best_feat_name]
print(f'Using: {best_feat_name} ({X_best.shape[1]} features)')
\end{lstlisting}
 
\begin{lstlisting}
# === ΠΕΙΡΑΜΑ 2: New Models ===
print('='*70)
print('ΠΕΙΡΑΜΑ 2: New Models')
print('='*70)
 
all_results = []
 
r = evaluate_model(
    lambda: LogisticRegression(max_iter=2000, class_weight='balanced',
                               random_state=SEED, solver='lbfgs',
                               multi_class='multinomial'),
    X_best, y, model_name='LR (baseline)', normalize=True)
all_results.append(r)
 
r = evaluate_model(
    lambda: SVC(kernel='rbf', class_weight='balanced',
                random_state=SEED, C=1.0, gamma='scale'),
    X_best, y, model_name='SVM (RBF)', normalize=True)
all_results.append(r)
 
r = evaluate_model(
    lambda: SVC(kernel='linear', class_weight='balanced',
                random_state=SEED, C=1.0),
    X_best, y, model_name='SVM (Linear)', normalize=True)
all_results.append(r)
 
r = evaluate_model(
    lambda: XGBClassifier(n_estimators=200, max_depth=3, learning_rate=0.1,
                          random_state=SEED, eval_metric='mlogloss',
                          use_label_encoder=False),
    X_best, y, model_name='XGBoost', normalize=False)
all_results.append(r)
 
r = evaluate_model(
    lambda: LGBMClassifier(n_estimators=200, max_depth=3, learning_rate=0.1,
                           random_state=SEED, verbose=-1,
                           class_weight='balanced'),
    X_best, y, model_name='LightGBM', normalize=False)
all_results.append(r)
 
r = evaluate_model(
    lambda: RandomForestClassifier(n_estimators=300, class_weight='balanced',
                                   random_state=SEED, n_jobs=-1),
    X_best, y, model_name='Random Forest', normalize=False)
all_results.append(r)
 
r = evaluate_model(
    lambda: ExtraTreesClassifier(n_estimators=300, class_weight='balanced',
                                  random_state=SEED, n_jobs=-1),
    X_best, y, model_name='Extra Trees', normalize=False)
all_results.append(r)
\end{lstlisting}
 
\begin{lstlisting}
# === ΠΕΙΡΑΜΑ 3: Hyperparameter Tuning ===
print('='*70)
print('ΠΕΙΡΑΜΑ 3: Hyperparameter Tuning (GridSearchCV)')
print('='*70)
 
scaler_full = StandardScaler()
X_best_scaled = scaler_full.fit_transform(X_best)
cv_inner = StratifiedKFold(n_splits=5, shuffle=True, random_state=SEED)
 
print('\n--- Logistic Regression ---')
gs_lr = GridSearchCV(
    LogisticRegression(max_iter=2000, class_weight='balanced',
                       random_state=SEED, multi_class='multinomial', solver='lbfgs'),
    {'C': [0.001, 0.01, 0.1, 0.5, 1.0, 5.0, 10.0, 50.0]},
    cv=cv_inner, scoring='accuracy', n_jobs=-1
)
gs_lr.fit(X_best_scaled, y)
print(f'  Best C={gs_lr.best_params_["C"]}, Acc={gs_lr.best_score_:.3f}')
 
print('\n--- SVM (RBF) ---')
gs_svm = GridSearchCV(
    SVC(kernel='rbf', class_weight='balanced', random_state=SEED),
    {'C': [0.1, 0.5, 1.0, 5.0, 10.0, 50.0],
     'gamma': ['scale', 'auto', 0.01, 0.1, 0.5]},
    cv=cv_inner, scoring='accuracy', n_jobs=-1
)
gs_svm.fit(X_best_scaled, y)
print(f'  Best C={gs_svm.best_params_["C"]}, gamma={gs_svm.best_params_["gamma"]}, Acc={gs_svm.best_score_:.3f}')
 
print('\n--- XGBoost ---')
gs_xgb = GridSearchCV(
    XGBClassifier(random_state=SEED, eval_metric='mlogloss', use_label_encoder=False),
    {'n_estimators': [50, 100, 200, 300], 'max_depth': [2, 3, 4, 5],
     'learning_rate': [0.01, 0.05, 0.1, 0.2]},
    cv=cv_inner, scoring='accuracy', n_jobs=-1
)
gs_xgb.fit(X_best, y)
print(f'  Best: {gs_xgb.best_params_}, Acc={gs_xgb.best_score_:.3f}')
 
print('\n--- LightGBM ---')
gs_lgbm = GridSearchCV(
    LGBMClassifier(random_state=SEED, verbose=-1, class_weight='balanced'),
    {'n_estimators': [50, 100, 200, 300], 'max_depth': [2, 3, 4, 5, -1],
     'learning_rate': [0.01, 0.05, 0.1, 0.2], 'num_leaves': [7, 15, 31]},
    cv=cv_inner, scoring='accuracy', n_jobs=-1
)
gs_lgbm.fit(X_best, y)
print(f'  Best: {gs_lgbm.best_params_}, Acc={gs_lgbm.best_score_:.3f}')
 
print('\n--- Summary ---')
print(f'  LR:      {gs_lr.best_score_:.3f}')
print(f'  SVM:     {gs_svm.best_score_:.3f}')
print(f'  XGBoost: {gs_xgb.best_score_:.3f}')
print(f'  LightGBM:{gs_lgbm.best_score_:.3f}')
\end{lstlisting}
 
\begin{lstlisting}
# === Tuned Models Evaluation ===
print('='*70)
print('TUNED MODELS (proper 5-Fold CV)')
print('='*70)
 
tuned_results = []
 
best_C_lr = gs_lr.best_params_['C']
r = evaluate_model(
    lambda: LogisticRegression(max_iter=2000, class_weight='balanced',
                               random_state=SEED, solver='lbfgs',
                               multi_class='multinomial', C=best_C_lr),
    X_best, y, model_name=f'LR tuned (C={best_C_lr})', normalize=True)
tuned_results.append(r)
 
best_C_svm = gs_svm.best_params_['C']
best_gamma = gs_svm.best_params_['gamma']
r = evaluate_model(
    lambda: SVC(kernel='rbf', class_weight='balanced', random_state=SEED,
                C=best_C_svm, gamma=best_gamma),
    X_best, y, model_name=f'SVM tuned (C={best_C_svm})', normalize=True)
tuned_results.append(r)
 
xgb_params = gs_xgb.best_params_
r = evaluate_model(
    lambda: XGBClassifier(**xgb_params, random_state=SEED,
                          eval_metric='mlogloss', use_label_encoder=False),
    X_best, y, model_name='XGBoost tuned', normalize=False)
tuned_results.append(r)
 
lgbm_params = gs_lgbm.best_params_
r = evaluate_model(
    lambda: LGBMClassifier(**lgbm_params, random_state=SEED, verbose=-1,
                           class_weight='balanced'),
    X_best, y, model_name='LightGBM tuned', normalize=False)
tuned_results.append(r)
\end{lstlisting}
 
\begin{lstlisting}
# === ΠΕΙΡΑΜΑ 4: Stacking & Voting ===
print('='*70)
print('ΠΕΙΡΑΜΑ 4: Stacking & Voting')
print('='*70)
 
r = evaluate_model(
    lambda: VotingClassifier(
        estimators=[
            ('lr', LogisticRegression(max_iter=2000, class_weight='balanced',
                                      random_state=SEED, C=best_C_lr,
                                      solver='lbfgs', multi_class='multinomial')),
            ('svm', SVC(kernel='rbf', class_weight='balanced', random_state=SEED,
                        C=best_C_svm, gamma=best_gamma, probability=True)),
            ('rf', RandomForestClassifier(n_estimators=300, class_weight='balanced',
                                          random_state=SEED)),
            ('xgb', XGBClassifier(**xgb_params, random_state=SEED,
                                  eval_metric='mlogloss', use_label_encoder=False)),
        ],
        voting='soft'
    ),
    X_best, y, model_name='Voting (LR+SVM+RF+XGB)', normalize=True)
tuned_results.append(r)
 
r = evaluate_model(
    lambda: StackingClassifier(
        estimators=[
            ('lr', LogisticRegression(max_iter=2000, class_weight='balanced',
                                      random_state=SEED, C=best_C_lr,
                                      solver='lbfgs', multi_class='multinomial')),
            ('svm', SVC(kernel='rbf', class_weight='balanced', random_state=SEED,
                        C=best_C_svm, gamma=best_gamma, probability=True)),
            ('rf', RandomForestClassifier(n_estimators=300, class_weight='balanced',
                                          random_state=SEED)),
            ('xgb', XGBClassifier(**xgb_params, random_state=SEED,
                                  eval_metric='mlogloss', use_label_encoder=False)),
        ],
        final_estimator=LogisticRegression(max_iter=2000, random_state=SEED),
        cv=3
    ),
    X_best, y, model_name='Stacking (LR+SVM+RF+XGB -> LR)', normalize=True)
tuned_results.append(r)
 
r = evaluate_model(
    lambda: StackingClassifier(
        estimators=[
            ('lr', LogisticRegression(max_iter=2000, class_weight='balanced',
                                      random_state=SEED, C=best_C_lr,
                                      solver='lbfgs', multi_class='multinomial')),
            ('svm', SVC(kernel='rbf', class_weight='balanced', random_state=SEED,
                        C=best_C_svm, gamma=best_gamma, probability=True)),
            ('rf', RandomForestClassifier(n_estimators=300, class_weight='balanced',
                                          random_state=SEED)),
        ],
        final_estimator=RandomForestClassifier(n_estimators=100, random_state=SEED),
        cv=3
    ),
    X_best, y, model_name='Stacking (LR+SVM+RF -> RF)', normalize=True)
tuned_results.append(r)
\end{lstlisting}
 
\begin{lstlisting}
# === ΠΕΙΡΑΜΑ 5: SMOTE ===
print('='*70)
print('ΠΕΙΡΑΜΑ 5: SMOTE')
print('='*70)
 
r = evaluate_model(
    lambda: LogisticRegression(max_iter=2000, class_weight='balanced',
                               random_state=SEED, C=best_C_lr,
                               solver='lbfgs', multi_class='multinomial'),
    X_best, y, model_name='LR tuned + SMOTE', normalize=True, use_smote=True)
tuned_results.append(r)
 
r = evaluate_model(
    lambda: SVC(kernel='rbf', class_weight='balanced', random_state=SEED,
                C=best_C_svm, gamma=best_gamma),
    X_best, y, model_name='SVM tuned + SMOTE', normalize=True, use_smote=True)
tuned_results.append(r)
 
r = evaluate_model(
    lambda: XGBClassifier(**xgb_params, random_state=SEED,
                          eval_metric='mlogloss', use_label_encoder=False),
    X_best, y, model_name='XGBoost tuned + SMOTE', normalize=False, use_smote=True)
tuned_results.append(r)
\end{lstlisting}
 
\begin{lstlisting}
# === ΠΕΙΡΑΜΑ 6: Feature Count Sweep ===
print('='*70)
print('ΠΕΙΡΑΜΑ 6: Feature Count Sweep')
print('='*70)
 
k_values = [5, 8, 10, 12, 15, 20, 25, 30, 40, 50]
k_results = []
 
for k in k_values:
    k_actual = min(k, X_raw.shape[2])
    sel = SelectKBest(f_classif, k=k_actual)
    sel.fit(X_raw[:, 0, :], y)
    idx = np.where(sel.get_support())[0]
    X_k = X_raw[:, :, idx]
    X_k_temp = np.hstack([X_k[:, 0, :], X_k[:, 1, :]])
 
    r = evaluate_model(
        lambda: LogisticRegression(max_iter=2000, class_weight='balanced',
                                   random_state=SEED, C=best_C_lr,
                                   solver='lbfgs', multi_class='multinomial'),
        X_k_temp, y, model_name=f'LR (K={k_actual})', normalize=True)
    k_results.append((k_actual, r['mean_acc'], r['std_acc']))
 
fig, ax = plt.subplots(figsize=(10, 5))
ks, accs, stds = zip(*k_results)
ax.errorbar(ks, accs, yerr=stds, marker='o', capsize=4, color='#1D9E75', linewidth=2)
ax.axhline(y=1/3, color='gray', linestyle='--', alpha=0.5, label='Random (33%)')
ax.set_xlabel('Number of Proteins (K)', fontsize=12)
ax.set_ylabel('Mean Accuracy', fontsize=12)
ax.set_title('Accuracy vs Number of Proteins (LR, Temporal)', fontsize=13)
ax.legend()
ax.set_xticks(ks)
plt.tight_layout()
plt.savefig('feature_count_sweep.png', dpi=150, bbox_inches='tight')
plt.show()
best_k = max(k_results, key=lambda x: x[1])
print(f'\n>>> Best K={best_k[0]}: Acc={best_k[1]:.3f}+/-{best_k[2]:.3f}')
\end{lstlisting}
 
\begin{lstlisting}
# === MASTER TABLE ===
print('\n' + '='*70)
print('ΤΕΛΙΚΗ ΣΥΓΚΡΙΣΗ')
print('='*70)
 
all_combined = all_results + tuned_results
all_combined.sort(key=lambda x: x['mean_acc'], reverse=True)
 
rows = []
for r in all_combined:
    rows.append({'Model': r['model'], 'Acc': f"{r['mean_acc']:.3f}",
                 'Std': f"{r['std_acc']:.3f}", 'F1': f"{r['mean_f1']:.3f}"})
rows.append({'Model': '--- Deep Learning ---', 'Acc': '---', 'Std': '---', 'F1': '---'})
rows.append({'Model': 'Vanilla RNN (50 prot)', 'Acc': '0.513', 'Std': '0.116', 'F1': '-'})
rows.append({'Model': 'LSTM (50 prot)', 'Acc': '0.336', 'Std': '0.081', 'F1': '-'})
 
print(pd.DataFrame(rows).to_string(index=False))
\end{lstlisting}
 
\begin{lstlisting}
# === Bar Chart Top 10 ===
top_n = min(10, len(all_combined))
top_models = all_combined[:top_n]
 
fig, ax = plt.subplots(figsize=(14, 6))
models = [r['model'] for r in top_models] + ['Vanilla RNN', 'LSTM']
accs = [r['mean_acc'] for r in top_models] + [0.513, 0.336]
stds = [r['std_acc'] for r in top_models] + [0.116, 0.081]
colors = ['#1D9E75'] * top_n + ['#D85A30'] * 2
 
bars = ax.bar(range(len(models)), accs, yerr=stds, capsize=4,
              color=colors, edgecolor='white', linewidth=0.5, alpha=0.85)
ax.axhline(y=1/3, color='gray', linestyle='--', alpha=0.5, label='Random (33%)')
ax.set_xticks(range(len(models)))
ax.set_xticklabels(models, rotation=40, ha='right', fontsize=9)
ax.set_ylabel('Mean Accuracy (5-Fold CV)')
ax.set_title('Top Models vs Deep Learning')
ax.set_ylim(0, 1.0)
ax.legend()
for bar, acc in zip(bars, accs):
    ax.text(bar.get_x() + bar.get_width()/2, bar.get_height() + 0.02,
            f'{acc:.1%}', ha='center', va='bottom', fontsize=8, fontweight='bold')
plt.tight_layout()
plt.savefig('optimization_comparison.png', dpi=150, bbox_inches='tight')
plt.show()
\end{lstlisting}
 
\begin{lstlisting}
# === Save Results ===
save_rows = []
for r in all_combined:
    row = {'Model': r['model'], 'Mean_Accuracy': r['mean_acc'],
           'Std_Accuracy': r['std_acc'], 'Mean_F1': r['mean_f1']}
    for i, acc in enumerate(r['fold_accs'], 1):
        row[f'Fold_{i}'] = acc
    save_rows.append(row)
df_save = pd.DataFrame(save_rows)
df_save.to_csv('optimization_results.csv', index=False)
print('Saved: optimization_results.csv')
print(f'Best: {all_combined[0]["model"]} = {all_combined[0]["mean_acc"]:.3f}')
\end{lstlisting}
 
\begin{lstlisting}
# === Best Model Confusion Matrix ===
best = all_combined[0]
cm = confusion_matrix(best['y_true'], best['y_pred'])
fig, ax = plt.subplots(figsize=(7, 6))
sns.heatmap(cm, annot=True, fmt='d', cmap='Blues',
            xticklabels=class_names, yticklabels=class_names, ax=ax)
ax.set_title(f"Best: {best['model']}\nAcc={best['mean_acc']:.3f}")
ax.set_xlabel('Predicted')
ax.set_ylabel('True')
plt.tight_layout()
plt.savefig('best_model_cm.png', dpi=150, bbox_inches='tight')
plt.show()
\end{lstlisting}
 

\subsection{A.5 Αυστηρή Στατιστική Αξιολόγηση}
 
\noindent \textbf{Πηγή:} \texttt{Final\_Rigorous\_Evaluation.ipynb}
 
 
\subsubsection*{Final Rigorous Evaluation}
 
\begin{lstlisting}
!pip install xgboost imbalanced-learn -q
print('Installations complete ✓')
\end{lstlisting}
 
\begin{lstlisting}
import numpy as np
import pandas as pd
import warnings
warnings.filterwarnings('ignore')
 
from sklearn.model_selection import (
    StratifiedKFold, RepeatedStratifiedKFold, GridSearchCV,
    cross_val_score, permutation_test_score, learning_curve
)
from sklearn.preprocessing import StandardScaler, LabelEncoder
from sklearn.feature_selection import SelectKBest, f_classif
from sklearn.impute import SimpleImputer
from sklearn.metrics import (
    accuracy_score, f1_score, confusion_matrix, brier_score_loss,
    classification_report, roc_auc_score
)
from sklearn.linear_model import LogisticRegression
from sklearn.ensemble import RandomForestClassifier
from sklearn.svm import SVC
from sklearn.pipeline import Pipeline
from sklearn.calibration import calibration_curve, CalibratedClassifierCV
 
import matplotlib.pyplot as plt
import seaborn as sns
from scipy import stats
 
SEED = 42
np.random.seed(SEED)
print('Setup complete ✓')
\end{lstlisting}
 
\begin{lstlisting}
SAMPLES_PATH = '/content/DIOPTRA_Samples_list_Oncology.xlsx'
PROTEINS_PATH = '/content/DIOPTRA_Protein_results_Oncology.xlsx'
\end{lstlisting}
 
\begin{lstlisting}
def load_data(samples_path, proteins_path, include_t00=True):
    samples = pd.read_excel(samples_path)
    proteins = pd.read_excel(proteins_path)
    data = samples.merge(proteins, on='SampleID')
    protein_cols = [c for c in proteins.columns if c != 'SampleID']
    t0_data = data[data['Timepoint'] == 'T0'].copy()
    t1_data = data[data['Timepoint'] == 'T1'].copy()
    t0_data['PatientID'] = t0_data['SampleID']
    t1_data['PatientID'] = t1_data['SampleID'].str.replace('-F', '', regex=False)
    paired = t0_data.merge(t1_data[['PatientID'] + protein_cols], on='PatientID', suffixes=('_T0', '_T1'))
    print(f'Paired: {len(paired)}')
    if include_t00:
        t00_data = data[data['Timepoint'] == 'T00'].copy()
        for col in protein_cols:
            t00_data[f'{col}_T0'] = t00_data[col]
            t00_data[f'{col}_T1'] = t00_data[col]
        t00_data['PatientID'] = t00_data['SampleID']
        cols_needed = ['PatientID', 'Group'] + [f'{c}_T0' for c in protein_cols] + [f'{c}_T1' for c in protein_cols]
        combined = pd.concat([paired[cols_needed], t00_data[cols_needed]], ignore_index=True)
    else:
        combined = paired
    combined = combined[combined['Group'] != 'NAA']
    X_T0 = combined[[f'{c}_T0' for c in protein_cols]].values
    X_T1 = combined[[f'{c}_T1' for c in protein_cols]].values
    X = np.stack([X_T0, X_T1], axis=1)
    label_encoder = LabelEncoder()
    y = label_encoder.fit_transform(combined['Group'].values)
    print(f'Class distribution:')
    for i, cls in enumerate(label_encoder.classes_):
        print(f'  {cls}: {(y == i).sum()}')
    return X, y, protein_cols, label_encoder
 
X_raw, y, protein_names, label_encoder = load_data(SAMPLES_PATH, PROTEINS_PATH)
class_names = list(label_encoder.classes_)
\end{lstlisting}
 
\begin{lstlisting}
# Preprocessing
n_samples, n_timepoints, n_proteins = X_raw.shape
valid_mask = ~np.all(np.isnan(X_raw), axis=(0, 1))
X_raw = X_raw[:, :, valid_mask]
protein_names = [p for p, v in zip(protein_names, valid_mask) if v]
imputer = SimpleImputer(strategy='median')
X_imp = np.zeros_like(X_raw)
for t in range(n_timepoints):
    X_imp[:, t, :] = imputer.fit_transform(X_raw[:, t, :])
X_raw = X_imp
 
# Feature Selection: Top 15 (όπως στο previous)
K = 15
selector = SelectKBest(f_classif, k=K)
selector.fit(X_raw[:, 0, :], y)
selected_idx = np.where(selector.get_support())[0]
X_selected = X_raw[:, :, selected_idx]
selected_proteins = [protein_names[i] for i in selected_idx]
 
# Temporal feature matrix
X = np.hstack([X_selected[:, 0, :], X_selected[:, 1, :]])
print(f'X shape: {X.shape}, N={len(y)}, Classes={class_names}')
\end{lstlisting}
 
\subsubsection*{1. Nested Cross-Validation (Αμερόληπτη Εκτίμηση)}
 
\begin{lstlisting}
# === Nested CV για LR ===
print('='*70)
print('NESTED CROSS-VALIDATION: Logistic Regression')
print('='*70)
 
outer_cv = StratifiedKFold(n_splits=5, shuffle=True, random_state=SEED)
inner_cv = StratifiedKFold(n_splits=3, shuffle=True, random_state=SEED)
 
lr_grid = {'clf__C': [0.001, 0.01, 0.1, 0.5, 1.0, 5.0, 10.0, 50.0]}
 
pipe = Pipeline([
    ('scaler', StandardScaler()),
    ('clf', LogisticRegression(
        max_iter=2000, class_weight='balanced',
        random_state=SEED, solver='lbfgs', multi_class='multinomial'
    ))
])
 
gs = GridSearchCV(pipe, lr_grid, cv=inner_cv, scoring='accuracy', n_jobs=-1)
 
# Nested CV
nested_scores = []
best_Cs = []
 
for fold, (train_idx, val_idx) in enumerate(outer_cv.split(X, y), 1):
    X_train, X_val = X[train_idx], X[val_idx]
    y_train, y_val = y[train_idx], y[val_idx]
 
    gs.fit(X_train, y_train)
    val_acc = gs.score(X_val, y_val)
    best_C = gs.best_params_['clf__C']
 
    nested_scores.append(val_acc)
    best_Cs.append(best_C)
    print(f'  Fold {fold}: Best C={best_C}, Val Acc={val_acc:.3f}')
 
nested_mean = np.mean(nested_scores)
nested_std = np.std(nested_scores)
print(f'\n>>> Nested CV: Acc = {nested_mean:.3f} ± {nested_std:.3f}')
print(f'>>> Naive tuning (prev notebook): 0.662 ± 0.095')
print(f'>>> Bias correction: {nested_mean - 0.662:+.3f}')
print(f'\nSelected C values per fold: {best_Cs}')
\end{lstlisting}
 
\subsubsection*{2. Repeated Cross-Validation}
 
\begin{lstlisting}
# === Repeated CV ===
print('='*70)
print('REPEATED CROSS-VALIDATION (10 × 5-fold)')
print('='*70)
 
# Pipeline (με τη βέλτιστη C από το naive tuning)
best_C = 0.5  # Από το προηγούμενο notebook
 
pipe_best = Pipeline([
    ('scaler', StandardScaler()),
    ('clf', LogisticRegression(
        max_iter=2000, class_weight='balanced',
        random_state=SEED, solver='lbfgs',
        multi_class='multinomial', C=best_C
    ))
])
 
rkf = RepeatedStratifiedKFold(n_splits=5, n_repeats=10, random_state=SEED)
rep_scores = cross_val_score(pipe_best, X, y, cv=rkf, scoring='accuracy', n_jobs=-1)
 
print(f'  Total folds: {len(rep_scores)}')
print(f'  Mean Acc: {rep_scores.mean():.3f} ± {rep_scores.std():.3f}')
print(f'  Min: {rep_scores.min():.3f}  |  Max: {rep_scores.max():.3f}')
print(f'  Median: {np.median(rep_scores):.3f}')
 
# Distribution plot
fig, ax = plt.subplots(figsize=(10, 4))
ax.hist(rep_scores, bins=15, color='#1D9E75', alpha=0.7, edgecolor='white')
ax.axvline(rep_scores.mean(), color='darkred', linestyle='--', linewidth=2,
           label=f'Mean = {rep_scores.mean():.3f}')
ax.axvline(1/3, color='gray', linestyle=':', label='Random (33%)')
ax.set_xlabel('Accuracy', fontsize=12)
ax.set_ylabel('Frequency', fontsize=12)
ax.set_title(f'Distribution of Accuracy across 50 folds (LR, C={best_C})', fontsize=12)
ax.legend()
plt.tight_layout()
plt.savefig('repeated_cv_distribution.png', dpi=150, bbox_inches='tight')
plt.show()
\end{lstlisting}
 
\subsubsection*{3. Permutation Test}
 
\begin{lstlisting}
# === Permutation Test ===
print('='*70)
print('PERMUTATION TEST (1000 permutations)')
print('='*70)
print('This will take a few minutes...')
 
score, perm_scores, pvalue = permutation_test_score(
    pipe_best, X, y,
    scoring='accuracy',
    cv=StratifiedKFold(n_splits=5, shuffle=True, random_state=SEED),
    n_permutations=1000,
    random_state=SEED,
    n_jobs=-1
)
 
print(f'\n  Real accuracy: {score:.3f}')
print(f'  Random permutation mean: {perm_scores.mean():.3f}')
print(f'  Permutation std: {perm_scores.std():.3f}')
print(f'  p-value: {pvalue:.4f}')
print(f'\n>>> {"STATISTICALLY SIGNIFICANT" if pvalue < 0.05 else "NOT significant"}: p {"<" if pvalue < 0.05 else ">="} 0.05')
 
# Plot
fig, ax = plt.subplots(figsize=(10, 4))
ax.hist(perm_scores, bins=30, color='gray', alpha=0.6, edgecolor='white',
        label=f'Permuted labels (n=1000)')
ax.axvline(score, color='darkred', linestyle='-', linewidth=3,
           label=f'Real score = {score:.3f}')
ax.axvline(1/3, color='blue', linestyle=':', alpha=0.5, label='Random (33%)')
ax.set_xlabel('Accuracy', fontsize=12)
ax.set_ylabel('Frequency', fontsize=12)
ax.set_title(f'Permutation Test: Is 66% significantly better than random?\n'
             f'p-value = {pvalue:.4f}', fontsize=12)
ax.legend()
plt.tight_layout()
plt.savefig('permutation_test.png', dpi=150, bbox_inches='tight')
plt.show()
\end{lstlisting}
 
\subsubsection*{4. Bootstrap Confidence Intervals}
 
\begin{lstlisting}
# === Bootstrap CI ===
print('='*70)
print('BOOTSTRAP 95% CONFIDENCE INTERVAL')
print('='*70)
 
# Χρησιμοποιούμε τα 50 scores από Repeated CV
n_bootstrap = 2000
bootstrap_means = []
 
rng = np.random.RandomState(SEED)
for _ in range(n_bootstrap):
    sample = rng.choice(rep_scores, size=len(rep_scores), replace=True)
    bootstrap_means.append(sample.mean())
 
bootstrap_means = np.array(bootstrap_means)
ci_lower = np.percentile(bootstrap_means, 2.5)
ci_upper = np.percentile(bootstrap_means, 97.5)
 
print(f'  Point estimate: {rep_scores.mean():.3f}')
print(f'  Bootstrap mean: {bootstrap_means.mean():.3f}')
print(f'  95% CI: [{ci_lower:.3f}, {ci_upper:.3f}]')
print(f'  Width: {ci_upper - ci_lower:.3f}')
 
# Plot
fig, ax = plt.subplots(figsize=(10, 4))
ax.hist(bootstrap_means, bins=40, color='#1D9E75', alpha=0.7, edgecolor='white')
ax.axvline(rep_scores.mean(), color='darkred', linestyle='-', linewidth=2, label='Point estimate')
ax.axvline(ci_lower, color='orange', linestyle='--', linewidth=2, label=f'2.5%: {ci_lower:.3f}')
ax.axvline(ci_upper, color='orange', linestyle='--', linewidth=2, label=f'97.5%: {ci_upper:.3f}')
ax.set_xlabel('Bootstrap Mean Accuracy', fontsize=12)
ax.set_ylabel('Frequency', fontsize=12)
ax.set_title(f'Bootstrap Distribution (n={n_bootstrap})\n95% CI: [{ci_lower:.3f}, {ci_upper:.3f}]', fontsize=12)
ax.legend()
plt.tight_layout()
plt.savefig('bootstrap_ci.png', dpi=150, bbox_inches='tight')
plt.show()
\end{lstlisting}
 
\subsubsection*{5. Learning Curve}
 
\begin{lstlisting}
# === Learning Curve ===
print('='*70)
print('LEARNING CURVE ANALYSIS')
print('='*70)
 
train_sizes = np.linspace(0.3, 1.0, 8)
train_sizes_abs, train_scores, val_scores = learning_curve(
    pipe_best, X, y,
    train_sizes=train_sizes,
    cv=StratifiedKFold(n_splits=5, shuffle=True, random_state=SEED),
    scoring='accuracy',
    n_jobs=-1,
    random_state=SEED
)
 
train_mean = train_scores.mean(axis=1)
train_std = train_scores.std(axis=1)
val_mean = val_scores.mean(axis=1)
val_std = val_scores.std(axis=1)
 
# Plot
fig, ax = plt.subplots(figsize=(10, 5))
ax.plot(train_sizes_abs, train_mean, 'o-', color='#E24B4A', label='Training')
ax.fill_between(train_sizes_abs, train_mean - train_std, train_mean + train_std,
                color='#E24B4A', alpha=0.2)
ax.plot(train_sizes_abs, val_mean, 'o-', color='#1D9E75', label='Validation')
ax.fill_between(train_sizes_abs, val_mean - val_std, val_mean + val_std,
                color='#1D9E75', alpha=0.2)
ax.axhline(y=1/3, color='gray', linestyle='--', alpha=0.5, label='Random (33%)')
ax.set_xlabel('Training Set Size', fontsize=12)
ax.set_ylabel('Accuracy', fontsize=12)
ax.set_title('Learning Curve (LR, C=0.5)', fontsize=13)
ax.legend(loc='lower right')
ax.grid(alpha=0.3)
plt.tight_layout()
plt.savefig('learning_curve.png', dpi=150, bbox_inches='tight')
plt.show()
 
gap = train_mean[-1] - val_mean[-1]
print(f'\n  Final train acc: {train_mean[-1]:.3f}')
print(f'  Final val acc:   {val_mean[-1]:.3f}')
print(f'  Train-Val gap:   {gap:.3f}')
if gap > 0.15:
    print('  >>> ΥΠΑΡΧΕΙ overfitting - περισσότερα δεδομένα θα βοηθούσαν')
elif gap > 0.08:
    print('  >>> Μικρό gap - μερικά παραπάνω δεδομένα θα βοηθούσαν')
else:
    print('  >>> Μικρό gap - το μοντέλο γενικεύει καλά')
\end{lstlisting}
 
\subsubsection*{6. Calibration Analysis}
 
\begin{lstlisting}
# === Calibration Analysis ===
print('='*70)
print('CALIBRATION ANALYSIS')
print('='*70)
 
from sklearn.model_selection import cross_val_predict
 
# Cross-validated probabilities
proba_cv = cross_val_predict(
    pipe_best, X, y, cv=5, method='predict_proba', n_jobs=-1
)
 
# Για κάθε κλάση, reliability diagram
fig, axes = plt.subplots(1, 3, figsize=(15, 5))
 
for idx, class_name in enumerate(class_names):
    y_binary = (y == idx).astype(int)
    proba_class = proba_cv[:, idx]
 
    # Bin predictions and compute empirical probability
    fraction_pos, mean_pred = calibration_curve(y_binary, proba_class, n_bins=5)
 
    axes[idx].plot([0, 1], [0, 1], 'k:', label='Perfect calibration')
    axes[idx].plot(mean_pred, fraction_pos, 'o-', color='#1D9E75',
                   linewidth=2, markersize=8, label='LR (C=0.5)')
    axes[idx].set_xlabel('Mean predicted probability', fontsize=11)
    axes[idx].set_ylabel('Fraction of positives', fontsize=11)
    axes[idx].set_title(f'Class: {class_name}', fontsize=12)
    axes[idx].legend()
    axes[idx].grid(alpha=0.3)
    axes[idx].set_xlim(0, 1)
    axes[idx].set_ylim(0, 1)
 
    # Brier score
    brier = brier_score_loss(y_binary, proba_class)
    print(f'  {class_name}: Brier score = {brier:.3f}')
 
plt.suptitle('Reliability Diagrams (per class)', fontsize=13, y=1.02)
plt.tight_layout()
plt.savefig('calibration.png', dpi=150, bbox_inches='tight')
plt.show()
 
print('\nNote: Brier score closer to 0 = better calibration (perfect = 0, random ≈ 0.25 for balanced binary)')
\end{lstlisting}
 
\subsubsection*{7. Per-Class Performance Analysis}
 
\begin{lstlisting}
# === Detailed Classification Report ===
print('='*70)
print('PER-CLASS PERFORMANCE (5-Fold CV predictions)')
print('='*70)
 
y_pred_cv = cross_val_predict(pipe_best, X, y, cv=5, n_jobs=-1)
 
print('\nClassification Report:')
print(classification_report(y, y_pred_cv, target_names=class_names, digits=3))
 
# Κλινικά σημαντικές μετρικές
print('\nΚλινικά Σημαντικές Μετρικές:')
for idx, cls in enumerate(class_names):
    y_bin = (y == idx).astype(int)
    y_pred_bin = (y_pred_cv == idx).astype(int)
 
    tp = ((y_bin == 1) & (y_pred_bin == 1)).sum()
    tn = ((y_bin == 0) & (y_pred_bin == 0)).sum()
    fp = ((y_bin == 0) & (y_pred_bin == 1)).sum()
    fn = ((y_bin == 1) & (y_pred_bin == 0)).sum()
 
    sensitivity = tp / (tp + fn) if (tp + fn) > 0 else 0  # Recall
    specificity = tn / (tn + fp) if (tn + fp) > 0 else 0
    ppv = tp / (tp + fp) if (tp + fp) > 0 else 0  # Precision
    npv = tn / (tn + fn) if (tn + fn) > 0 else 0
 
    print(f'\n  {cls} (n={y_bin.sum()}):')
    print(f'    Sensitivity (Recall): {sensitivity:.3f}')
    print(f'    Specificity:          {specificity:.3f}')
    print(f'    PPV (Precision):      {ppv:.3f}')
    print(f'    NPV:                  {npv:.3f}')
\end{lstlisting}
 
\subsubsection*{8. Τελική Σύνοψη}
 
\begin{lstlisting}
# === Summary ===
print('='*70)
print('FINAL SUMMARY: Rigorous Evaluation of LR (C=0.5)')
print('='*70)
 
summary = {
    'Metric': [
        'Naive 5-Fold CV (biased)',
        'Nested CV (unbiased)',
        'Repeated 10×5-Fold CV',
        'Bootstrap 95% CI',
        'Permutation Test p-value',
        'Train-Val Gap',
    ],
    'Value': [
        '0.662 ± 0.095',
        f'{nested_mean:.3f} ± {nested_std:.3f}',
        f'{rep_scores.mean():.3f} ± {rep_scores.std():.3f}',
        f'[{ci_lower:.3f}, {ci_upper:.3f}]',
        f'{pvalue:.4f} {"(significant)" if pvalue < 0.05 else "(NOT significant)"}',
        f'{gap:.3f}',
    ]
}
 
df_summary = pd.DataFrame(summary)
print(df_summary.to_string(index=False))
 
# Save
df_summary.to_csv('final_evaluation.csv', index=False)
 
# Συμπέρασμα
print('\n' + '='*70)
print('ΣΥΜΠΕΡΑΣΜΑ:')
print('='*70)
if pvalue < 0.05:
    print(f'✓ Το μοντέλο είναι στατιστικά σημαντικό (p={pvalue:.4f})')
if ci_lower > 1/3:
    print(f'✓ 95% CI πάνω από τυχαίο ({ci_lower:.3f} > 0.333)')
print(f'✓ Αμερόληπτη εκτίμηση: {nested_mean:.3f} (vs naive {0.662:.3f})')
print(f'✓ Σταθερή απόδοση σε 50 folds: {rep_scores.mean():.3f} ± {rep_scores.std():.3f}')
\end{lstlisting}




\end{document}
